% Les rites de veuvage — Allemand Tle A — Cours signé OnBuch+
% Programme officiel MINESEC. Compiler avec : tectonic AL03-les-rites-de-veuvage.tex
\newcommand{\DOCMATIERE}{Allemand}
\newcommand{\DOCNIVEAU}{Tle A}
\newcommand{\DOCMODULE}{Chapitre 1 : Vie familiale et sociale}
\newcommand{\DOCLECON}{AL03}
\newcommand{\DOCTITRE}{Les rites de veuvage}
% OnBuch+ — préambule « cours signés » ENRICHI (compilé avec tectonic / XeLaTeX)
% Système visuel « Pop » d'OnBuch+ : fond crème (jamais blanc), bordures encre
% épaisses, ombres « dures » décalées, titres Archivo Black en capitales.
% Version MATHÉMATIQUES : graphes (pgfplots/tikz), boîtes pédagogiques
% (définition, propriété, méthode, attention, le savais-tu, exercice résolu,
% exercice, TP, bilan), page de garde et signature OnBuch+ ancrée en bas.
%
% Macros attendues dans main.tex AVANT \input{preamble} :
%   \DOCMATIERE  \DOCNIVEAU  \DOCMODULE  \DOCLECON (numéro)  \DOCTITRE
\documentclass[11pt]{article}

\usepackage{fontspec}
\usepackage[ngerman,french]{babel} % français = langue principale ; l'allemand via \all{..} / textebox
\usepackage[a4paper,margin=2cm,top=3.4cm,bottom=2.8cm,headheight=1.4cm,footskip=1.3cm]{geometry}
\usepackage{amsmath,amssymb,amsfonts}
\usepackage{mathtools} % \xmapsto, \coloneqq... souvent utilisés par les modèles
\usepackage{cancel} % simplifications barrées (bilans de forces)
% \abs{z} (module, valeur absolue) et \norm{u} (norme), utilisés par les modèles
\providecommand{\abs}{}\let\abs\relax\DeclarePairedDelimiter{\abs}{\lvert}{\rvert}
\providecommand{\norm}{}\let\norm\relax\DeclarePairedDelimiter{\norm}{\lVert}{\rVert}
\usepackage[table]{xcolor} % [table] : fournit aussi \rowcolors (alternance de lignes)
\usepackage{pagecolor}
\usepackage{tikz}
\usetikzlibrary{arrows.meta,positioning,calc,shapes.geometric,shapes.misc,shapes.arrows,decorations.pathmorphing,decorations.pathreplacing,decorations.markings,patterns,shadows,mindmap,trees,fit,backgrounds,matrix,intersections,angles,quotes}
\usepackage{pgfplots}
\usepackage[version=4]{mhchem} % équations nucléaires \ce{^{235}_{92}U}
\usepackage[european,siunitx]{circuitikz} % schémas électriques
\pgfplotsset{compat=1.18}
\usepgfplotslibrary{fillbetween,groupplots} % aires entre courbes (calcul intégral)
\usepackage{enumitem}
\usepackage{fancyhdr}
% [most] charge toutes les bibliothèques tcolorbox (listings, minted, raster,
% documentation...) et gonfle inutilement la mémoire TeX sur un document long
% (~300 boîtes) : on ne charge que ce qui est réellement utilisé.
\usepackage[skins,breakable]{tcolorbox}
\usepackage{graphicx}
\usepackage{colortbl}
\usepackage{booktabs}
\usepackage{tabularx}
\usepackage{diagbox} % case d'en-tête diagonale des tableaux à double entrée
% Colonnes extensibles centrée / gauche / droite (C, L, R), souvent utilisées par les modèles
\newcolumntype{C}{>{\centering\arraybackslash}X}
\newcolumntype{L}{>{\raggedright\arraybackslash}X}
\newcolumntype{R}{>{\raggedleft\arraybackslash}X}
\usepackage{array}
\usepackage{multicol}
\usepackage{multirow}
\usepackage{siunitx}
% parse-numbers=false : accepte aussi \SI{2,5 \times 10^{-3}}{...}
\sisetup{locale=FR,output-decimal-marker={,},group-minimum-digits=4,parse-numbers=false}
\DeclareSIUnit\ohm{\ensuremath{\symup{\Omega}}} % U+2126 absent de la police
\DeclareSIUnit\division{div} % réglages d'oscilloscope (ms/div, V/div)
\DeclareSIUnit\tr{tr} % tours (vitesse de rotation, ex. \SI{1500}{\tr\per\minute})
\DeclareSIUnit\molar{M} % concentration molaire (ex. \SI{3}{\molar}, \SI{1}{\micro\molar})
\DeclareSIUnit\an{an} % année (le modèle confond parfois avec \year, un compteur TeX) — ex. \SI{5}{\kilo\meter\per\an}
\DeclareSIUnit\FCFA{FCFA} % franc CFA (ex. \SI{500}{\FCFA\per\kilo\gram})
\DeclareSIUnit\degreeBrix{\textdegree Brix} % teneur en sucre d'un sirop/jus (réfractométrie)
\DeclareSIUnit\month{mois} % le modèle utilise parfois \month, un compteur TeX, comme unité de durée
\usepackage{qrcode}
% \degree : commande fréquemment utilisée par les modèles pour un angle ou une
% température en dehors de \SI{}{} (non fournie par les paquets chargés).
\providecommand{\degree}{^{\circ}}
% \qed (fin de démonstration), utilisé par les modèles sans amsthm
\providecommand{\qed}{\hfill$\square$}
% \barre : double barre des valeurs interdites dans un tableau de variations (tkz-tab)
\providecommand{\barre}{\ensuremath{\|}}
% environnement proof (amsthm non chargé) utilisé par les modèles
\ifdefined\proof\else\newenvironment{proof}[1][Démonstration]{\par\noindent\textit{#1.}\ }{\qed\par}\fi
\usepackage{lastpage}
\usepackage{hyperref}
\hypersetup{hidelinks}

% ============================================================
% Palette « Pop » OnBuch+ — mêmes hex que la classe P (plus_ui.dart)
% ============================================================
\definecolor{poporange}{HTML}{F59321}
\definecolor{popdark}{HTML}{E07A0C}
\definecolor{popcream}{HTML}{FFF6E8}
\definecolor{popink}{HTML}{1C1714}
\definecolor{poppurple}{HTML}{7A5AE0}
\definecolor{popgreen}{HTML}{2E9E6A}
\definecolor{poppink}{HTML}{E0457B}
\definecolor{popblue}{HTML}{2F7DE1}
\definecolor{popgold}{HTML}{C98A12}
\definecolor{popmuted}{HTML}{8A7F76}
\definecolor{popline}{HTML}{EADFD2}
\definecolor{popgray}{HTML}{8A7F76}
\colorlet{popgrey}{popgray}
% Alias tolérants (le modèle nomme parfois les couleurs différemment)
\colorlet{popwhite}{white}
\colorlet{popblack}{popink}
\colorlet{popred}{poppink}
\colorlet{popyellow}{popgold}
% Teintes claires (fonds de tableaux/encadrés) — le modèle nomme parfois
% les couleurs avec un suffixe L (ex. popblueL) sans les avoir définies.
\colorlet{poporangeL}{poporange!20}
\colorlet{popdarkL}{popdark!20}
\colorlet{poppurpleL}{poppurple!20}
\colorlet{popgreenL}{popgreen!20}
\colorlet{poppinkL}{poppink!20}
\colorlet{popblueL}{popblue!20}
\colorlet{popgoldL}{popgold!20}
\colorlet{popredL}{popred!20}
\colorlet{popgrayL}{popgray!20}
% Teintes claires pour fonds de tableaux / zones
\colorlet{popblueL}{popblue!10!white}
\colorlet{poporangeL}{poporange!14!white}
\colorlet{popgreenL}{popgreen!12!white}
\colorlet{poppurpleL}{poppurple!12!white}
\colorlet{poppinkL}{poppink!12!white}
\colorlet{popgoldL}{popgold!14!white}

\pagecolor{popcream}

% ============================================================
% Typographie
% ============================================================
\defaultfontfeatures{Ligatures=TeX}
\setmainfont{PlusJakartaSans-Medium.ttf}[Path=../../../fonts/,BoldFont=PlusJakartaSans-Bold.ttf]
% Maths en sans-serif (Fira Math) avec chiffres et lettres droites de Plus
% Jakarta, pour que les formules se fondent dans le texte.
\usepackage[math-style=ISO,bold-style=ISO]{unicode-math}
\setmathfont{FiraMath-Regular.otf}[Path=../../../fonts/]
\setmathfont{PlusJakartaSans-Medium.ttf}[Path=../../../fonts/,range={up/num,up/{Latin,latin},\mathup}]
\usepackage{icomma} % virgule décimale sans espace en mode math
% Caractères mathématiques Unicode absents des polices de texte (Plus Jakarta,
% Archivo Black) : rendus par leur équivalent mathématique, en texte comme en
% formule — sinon ils s'affichent en carré vide (ex. « Arithmétique dans ℤ »).
\usepackage{newunicodechar}
\newunicodechar{ℝ}{\ensuremath{\mathbb{R}}}
\newunicodechar{ℤ}{\ensuremath{\mathbb{Z}}}
\newunicodechar{ℂ}{\ensuremath{\mathbb{C}}}
\newunicodechar{ℕ}{\ensuremath{\mathbb{N}}}
\newunicodechar{ℚ}{\ensuremath{\mathbb{Q}}}
\newunicodechar{𝔻}{\ensuremath{\mathbb{D}}}
\newunicodechar{α}{\ensuremath{\alpha}}
\newunicodechar{β}{\ensuremath{\beta}}
\newunicodechar{γ}{\ensuremath{\gamma}}
\newunicodechar{δ}{\ensuremath{\delta}}
\newunicodechar{ε}{\ensuremath{\varepsilon}}
\newunicodechar{θ}{\ensuremath{\theta}}
\newunicodechar{λ}{\ensuremath{\lambda}}
\newunicodechar{μ}{\ensuremath{\mu}}
\newunicodechar{σ}{\ensuremath{\sigma}}
\newunicodechar{τ}{\ensuremath{\tau}}
\newunicodechar{φ}{\ensuremath{\varphi}}
\newunicodechar{ψ}{\ensuremath{\psi}}
\newunicodechar{ω}{\ensuremath{\omega}}
\newunicodechar{Σ}{\ensuremath{\Sigma}}
% majuscules produites par \MakeUppercase dans les titres (α-aminés -> Α)
\newunicodechar{Α}{\ensuremath{\alpha}}
\newunicodechar{Β}{\ensuremath{\beta}}
\newunicodechar{Γ}{\ensuremath{\Gamma}}
\newunicodechar{Δ}{\ensuremath{\Delta}}
\newunicodechar{Ε}{\ensuremath{\varepsilon}}
\newunicodechar{Θ}{\ensuremath{\Theta}}
\newunicodechar{Λ}{\ensuremath{\Lambda}}
\newunicodechar{Μ}{\ensuremath{\mu}}
\newunicodechar{Τ}{\ensuremath{\tau}}
\newunicodechar{Φ}{\ensuremath{\Phi}}
\newunicodechar{Ψ}{\ensuremath{\Psi}}
\newunicodechar{Ω}{\ensuremath{\Omega}}
\newunicodechar{↦}{\ensuremath{\mapsto}}
\newunicodechar{∈}{\ensuremath{\in}}
\newunicodechar{∉}{\ensuremath{\notin}}
\newunicodechar{∧}{\ensuremath{\wedge}}
\newunicodechar{⇔}{\ensuremath{\Leftrightarrow}}
\newunicodechar{⟺}{\ensuremath{\Longleftrightarrow}}
\newunicodechar{⇒}{\ensuremath{\Rightarrow}}
\newunicodechar{⟹}{\ensuremath{\Longrightarrow}}
\newunicodechar{∘}{\ensuremath{\circ}}
\newunicodechar{∪}{\ensuremath{\cup}}
\newunicodechar{∩}{\ensuremath{\cap}}
\newunicodechar{≡}{\ensuremath{\equiv}}
\newunicodechar{⊂}{\ensuremath{\subset}}
\newunicodechar{⊥}{\ensuremath{\perp}}
\newunicodechar{∃}{\ensuremath{\exists}}
\newunicodechar{∀}{\ensuremath{\forall}}
\newunicodechar{∛}{\ensuremath{\sqrt[3]{\,}}}
\newunicodechar{∜}{\ensuremath{\sqrt[4]{\,}}}
\newunicodechar{⃗}{\ensuremath{{}^{\to}}}
\newunicodechar{ⁿ}{\ifmmode^{n}\else\textsuperscript{\ensuremath{n}}\fi}
\newunicodechar{ˣ}{\ifmmode^{x}\else\textsuperscript{\ensuremath{x}}\fi}
\newunicodechar{ᵇ}{\ifmmode^{b}\else\textsuperscript{\ensuremath{b}}\fi}
\newunicodechar{ʳ}{\ifmmode^{r}\else\textsuperscript{\ensuremath{r}}\fi}
\newunicodechar{ᵘ}{\ifmmode^{u}\else\textsuperscript{\ensuremath{u}}\fi}
\newunicodechar{ʸ}{\ifmmode^{y}\else\textsuperscript{\ensuremath{y}}\fi}
\newunicodechar{ᵃ}{\ifmmode^{a}\else\textsuperscript{\ensuremath{a}}\fi}
\newunicodechar{ᵉ}{\ifmmode^{e}\else\textsuperscript{\ensuremath{e}}\fi}
\newunicodechar{ᵏ}{\ifmmode^{k}\else\textsuperscript{\ensuremath{k}}\fi}
\newunicodechar{ᵖ}{\ifmmode^{p}\else\textsuperscript{\ensuremath{p}}\fi}
\newunicodechar{ᴺ}{\ifmmode^{N}\else\textsuperscript{\ensuremath{N}}\fi}
\newunicodechar{ᶜ}{\ifmmode^{c}\else\textsuperscript{\ensuremath{c}}\fi}
\newunicodechar{ʰ}{\ifmmode^{h}\else\textsuperscript{\ensuremath{h}}\fi}
\newunicodechar{ᵠ}{\ifmmode^{q}\else\textsuperscript{\ensuremath{q}}\fi}
\newunicodechar{ₙ}{\ifmmode_{n}\else\textsubscript{\ensuremath{n}}\fi}
\newunicodechar{ₐ}{\ifmmode_{a}\else\textsubscript{\ensuremath{a}}\fi}
\newunicodechar{ᵢ}{\ifmmode_{i}\else\textsubscript{\ensuremath{i}}\fi}
\newunicodechar{ᵣ}{\ifmmode_{r}\else\textsubscript{\ensuremath{r}}\fi}
\newunicodechar{ₖ}{\ifmmode_{k}\else\textsubscript{\ensuremath{k}}\fi}
\newunicodechar{ᵦ}{\ifmmode_{\beta}\else\textsubscript{\ensuremath{\beta}}\fi}
\newunicodechar{ₓ}{\ifmmode_{x}\else\textsubscript{\ensuremath{x}}\fi}
\newfontfamily\popdisplay{ArchivoBlack-Regular.ttf}[Path=../../../fonts/]
\newfontfamily\poplogo{SpaceGrotesk-Bold.ttf}[Path=../../../fonts/]
\newfontfamily\popsemi{PlusJakartaSans-SemiBold.ttf}[Path=../../../fonts/]
\newfontfamily\popxbold{PlusJakartaSans-ExtraBold.ttf}[Path=../../../fonts/]
\newfontfamily\popxboldls{PlusJakartaSans-ExtraBold.ttf}[Path=../../../fonts/,LetterSpace=3.0]

\setlist[enumerate]{leftmargin=1.7em,itemsep=5pt,topsep=4pt}
\setlist[itemize]{leftmargin=1.7em,itemsep=4pt,topsep=4pt,label={\color{poporange}\textbullet}}
\setlength{\parindent}{0pt}
\setlength{\parskip}{5pt}
\renewcommand{\arraystretch}{1.25}

% pgfplots : style Pop par défaut
\pgfplotsset{
  every axis/.append style={axis line style={popink,line width=1pt},tick style={popink},
    grid style={popline,line width=0.5pt},label style={font=\small},tick label style={font=\footnotesize}},
  popcurve/.style={poporange,line width=1.8pt,no markers,smooth},
}
% Même style disponible dans un \draw TikZ simple (hors pgfplots)
\tikzset{popcurve/.style={poporange,line width=1.8pt}}
\tikzset{popgrid/.style={popline,very thin}}
\colorlet{popcurve}{poporange} % le nom du style est parfois utilisé comme couleur

% ============================================================
% Pill / squiggle
% ============================================================
\newcommand{\poppill}[3]{%
  \tikz[baseline=-0.55ex]\node[draw=popink,line width=1.2pt,rounded corners=2.6pt,%
    fill=#2,inner xsep=6.5pt,inner ysep=3pt]{%
    \popxboldls\fontsize{7}{7}\selectfont\color{#3}\MakeUppercase{#1}};%
}
\newcommand{\popsquiggle}[1]{%
  \tikz[baseline=-0.4ex]{
    \draw[#1,line width=2.6pt,line cap=round]
      plot [smooth] coordinates {(0,0.09) (0.34,0.01) (0.68,0.21) (1.02,0.01) (1.36,0.10)};
  }%
}
% Mot-clé surligné (marqueur orange sous le texte)
\newcommand{\cle}[1]{{\popxbold\color{popdark}#1}}

% ============================================================
% En-tête / pied
% ============================================================
\pagestyle{fancy}
\fancyhf{}
\renewcommand{\headrulewidth}{0pt}
\renewcommand{\footrulewidth}{0pt}
\lhead{\raisebox{-0.15ex}{\poplogo\fontsize{13}{13}\selectfont\color{popink}OnBuch\color{poporange}+}}
\rhead{\poppill{\DOCMATIERE\ \textbullet\ \DOCNIVEAU\ \textbullet\ Leçon \DOCLECON}{white}{popink}}
\lfoot{\popxbold\fontsize{7.6}{7.6}\selectfont\color{popmuted}\MakeUppercase{Cours signé OnBuch+}\ \textbullet\ {\popsemi\DOCTITRE}}
\rfoot{\tikz[baseline=-0.6ex]\node[rounded corners=8.5pt,fill=popink,minimum height=17pt,minimum width=17pt,inner xsep=5pt,inner sep=0]{\popxbold\fontsize{8}{8}\selectfont\color{popcream}\thepage\,/\,\pageref*{LastPage}};}

% ============================================================
% Titres
% ============================================================
\newcommand{\chaptitre}[2]{%
  \begin{tcolorbox}[enhanced,colback=poporange,colframe=popink,boxrule=2.6pt,arc=11pt,
      drop shadow=popink,left=15pt,right=15pt,top=12pt,bottom=11pt,before skip=3pt,after skip=18pt]
    \poppill{#2}{popink}{popcream}\par\vspace{9pt}
    {\raggedright\hyphenpenalty=10000\exhyphenpenalty=10000\popdisplay\fontsize{22}{24}\selectfont\color{popcream}\MakeUppercase{#1}\par}
  \end{tcolorbox}
}
% Section numérotée : pastille orange avec le numéro + titre Archivo
\newcounter{popsec}
\newcommand{\coursec}[1]{%
  \stepcounter{popsec}\setcounter{popsub}{0}%
  \par\vspace{16pt}\noindent
  \begin{minipage}{\linewidth}
  \tikz[baseline=-0.7ex]\node[draw=popink,line width=1.6pt,fill=poporange,rounded corners=4pt,minimum size=22pt,inner sep=0]{\popdisplay\fontsize{12}{12}\selectfont\color{popcream}\thepopsec};%
  \hspace{8pt}{\popdisplay\fontsize{15}{17}\selectfont\color{popink}\MakeUppercase{#1}}\par
  \vspace{4pt}\noindent{\color{poporange}\rule{36pt}{3.6pt}}
  \end{minipage}\par\nopagebreak\vspace{8pt}
}
\newcounter{popsub}
\newcommand{\courssub}[1]{%
  \stepcounter{popsub}%
  \par\vspace{9pt}\noindent{\popxbold\fontsize{11.5}{13}\selectfont\color{popink}{\color{poporange}\thepopsec.\thepopsub}\hspace{6pt}#1}\par\nopagebreak\vspace{4pt}
}

% ============================================================
% Boîtes pédagogiques — même recette : fond blanc, bordure épaisse colorée,
% ombre dure encre, badge plein. Titre optionnel [#1].
% ============================================================
\tcbset{popbase/.style={breakable,enhanced,colback=white,boxrule=2.3pt,arc=9pt,
  shadow={3pt}{-3pt}{0pt}{popink},left=14pt,right=14pt,top=11pt,bottom=11pt,before skip=11pt,after skip=15pt,
  fontupper=\popsemi\fontsize{10.6}{14.5}\selectfont\color{popink},
  fontlower=\popsemi\fontsize{10.6}{14.5}\selectfont\color{popink}}}
% Coche dessinée (le glyphe ✓ n'existe pas dans les polices du document)
\newcommand{\popcheck}[1]{\tikz[baseline=-0.5ex]\draw[#1,line width=1.6pt,line cap=round,line join=round](-0.13,0.0)--(-0.04,-0.09)--(0.13,0.1);}
\providecommand{\coche}{\popcheck{popgreen}}
\providecommand{\resultat}[1]{\par\smallskip\noindent\fcolorbox{popgreen}{popgreenL}{\parbox{\dimexpr\linewidth-2\fboxsep-2\fboxrule\relax}{\textbf{Résultat.} #1}}\par\smallskip}
\newcommand{\popboxhead}[3]{\poppill{#1}{#2}{white}\ifx\relax#3\relax\else\hspace{6pt}{\popxbold\fontsize{10.6}{12}\selectfont\color{popink}#3}\fi\par\vspace{7pt}}

\newtcolorbox{aretenir}[1][]{popbase,colframe=popgreen,before upper={\popboxhead{À retenir}{popgreen}{#1}}}
% Texte en allemand (document de lecture, dialogue, extrait) : cadre
% d'étude. Un titre optionnel (source, auteur) s'affiche dans l'en-tête.
\newtcolorbox{textebox}[1][]{popbase,colframe=popblue,before upper={\popboxhead{Text}{popblue}{#1}\begin{otherlanguage}{ngerman}},after upper={\end{otherlanguage}}}
% Marques ✓ / ✗ (forme correcte / fautive) souvent utilisées par les modèles
\providecommand{\cross}{\textcolor{poppink}{\ensuremath{\times}}}
\providecommand{\xmark}{\cross}\providecommand{\crossmark}{\cross}\providecommand{\cmark}{\ensuremath{\checkmark}}
% Ligne de réponse / justification à compléter (inventée par les modèles)
\providecommand{\justification}[1]{\par\noindent\textit{Justification~:} \dotfill\par}
% \textipa (package tipa non chargé) : la police ne couvre pas l'API -> simple passage
\providecommand{\textipa}[1]{#1}
% Soulignement (ulem non chargé) : \uline, \ul
\providecommand{\uline}[1]{\underline{#1}}\providecommand{\ul}[1]{\underline{#1}}
% variantes ASCII d'ordinaux écrites en macro
\providecommand{\iere}{\textsuperscript{re}}\providecommand{\ieme}{\textsuperscript{e}}\providecommand{\ere}{\textsuperscript{re}}\providecommand{\eme}{\textsuperscript{e}}\providecommand{\er}{\textsuperscript{er}}\providecommand{\ie}{\textsuperscript{e}}
% Ordinaux français écrits en macro par les modèles : 1\ière, 3\ième
\newcommand{\ière}{\textsuperscript{re}}\newcommand{\ième}{\textsuperscript{e}}
% \exemple (inventée par les modèles) : étiquette « Exemple : »
\providecommand{\exemple}{\textbf{Exemple~:}\ }
% \square utilisable aussi hors mode math (cases à cocher)
\AtBeginDocument{\let\squareMath\square\renewcommand{\square}{\ensuremath{\squareMath}}}
% Mot ou phrase en allemand à l'intérieur d'un paragraphe français
\newcommand{\all}[1]{\ifmmode\text{\textit{#1}}\else\begingroup\selectlanguage{ngerman}\itshape #1\endgroup\fi}
\let\esp\all % alias toléré
\newtcolorbox{exemplebox}[1][]{popbase,colframe=poporange,before upper={\popboxhead{Exemple}{poporange}{#1}}}
\newtcolorbox{definition}[1][]{popbase,colframe=popblue,before upper={\popboxhead{Définition}{popblue}{#1}}}
\newtcolorbox{propriete}[1][]{popbase,colframe=poppurple,before upper={\popboxhead{Propriété}{poppurple}{#1}}}
\newtcolorbox{methode}[1][]{popbase,colframe=popgold,before upper={\popboxhead{Méthode}{popgold}{#1}}}
\newtcolorbox{attention}[1][]{popbase,colframe=poppink,before upper={\popboxhead{Attention}{poppink}{#1}}}
\newtcolorbox{experience}[1][]{popbase,colframe=popblue,colback=popblueL,before upper={\popboxhead{Expérience}{popblue}{#1}}}
% « Le savais-tu ? » : carte violette pleine (contexte, culture, Cameroun)
\newtcolorbox{savaistu}[1][]{popbase,colback=poppurple,colframe=popink,
  fontupper=\popsemi\fontsize{10.4}{14.2}\selectfont\color{white},
  before upper={\poppill{Le savais-tu\,?}{white}{poppurple}\ifx\relax#1\relax\else\hspace{6pt}{\popxbold\color{white}#1}\fi\par\vspace{7pt}}}
% Exercice résolu : énoncé en haut, \tcblower puis la solution sur fond crème
\newcounter{popexo}\newcounter{popexr}
\newtcolorbox{exoresolu}[1][]{popbase,colframe=poporange,colbacklower=popcream,
  segmentation style={popink,line width=1.2pt,dashed},
  before upper={\stepcounter{popexr}\popboxhead{Exercice résolu \thepopexr}{poporange}{#1}},
  before lower={\poppill{Solution}{popink}{popcream}\par\vspace{6pt}}}
% Exercice (non résolu en place) — la correction est dans la partie Corrigés
\newtcolorbox{exercice}[1][]{popbase,colframe=popink,
  before upper={\stepcounter{popexo}\popboxhead{Exercice \thepopexo}{popink}{#1}}}
% Corrigé
\newtcolorbox{corrige}[1][]{popbase,colframe=popgreen,colback=popgreenL,
  before upper={\popboxhead{Corrigé}{popgreen}{#1}}}
% Objectifs / prérequis / bilan
\newtcolorbox{objectifs}{popbase,colframe=popink,colback=poporangeL,
  before upper={\popboxhead{Objectifs de la leçon}{popink}{}}}
\newtcolorbox{prerequis}{popbase,colframe=popink,colback=popblueL,
  before upper={\popboxhead{Prérequis}{popblue}{}}}
\newtcolorbox{bilan}{popbase,colframe=popink,colback=white,boxrule=2.8pt,
  before upper={\popboxhead{Fiche bilan}{poporange}{L'essentiel à connaître pour l'examen}}}
% Figure encadrée avec légende
\newtcolorbox{popfigure}[1][]{enhanced,breakable=false,colback=white,colframe=popline,boxrule=1.4pt,arc=8pt,
  left=10pt,right=10pt,top=10pt,bottom=8pt,before skip=10pt,after skip=12pt,halign=center,
  lower separated=false,halign lower=center,
  fontlower=\popsemi\fontsize{9}{11.5}\selectfont\color{popmuted},
  #1}
\newcommand{\legende}[1]{\par\vspace{6pt}{\popsemi\fontsize{9}{11.5}\selectfont\color{popmuted}#1}}

% ============================================================
% Signature OnBuch+
% ============================================================
\newcommand{\poplogoinline}[1]{{\poplogo\fontsize{#1}{#1}\selectfont\color{popink}OnBuch\color{poporange}+}}
\newcommand{\signatureonbuch}{%
  \begin{tcolorbox}[enhanced,colback=popink,colframe=popink,boxrule=2.2pt,arc=10pt,
      drop shadow=poporange,left=14pt,right=14pt,top=10pt,bottom=10pt,before skip=0pt,after skip=0pt]
    \tikz[baseline=-0.7ex]\node[circle,fill=popgreen,draw=popcream,line width=1.3pt,minimum size=22pt,inner sep=0]{\popcheck{white}};%
    \hspace{9pt}\begin{minipage}[c]{0.62\linewidth}
      {\popxbold\fontsize{12}{13}\selectfont\color{popcream}Signé\ \textbullet\ {\poplogo OnBuch\color{poporange}+}}\par\vspace{2pt}
      {\raggedright\popsemi\fontsize{8}{10.5}\selectfont\color{popline}\MakeUppercase{Équipe pédagogique OnBuch+}\\\MakeUppercase{Conforme au programme officiel MINESEC}\par}
    \end{minipage}\hfill
    {\poplogo\fontsize{22}{22}\selectfont\color{popcream}OnBuch\color{poporange}+}
  \end{tcolorbox}%
}

% ============================================================
% Page de garde
% #1 titre · #2 matière · #3 niveau · #4 module · #5 n° leçon · #6 durée ·
% #7 description / objectifs (texte ou itemize)
% ============================================================
\newcommand{\pagegarde}[7]{%
  \thispagestyle{empty}
  \newgeometry{a4paper,margin=1.7cm,top=1.6cm,bottom=1.4cm}%
  \begin{tikzpicture}[remember picture,overlay]
    \fill[poporange!18] ([xshift=-3.2cm,yshift=-3.6cm]current page.north east) circle (4.6cm);
    \fill[poppurple!14] ([xshift=2.4cm,yshift=7.6cm]current page.south west) circle (3.8cm);
  \end{tikzpicture}%
  {\poplogo\fontsize{30}{30}\selectfont\color{popink}OnBuch\color{poporange}+}\hfill
  \raisebox{0.6ex}{\poppill{Cours signé \textbullet\ Premium}{popink}{popcream}}\par
  \vspace{1pt}\popsquiggle{poppurple}\par
  \vspace{8pt}
  \begin{tcolorbox}[enhanced,colback=poporange,colframe=popink,boxrule=2.8pt,arc=13pt,
      drop shadow=popink,left=18pt,right=18pt,top=12pt,bottom=13pt,after skip=10pt]
    \poppill{#2 \textbullet\ #3}{popink}{popcream}\hspace{5pt}\poppill{#4}{white}{popink}\par\vspace{8pt}
    {\popdisplay\fontsize{14}{14}\selectfont\color{popink}LEÇON #5}\par\vspace{4pt}
    {\raggedright\hyphenpenalty=10000\exhyphenpenalty=10000\popdisplay\fontsize{25}{27}\selectfont\color{popcream}\MakeUppercase{#1}\par}
    \vspace{8pt}
    {\popxbold\fontsize{9.5}{9.5}\selectfont\color{popink}\MakeUppercase{Durée conseillée : #6}}
  \end{tcolorbox}
  \begin{tcolorbox}[enhanced,colback=white,colframe=popink,boxrule=2.2pt,arc=10pt,
      drop shadow=popink,left=16pt,right=16pt,top=10pt,bottom=10pt,after skip=8pt]
    \poppill{Dans cette leçon}{poppurple}{white}\par\vspace{5pt}
    \popsemi\fontsize{9.3}{12.6}\selectfont\color{popink}#7
  \end{tcolorbox}
  \noindent\begin{minipage}[t]{0.487\linewidth}
    \begin{tcolorbox}[enhanced,colback=popgreen,colframe=popink,boxrule=2.2pt,arc=10pt,
        drop shadow=popink,left=11pt,right=11pt,top=10pt,bottom=10pt,height=3.1cm,valign=center]
      \tcbox[colback=white,colframe=popink,boxrule=1.3pt,arc=5pt,boxsep=0pt,nobeforeafter,
        left=3pt,right=3pt,top=3pt,bottom=3pt,valign=center]{\qrcode[height=1.9cm]{https://play.google.com/store/apps/details?id=com.luvvix.onbuch&hl=fr}}%
      \hspace{8pt}\begin{minipage}[c]{0.5\linewidth}
        {\popdisplay\fontsize{11}{12.5}\selectfont\color{white}\MakeUppercase{Télécharge OnBuch}}\par\vspace{4pt}
        {\raggedright\popsemi\fontsize{8.4}{11}\selectfont\color{white}Gratuit \textbullet\ Google Play\\Tuteur IA, annales, concours, résultats\par}
      \end{minipage}
    \end{tcolorbox}
  \end{minipage}\hfill
  \begin{minipage}[t]{0.487\linewidth}
    \begin{tcolorbox}[enhanced,colback=poppurple,colframe=popink,boxrule=2.2pt,arc=10pt,
        drop shadow=popink,left=11pt,right=11pt,top=10pt,bottom=10pt,height=3.1cm,valign=center]
      \tikz[baseline=-0.7ex]\node[circle,fill=white,draw=popink,line width=1.3pt,minimum size=1.6cm,inner sep=0]{\popdisplay\fontsize{24}{24}\selectfont\color{poppurple}+};%
      \hspace{8pt}\begin{minipage}[c]{0.6\linewidth}
        {\popdisplay\fontsize{11}{12.5}\selectfont\color{white}\MakeUppercase{Passe à OnBuch+}}\par\vspace{4pt}
        {\raggedright\popsemi\fontsize{8.4}{11}\selectfont\color{white}Tous les cours signés, Tuteur IA illimité, fiches PDF — onglet OnBuch+ de l'app\par}
      \end{minipage}
    \end{tcolorbox}
  \end{minipage}
  \vspace{8pt}
  \popcontenu
  \vfill
  \signatureonbuch
  \restoregeometry
  \newpage
}

% Rangée « Ce cours contient » (page de garde)
\newcommand{\poptile}[3]{% #1 couleur #2 gros texte #3 légende
  \begin{tcolorbox}[enhanced,colback=#1,colframe=popink,boxrule=2pt,arc=9pt,shadow={2.5pt}{-2.5pt}{0pt}{popink},
      left=6pt,right=6pt,top=9pt,bottom=9pt,halign=center,valign=center,height=1.75cm,width=\linewidth,nobeforeafter]
    {\popdisplay\fontsize{12.5}{13}\selectfont\color{white}\MakeUppercase{#2}}\par\vspace{3pt}
    {\centering\popsemi\fontsize{7.8}{9.5}\selectfont\color{white}#3\par}
  \end{tcolorbox}}
\newcommand{\popcontenu}{%
  \par\noindent{\popxboldls\fontsize{7.5}{7.5}\selectfont\color{popmuted}CE COURS SIGNÉ CONTIENT}\par\vspace{7pt}
  \noindent\begin{minipage}[t]{0.235\linewidth}\poptile{popblue}{Cours}{complet, illustré, pas à pas}\end{minipage}\hfill
  \begin{minipage}[t]{0.235\linewidth}\poptile{poporange}{Méthodes}{et exercices résolus}\end{minipage}\hfill
  \begin{minipage}[t]{0.235\linewidth}\poptile{poppink}{Exercices}{type examen + corrigés}\end{minipage}\hfill
  \begin{minipage}[t]{0.235\linewidth}\poptile{popgreen}{Fiche bilan}{l'essentiel à retenir}\end{minipage}\par}

% Sommaire
\newtcolorbox{sommaire}{enhanced,colback=white,colframe=popink,boxrule=2.2pt,arc=10pt,
  drop shadow=popink,left=18pt,right=18pt,top=13pt,bottom=13pt,before skip=6pt,after skip=20pt,
  before upper={\poppill{Sommaire}{poppurple}{white}\par\vspace{9pt}\popsemi\fontsize{10.6}{14.5}\selectfont\color{popink}}}

% Signature de fin de cours — poussée tout en bas de la dernière page
\newcommand{\findecours}{%
  \par\vfill\nopagebreak
  \begin{center}\popsquiggle{poporange}\end{center}\vspace{4pt}
  \signatureonbuch
}
\AtBeginDocument{\def\llbracket{\mathopen{[\mkern-3.5mu[}}\def\rrbracket{\mathclose{]\mkern-3.5mu]}}} % intervalles d'entiers (glyphe ⟦ absent de la police)
% Glyphes absents de Fira Math : redéfinis à partir de glyphes disponibles
\AtBeginDocument{%
  \def\setminus{\mathbin{\backslash}}\let\smallsetminus\setminus
  \def\vdots{\mathord{\vbox{\baselineskip3.2pt\lineskiplimit0pt\kern4pt\hbox{.}\hbox{.}\hbox{.}}}}%
  \def\ddots{\mathinner{\mkern1mu\raise6.5pt\hbox{.}\mkern2mu\raise3.6pt\hbox{.}\mkern2mu\raise0.7pt\hbox{.}\mkern1mu}}%
  \def\triangle{\mathord{\scalebox{0.9}{$\symup{\Delta}$}}}%
  \def\nabla{\mathord{\raisebox{\depth}{\scalebox{1}[-1]{$\symup{\Delta}$}}}}%
  \def\checkmark{\popcheck{popdark}}%
  \def\star{\mathbin{\ast}}\def\bigstar{\mathord{\ast}}%
  \def\diamond{\mathbin{\scalebox{0.6}{$\Diamond$}}}%
  \def\bigcup{\mathop{\mathchoice{\raisebox{-0.3em}{\scalebox{1.7}{$\cup$}}}{\scalebox{1.25}{$\cup$}}{\cup}{\cup}}\displaylimits}%
  \def\bigcap{\mathop{\mathchoice{\raisebox{-0.3em}{\scalebox{1.7}{$\cap$}}}{\scalebox{1.25}{$\cap$}}{\cap}{\cap}}\displaylimits}%
  \def\lesseqqgtr{\mathrel{\lessgtr}}\def\bigoplus{\mathop{\oplus}\displaylimits}%
}
\providecommand{\ssi}{si et seulement si} % abréviation fréquente
\AtBeginDocument{\providecommand{\wideparen}{\overparen}} % arc de cercle

%% FIGFIX-BEGIN v5 — correctifs globaux des figures (docs/onbuch-generation/tools/figfix)
\usepackage{adjustbox}\usepackage{etoolbox}\tcbuselibrary{raster}
% toute figure TikZ/pgfplots plus large que la ligne est réduite à la largeur disponible
\BeforeBeginEnvironment{tikzpicture}{\begin{adjustbox}{max width=\linewidth}}
\AfterEndEnvironment{tikzpicture}{\end{adjustbox}}
% légendes pgfplots lisibles par-dessus les courbes
\pgfplotsset{every axis/.append style={legend style={fill=white,fill opacity=0.92,text opacity=1,draw=black!15,inner sep=3pt}}}
% étiquettes d'arêtes (syntaxe edge["..."]) avec fond blanc
\tikzset{every edge quotes/.append style={fill=white,inner sep=1.2pt}}
%% FIGFIX-END
\begin{document}

\pagegarde{Les rites de veuvage}{Allemand}{Tle A}{Chapitre 1 : Vie familiale et sociale}{AL03}{2 h}{Cette leçon de type texte (Tle A) exploite un texte sur les rites de veuvage pour construire le vocabulaire du deuil et des traditions, maîtriser le passif (Vorgangspassiv), les propositions relatives et le prétérit des verbes forts. Elle aboutit à une production guidée de comparaison interculturelle Cameroun–pays germanophones.
\vspace{4pt}
\begin{multicols}{2}\small
\begin{itemize}
\item À la fin de cette leçon, je sais comprendre globalement puis en détail un texte allemand sur les rites de veuvage.
\item À la fin de cette leçon, je sais employer le vocabulaire thématique du deuil et des traditions (die Witwe, der Brauch, die Trauer...) avec articles et pluriels corrects.
\item À la fin de cette leçon, je sais former et utiliser le Vorgangspassiv (présent, prétérit, parfait) pour décrire les étapes d'un rite.
\item À la fin de cette leçon, je sais construire des propositions relatives avec les pronoms relatifs déclinés.
\item À la fin de cette leçon, je sais conjuguer les verbes forts au prétérit et former les participes passés (avec et sans ge-).
\item À la fin de cette leçon, je sais répondre à des questions de compréhension et rédiger un résumé structuré du texte.
\end{itemize}\end{multicols}}

\begin{sommaire}
\begin{multicols}{2}
\begin{enumerate}
\item Texte support : « Witwenrituale in Kamerun und in Deutschland »
\item Compréhension du texte
\item Lexique thématique : deuil, veuvage et traditions
\item Grammaire 1 : le passif (Vorgangspassiv) pour décrire un rite
\item Grammaire 2 : propositions relatives, prétérit des verbes forts et participes
\item Méthode du commentaire, commentaire modèle et production guidée
\item Activité d'intégration
\item Méthodes et exercices résolus
\item Exercices
\item Corrigés des exercices
\item Fiche bilan
\end{enumerate}
\end{multicols}
\end{sommaire}

\begin{objectifs}
\begin{itemize}
\item À la fin de cette leçon, je sais comprendre globalement puis en détail un texte allemand sur les rites de veuvage.
\item À la fin de cette leçon, je sais employer le vocabulaire thématique du deuil et des traditions (die Witwe, der Brauch, die Trauer...) avec articles et pluriels corrects.
\item À la fin de cette leçon, je sais former et utiliser le Vorgangspassiv (présent, prétérit, parfait) pour décrire les étapes d'un rite.
\item À la fin de cette leçon, je sais construire des propositions relatives avec les pronoms relatifs déclinés.
\item À la fin de cette leçon, je sais conjuguer les verbes forts au prétérit et former les participes passés (avec et sans ge-).
\item À la fin de cette leçon, je sais répondre à des questions de compréhension et rédiger un résumé structuré du texte.
\item À la fin de cette leçon, je sais rédiger un court paragraphe comparant un rite de veuvage camerounais et un rite germanophone.
\item À la fin de cette leçon, je sais commenter un texte documentaire selon la méthode en trois temps (présentation, analyse, bilan).
\end{itemize}
\end{objectifs}

\begin{prerequis}
\begin{itemize}
\item Le prétérit des verbes faibles et des auxiliaires (haben, sein, werden) vu en Première
\item La formation du participe passé des verbes faibles (gemacht, gearbeitet)
\item La déclinaison des articles définis et indéfinis (der/die/das, Nominatif–Accusatif–Datif)
\item La phrase déclarative et subordonnée simple (weil, dass) avec place du verbe
\item Le vocabulaire de la famille et de la vie sociale (die Familie, die Hochzeit, der Tod)
\end{itemize}
\end{prerequis}

\newpage

\coursec{Texte support : « Witwenrituale in Kamerun und in Deutschland »}

\courssub{Situation d'accroche et problématique}

Au Cameroun, quand un chef de famille décède, la veuve observe souvent des rites précis : port de vêtements particuliers, période de deuil, cérémonies de purification, parfois confinement dans la maison pendant plusieurs mois. À Yaoundé comme dans les villages de l'Ouest ou du Nord-Ouest, ces coutumes structurent la vie de la famille endeuillée. En Allemagne aussi, le veuvage est marqué par des traditions : la \all{Trauerzeit} (période de deuil), les vêtements sombres, la messe d'anniversaire du décès (\all{Todestagsmesse}) en Autriche catholique, l'annonce de décès dans le journal (\all{Todesanzeige}).

\textbf{Problématique :} Comment parler de ces coutumes en allemand ? Comment décrire un rite avec le passif (\all{Vorgangspassiv}) et raconter ce qui s'est passé au prétérit ? Quelles propositions relatives permettent de préciser qui fait quoi dans la cérémonie ?

Cette leçon nous donne les mots et les structures pour comparer les rites de veuvage au Cameroun et dans les pays germanophones, avec respect et précision linguistique.

\courssub{Première lecture globale : survol du texte}

\begin{methode}[Lire un texte documentaire en trois passes]
\begin{enumerate}
\item \textbf{Survol} (30 secondes) : lire le titre, la première et la dernière phrase. Repérer le thème général, le narrateur, les lieux. Ne pas s'arrêter sur les mots inconnus.
\item \textbf{Lecture avec glossaire} : lire phrase par phrase en s'aidant du glossaire. Souligner les mots du champ lexical du deuil. Comprendre le sens global de chaque paragraphe.
\item \textbf{Relevé des structures} : chercher les formes grammaticales cibles de la leçon — ici le passif (\all{wurde geschnitten}), les propositions relatives (\all{die bei uns beobachtet werden}) et les verbes forts au prétérit (\all{starb}).
\end{enumerate}
\end{methode}

Avant de lire, observez le titre : \all{Witwenrituale in Kamerun und in Deutschland}. Deux questions guident la première lecture : \emph{Qui parle ?} et \emph{Quels sont les deux espaces comparés ?}

\courssub{Le texte support}

\begin{textebox}[Witwenrituale in Kamerun und in Deutschland — Erzählung einer Witwe aus Bafoussam]
Als mein Mann vor zwei Jahren starb, wurde ich von den älteren Frauen des Dorfes empfangen. Sie nahmen mich in ein kleines Haus hinter der Familienkonzession, wo ich sieben Tage bleiben musste. Zuerst wurden meine Haare geschnitten; danach wurde mir ein weißes Tuch gegeben, das ich um die Hüften tragen sollte. Während der Trauerzeit durfte ich das Haus nicht verlassen und kein Fest besuchen.

Die Verwandten meines Mannes kamen jeden Tag. Essen wurde von den Nachbarinnen gekocht, und die Kinder wurden von meiner Schwester betreut. Viele Bräuche wurden beobachtet: Am Morgen wurde laut geweint, am Abend wurden Gebete gesprochen. Man sagte, diese Riten zeigen den Respekt vor dem Verstorbenen und schützen die Familie.

Nach einem Jahr wurde die Trauerzeit offiziell beendet. Bei einer großen Zeremonie wurden die schwarzen und weißen Kleider abgelegt, und ich durfte wieder bunte Stoffe tragen. Erst danach wurde das Erbe meines Mannes unter den Kindern geteilt.

In Deutschland, wo meine Tochter jetzt lebt, ist das anders. Dort trägt eine Witwe meistens nur zur Beerdigung Schwarz; die Trauerzeit wird nicht so streng befolgt. Die Todesanzeige wird in der Zeitung veröffentlicht, und nach einigen Wochen geht das Leben weiter. Die Riten, die bei uns beobachtet werden, sind strenger, aber sie geben der Witwe auch den Schutz und die Achtung der ganzen Gemeinschaft.
\end{textebox}

\textbf{Traduction du texte.} « Quand mon mari est mort il y a deux ans, j'ai été accueillie par les femmes âgées du village. Elles m'ont emmenée dans une petite maison derrière la concession familiale, où je devais rester sept jours. D'abord mes cheveux ont été coupés ; ensuite on m'a donné un pagne blanc que je devais porter autour des hanches. Pendant la période de deuil, je n'avais pas le droit de quitter la maison ni d'assister à une fête.

Les parents de mon mari venaient chaque jour. Les repas étaient préparés par les voisines, et les enfants étaient gardés par ma sœur. Beaucoup de coutumes étaient observées : le matin on pleurait à haute voix, le soir on récitait des prières. On disait que ces rites montrent le respect envers le défunt et protègent la famille.

Au bout d'un an, la période de deuil a été officiellement terminée. Lors d'une grande cérémonie, les vêtements noirs et blancs ont été retirés, et j'ai pu porter à nouveau des tissus colorés. Ce n'est qu'ensuite que l'héritage de mon mari a été partagé entre les enfants.

En Allemagne, où ma fille vit maintenant, c'est différent. Là-bas, une veuve ne porte le noir le plus souvent qu'à l'enterrement ; la période de deuil n'est pas suivie de manière aussi stricte. L'avis de décès est publié dans le journal, et après quelques semaines la vie continue. Les rites qui sont observés chez nous sont plus stricts, mais ils donnent aussi à la veuve la protection et la considération de toute la communauté. »

\courssub{Glossaire du texte}

\begin{center}
\begin{tabularx}{\linewidth}{|l|X|X|}
\hline
\rowcolor{popblueL} \textbf{Deutsch} & \textbf{Pluriel / formes} & \textbf{Français} \\
\hline
die Witwe & die Witwen & la veuve \\
\hline
der Witwer & die Witwer & le veuf \\
\hline
der Brauch & die Bräuche & la coutume, l'usage \\
\hline
die Trauer & (sans pluriel) & le deuil, la tristesse \\
\hline
das Ritual & die Rituale & le rite, le rituel \\
\hline
die Trauerzeit & die Trauerzeiten & la période de deuil \\
\hline
die Beerdigung & die Beerdigungen & l'enterrement, les funérailles \\
\hline
der Verstorbene & die Verstorbenen & le défunt, la défunte \\
\hline
das Erbe & (sans pluriel) & l'héritage \\
\hline
die Todesanzeige & die Todesanzeigen & l'avis de décès (faire-part) \\
\hline
beerdigen & beerdigte, hat beerdigt & enterrer, inhumer \\
\hline
erben & erbte, hat geerbt & hériter \\
\hline
schneiden (die Haare) & schnitt, hat geschnitten & couper (les cheveux) \\
\hline
tragen (Trauerkleidung) & trug, hat getragen & porter (des vêtements de deuil) \\
\hline
betreuen & betreute, hat betreut & garder, s'occuper de \\
\hline
teilen & teilte, hat geteilt & partager \\
\hline
veröffentlichen & veröffentlichte, hat veröffentlicht & publier \\
\hline
\end{tabularx}
\end{center}

\begin{definition}[der Brauch, die Bräuche]
\all{Der Brauch} (pluriel \all{die Bräuche}) signifie « la coutume, l'usage traditionnel ». Exemple : \all{Bei uns ist es ein alter Brauch, ein Jahr lang Trauerkleidung zu tragen.} « Chez nous, c'est une vieille coutume de porter des vêtements de deuil pendant un an. » Attention au pluriel avec Umlaut : \all{Brauch} devient \all{Bräuche}.
\end{definition}

\begin{attention}[Faux ami : der Brauch / brauchen]
Ne confondez pas \all{der Brauch} (la coutume, un nom) avec le verbe \all{brauchen} (avoir besoin) ! \all{Ich brauche ein Tuch.} « J'ai besoin d'un pagne. » — \all{Das ist ein alter Brauch.} « C'est une vieille coutume. » Les francophones traduisent aussi parfois « coutume » par \all{die Kostüm} : erreur ! \all{Das Kostüm} signifie « le costume, le déguisement ». La coutume = \all{der Brauch} ou \all{die Sitte}.
\end{attention}

\begin{attention}[die Witwe / der Witwer : le -r du masculin]
\all{Die Witwe} (la veuve) est féminin ; le masculin est \all{der Witwer} (le veuf), avec un \textbf{-r} final. Les francophones oublient souvent ce -r et disent \all{der Witwe}, ce qui est faux. Retenez la paire : \all{die Witwe — der Witwer}, comme \all{die Witwe} + \textbf{r}.
\end{attention}

\courssub{Deuxième lecture détaillée : le champ lexical du deuil}

Soulignons dans le texte tous les mots qui appartiennent au champ lexical du deuil et du veuvage. On peut les classer en quatre familles :

\begin{center}
\begin{tabularx}{\linewidth}{|l|X|X|}
\hline
\rowcolor{popblueL} \textbf{Famille} & \textbf{Mots du texte} & \textbf{Français} \\
\hline
Personnes & die Witwe, der Witwer, der Verstorbene, die Verwandten & veuve, veuf, défunt, parents \\
\hline
Rites et actions & Haare schneiden, Trauerkleidung tragen, Gebete sprechen, das Erbe teilen & couper les cheveux, porter le deuil, prier, partager l'héritage \\
\hline
Temps & die Trauerzeit, sieben Tage, nach einem Jahr & période de deuil, sept jours, au bout d'un an \\
\hline
Sentiments & die Trauer, der Schmerz, der Respekt, die Achtung & deuil, douleur, respect, considération \\
\hline
\end{tabularx}
\end{center}

\begin{popfigure}
\begin{center}
\begin{center}\tcbox[colback=popblue,colframe=popblue,coltext=white,arc=9pt,boxsep=4pt,left=6pt,right=6pt]{\bfseries\small das Witwensein (le veuvage)}\end{center}
\vspace{-2pt}
\begin{tcbraster}[raster columns=2,raster equal height=rows,raster column skip=6pt,raster row skip=6pt,enhanced,blankest]
\begin{tcolorbox}[enhanced,colback=white,colframe=poporange,boxrule=0.9pt,arc=3pt,title={Personen},fonttitle=\bfseries\scriptsize,coltitle=white,colbacktitle=poporange,left=4pt,right=4pt,top=2pt,bottom=2pt,boxsep=1pt,toptitle=1.5pt,bottomtitle=1.5pt,halign=left]
\begin{itemize}[nosep,leftmargin=*,itemsep=1pt,font=\scriptsize]
\item die Witwe der Witwer
\item die Verwandten der Verstorbene
\end{itemize}
\end{tcolorbox}
\begin{tcolorbox}[enhanced,colback=white,colframe=popgreen,boxrule=0.9pt,arc=3pt,title={Riten},fonttitle=\bfseries\scriptsize,coltitle=white,colbacktitle=popgreen,left=4pt,right=4pt,top=2pt,bottom=2pt,boxsep=1pt,toptitle=1.5pt,bottomtitle=1.5pt,halign=left]
\begin{itemize}[nosep,leftmargin=*,itemsep=1pt,font=\scriptsize]
\item Haare schneiden
\item Trauerkleidung tragen
\end{itemize}
\end{tcolorbox}
\begin{tcolorbox}[enhanced,colback=white,colframe=poppurple,boxrule=0.9pt,arc=3pt,title={Zeit},fonttitle=\bfseries\scriptsize,coltitle=white,colbacktitle=poppurple,left=4pt,right=4pt,top=2pt,bottom=2pt,boxsep=1pt,toptitle=1.5pt,bottomtitle=1.5pt,halign=left]
\begin{itemize}[nosep,leftmargin=*,itemsep=1pt,font=\scriptsize]
\item die Trauerzeit
\item ein Jahr
\end{itemize}
\end{tcolorbox}
\begin{tcolorbox}[enhanced,colback=white,colframe=poppink,boxrule=0.9pt,arc=3pt,title={Gefühle},fonttitle=\bfseries\scriptsize,coltitle=white,colbacktitle=poppink,left=4pt,right=4pt,top=2pt,bottom=2pt,boxsep=1pt,toptitle=1.5pt,bottomtitle=1.5pt,halign=left]
\begin{itemize}[nosep,leftmargin=*,itemsep=1pt,font=\scriptsize]
\item die Trauer
\item der Schmerz
\end{itemize}
\end{tcolorbox}
\end{tcbraster}

\end{center}
\legende{Carte mentale du champ lexical du veuvage : quatre branches (personnes, rites, temps, sentiments) à mémoriser avec articles et pluriels.}
\end{popfigure}

\courssub{Repérage des marqueurs chronologiques du rite}

Un rite se raconte dans un ordre précis. Le texte utilise des \cle{marqueurs chronologiques} qu'il faut savoir employer pour décrire une cérémonie :

\begin{center}
\begin{tabularx}{\linewidth}{|l|X|X|}
\hline
\rowcolor{popblueL} \textbf{Marqueur} & \textbf{Français} & \textbf{Exemple tiré ou inspiré du texte} \\
\hline
zuerst & d'abord & \all{Zuerst wurden meine Haare geschnitten.} \\
\hline
danach & ensuite & \all{Danach wurde mir ein weißes Tuch gegeben.} \\
\hline
während (+ génitif) & pendant & \all{Während der Trauerzeit durfte ich das Haus nicht verlassen.} \\
\hline
nach einem Jahr & au bout d'un an & \all{Nach einem Jahr wurde die Trauerzeit beendet.} \\
\hline
schließlich / erst danach & finalement / seulement ensuite & \all{Erst danach wurde das Erbe geteilt.} \\
\hline
\end{tabularx}
\end{center}

\begin{aretenir}[Ordre des mots avec les marqueurs]
Quand la phrase commence par un marqueur chronologique (\all{zuerst}, \all{danach}, \all{nach einem Jahr}), le verbe conjugué vient en \textbf{deuxième position} et le sujet \textbf{après} le verbe : \all{Nach einem Jahr \textbf{wurde} die Trauerzeit \textbf{beendet}.} Ne dites jamais \all{Nach einem Jahr die Trauerzeit wurde beendet} : c'est le piège numéro un des francophones.
\end{aretenir}

\courssub{Questions de compréhension globale}

\begin{exoresolu}[Compréhension globale du texte]
\textbf{Énoncé.} Répondez en allemand par des phrases complètes :
\begin{enumerate}
\item Wer erzählt in diesem Text ?
\item Wie lange dauerte die Trauerzeit der Erzählerin ?
\item Was wurde zuerst mit den Haaren der Witwe gemacht ?
\item Was ist in Deutschland anders ?
\end{enumerate}
\tcblower
\textbf{Solution détaillée.}
\begin{enumerate}
\item \all{Eine Witwe aus Bafoussam erzählt; ihr Mann ist vor zwei Jahren gestorben.} (Une veuve de Bafoussam raconte ; son mari est mort il y a deux ans.)
\item \all{Die Trauerzeit dauerte ein Jahr.} (La période de deuil a duré un an.) — On reprend le marqueur \all{nach einem Jahr} du texte.
\item \all{Zuerst wurden ihre Haare geschnitten.} (D'abord, ses cheveux ont été coupés.) — Réponse au passif, comme dans le texte.
\item \all{In Deutschland trägt eine Witwe meistens nur zur Beerdigung Schwarz, und die Trauerzeit wird nicht so streng befolgt.} (En Allemagne, une veuve ne porte le noir le plus souvent qu'à l'enterrement, et la période de deuil n'est pas suivie de manière aussi stricte.)
\end{enumerate}
Méthode : pour chaque question, repérez la phrase du texte qui contient la réponse, puis reformulez-la avec le bon ordre des mots (verbe en deuxième position).
\end{exoresolu}

\begin{exercice}[À vous de jouer — Compréhension et lexique]
\begin{enumerate}
\item Soulignez dans le texte trois autres formes du passif et traduisez-les.
\item Trouvez dans le texte deux marqueurs chronologiques et construisez avec chacun une phrase personnelle sur un rite de votre région.
\item Donnez l'article et le pluriel de : \all{Brauch}, \all{Ritual}, \all{Beerdigung}, \all{Witwe}.
\item Traduisez : « D'abord les cheveux de la veuve ont été coupés, ensuite l'héritage a été partagé. »
\end{enumerate}
\end{exercice}

\begin{corrige}[Exercice « À vous de jouer — Compréhension et lexique »]
\begin{enumerate}
\item Par exemple : \all{Essen wurde von den Nachbarinnen gekocht} « les repas étaient préparés par les voisines » ; \all{die Kinder wurden von meiner Schwester betreut} « les enfants étaient gardés par ma sœur » ; \all{am Abend wurden Gebete gesprochen} « le soir, des prières étaient récitées ».
\item Exemples : \all{Zuerst wurde die Familie informiert.} « D'abord la famille a été informée. » \all{Schließlich wurde ein Fest vorbereitet.} « Finalement, une fête a été préparée. »
\item \all{der Brauch, die Bräuche} ; \all{das Ritual, die Rituale} ; \all{die Beerdigung, die Beerdigungen} ; \all{die Witwe, die Witwen}.
\item \all{Zuerst wurden die Haare der Witwe geschnitten, danach wurde das Erbe geteilt.} Attention : verbe en deuxième position après \all{zuerst} et \all{danach}, participe passé en fin de phrase.
\end{enumerate}
\end{corrige}

\coursec{Grammaire 1 : Le passif (Vorgangspassiv) pour décrire un rite}

\courssub{Observation et règle de formation}

Dans le texte, l'accent est mis sur ce qui \emph{arrive} à la veuve, non sur qui accomplit l'action. C'est la fonction du \cle{Vorgangspassiv} (passif de procès / passif d'action).

\begin{propriete}[Formation du Vorgangspassiv]
\textbf{Structure :} \cle{werden} (conjugué au temps voulu) + \cle{Partizip II} (participe passé) en position finale.

\begin{center}
\begin{tabular}{|l|l|l|l|}
\hline
\rowcolor{popblueL} \textbf{Temps} & \textbf{Allemand} & \textbf{Français} & \textbf{Exemple du texte} \\
\hline
Présent & \all{werden} + Partizip II & est / sont + participe & \all{Die Todesanzeige \textbf{wird veröffentlicht}.} \\
\hline
Prétérit & \all{wurden} + Partizip II & a/ont été + participe & \all{Die Haare \textbf{wurden geschnitten}.} \\
\hline
Parfait & \all{sind geworden} + Partizip II & ont été + participe & \all{Die Haare \textbf{sind geschnitten worden}.} \\
\hline
Futur & \all{werden} + Partizip II + \all{werden} & seront + participe & \all{Die Haare \textbf{werden geschnitten werden}.} \\
\hline
\end{tabular}
\end{center}

\textbf{Règle d'accord :} \all{werden} s'accorde avec le \textbf{sujet patient} (celui qui subit l'action).
\begin{itemize}
\item Singulier : \all{Das Tuch \textbf{wurde} gegeben.} / \all{Die Trauerzeit \textbf{wurde} beendet.}
\item Pluriel : \all{Die Haare \textbf{wurden} geschnitten.} / \all{Die Riten \textbf{werden} beobachtet.}
\end{itemize}
\end{propriete}

\begin{attention}[Passif impersonnel vs personnel]
\begin{itemize}
\item \textbf{Passif personnel} (sujet = objet actif) : \all{Man schneidet die Haare.} $\rightarrow$ \all{Die Haare werden geschnitten.}
\item \textbf{Passif impersonnel} (verbe intransitif ou pas d'objet) : \all{Man weint.} $\rightarrow$ \all{Es wird geweint.} (Il est pleuré / On pleure). Le sujet formel est \all{es} (ou omission en position 1 : \all{Am Morgen wurde laut geweint.}).
\end{itemize}
Au Baccalauréat, sachez transformer une phrase active en passif personnel \textbf{et} reconnaître le passif impersonnel.
\end{attention}

\courssub{Le complément d'agent : von + datif}

Si l'on veut mentionner l'auteur de l'action, on utilise \all{von} + \textbf{datif}.

\begin{exemplebox}[Complément d'agent]
\begin{itemize}
\item \all{Essen wurde \textbf{von den Nachbarinnen} gekocht.} (par les voisines — datif pluriel : \all{den Nachbarinnen})
\item \all{Die Kinder wurden \textbf{von meiner Schwester} betreut.} (par ma sœur — datif singulier : \all{meiner Schwester})
\item \all{Ich wurde \textbf{von den älteren Frauen} empfangen.} (par les femmes âgées)
\end{itemize}
\emph{Attention :} \all{durch} + accusatif s'utilise pour l'instrument ou la cause intermédiaire (\all{Das Feuer wurde durch Blitzschlag verursacht}), mais pour l'agent humain, c'est \all{von}.
\end{exemplebox}

\courssub{Exercices d'application : Passif}

\begin{exercice}[Transformation active $\rightarrow$ passif (Prétérit)]
Transformez les phrases actives en passif (Prétérit). Mettez le complément d'agent entre parenthèses avec \all{von} + datif.
\begin{enumerate}
\item Die älteren Frauen empfangen die Witwe. (von den älteren Frauen)
\item Die Nachbarinnen kochen das Essen. (von den Nachbarinnen)
\item Meine Schwester betreut die Kinder. (von meiner Schwester)
\item Man schneidet der Witwe die Haare. (Achtung: deux objets !)
\item Man gibt ihr ein weißes Tuch. (von den Frauen)
\end{enumerate}
\end{exercice}

\begin{corrige}[Exercice « Transformation active $\rightarrow$ passif »]
\begin{enumerate}
\item \all{Die Witwe wurde von den älteren Frauen empfangen.}
\item \all{Das Essen wurde von den Nachbarinnen gekocht.}
\item \all{Die Kinder wurden von meiner Schwester betreut.}
\item \all{Der Witwe wurden die Haare geschnitten.} (Passif avec double objet : l'objet datif \all{der Witwe} reste datif, l'objet accusatif \all{die Haare} devient sujet nominatif pluriel $\rightarrow$ \all{wurden}). \emph{Variante impersonnelle :} \all{Es wurde der Witwe die Haare geschnitten.}
\item \all{Ein weißes Tuch wurde ihr von den Frauen gegeben.} (Sujet : \all{ein weißes Tuch} ; \all{ihr} = datif inchangé).
\end{enumerate}
\end{corrige}

\begin{exercice}[Passif au présent et au parfait]
Conjuguez le verbe entre parenthèses au passif (temps indiqué).
\begin{enumerate}
\item Die Todesanzeige \_\_\_\_\_\_\_\_\_\_ in der Zeitung \textbf{veröffentlichen} (Présent).
\item Die Trauerzeit \_\_\_\_\_\_\_\_\_\_ offiziell \textbf{beenden} (Prétérit).
\item Das Erbe \_\_\_\_\_\_\_\_\_\_ unter den Kindern \textbf{teilen} (Parfait).
\item Bei der Zeremonie \_\_\_\_\_\_\_\_\_\_ die Kleider \_\_\_\_\_\_\_\_\_\_ \textbf{ablegen} (Prétérit, verbe à particule).
\end{enumerate}
\end{exercice}

\begin{corrige}[Exercice « Passif présent / prétérit / parfait »]
\begin{enumerate}
\item \all{wird veröffentlicht}
\item \all{wurde beendet}
\item \all{ist geteilt worden} (Partizip II de \all{werden} = \all{geworden}, mais au passif on omet \all{ge-} : \all{worden})
\item \all{wurden ... abgelegt} (verbe séparable : \all{ab|legen} $\rightarrow$ Partizip II \all{abgelegt})
\end{enumerate}
\end{corrige}

\begin{aretenir}[Bilan : Le passif pour décrire un rite]
\begin{itemize}
\item Utilisez le \textbf{Prétérit passif} (\all{wurden} + Partizip II) pour raconter un rite passé : \all{Zuerst wurden die Haare geschnitten.}
\item Utilisez le \textbf{Présent passif} (\all{werden} + Partizip II) pour décrire une coutume générale : \all{Die Todesanzeige wird veröffentlicht.}
\item L'agent humain : \all{von} + \textbf{Datif}.
\item Ordre des mots : Marqueur (Zuerst) — \textbf{Verbe conjugué} — Sujet — ... — \textbf{Partizip II}.
\end{itemize}
\end{aretenir}

\coursec{Grammaire 2 : Propositions relatives, Prétérit des verbes forts, Participes passés}

\courssub{Les propositions relatives (Relativsätze)}

Une proposition relative précise un antécédent (nom ou pronom). En allemand, le verbe conjugué va à la \textbf{fin} de la subordonnée.

\begin{propriete}[Le pronom relatif : déclinaison et emploi]
Le pronom relatif s'accorde en \textbf{genre} et \textbf{nombre} avec son antécédent. Son \textbf{cas} dépend de sa fonction \emph{dans la relative}.

\begin{center}
\begin{tabular}{|l|c|c|c|c|}
\hline
\rowcolor{popblueL} \textbf{Cas} & \textbf{Masculin} & \textbf{Féminin} & \textbf{Neutre} & \textbf{Pluriel} \\
\hline
Nominatif & der & die & das & die \\
\hline
Accusatif & den & die & das & die \\
\hline
Datif & dem & der & dem & \textbf{denen} \\
\hline
Génitif & dessen & deren & dessen & deren \\
\hline
\end{tabular}
\end{center}

\end{propriete}

\begin{exemplebox}[Exemples du texte et variations]
\begin{itemize}
\item \textbf{Nominatif (sujet) :} \all{Die Riten, \textbf{die} bei uns beobachtet werden, sind strenger.} (\all{die Riten} f. pl. $\rightarrow$ \all{die} ; fonction sujet dans la relative).
\item \textbf{Accusatif (COD) :} \all{Ein weißes Tuch, \textbf{das} ich um die Hüften tragen \textbf{sollte}, ...} (\all{das Tuch} n. sg. $\rightarrow$ \all{das} ; fonction COD de \all{tragen}).
\item \textbf{Datif (COI / préposition) :} \all{Die Frau, \textbf{der} ich das Tuch gab, war alt.} (\all{die Frau} f. sg. $\rightarrow$ \all{der} ; fonction COI de \all{geben}).
\item \textbf{Génitif (possession) :} \all{Der Mann, \textbf{dessen} Frau starb, weinte.} (\all{der Mann} m. sg. $\rightarrow$ \all{dessen}).
\item \textbf{Préposition + pronom relatif :} \all{Das Haus, \textbf{in dem} ich blieb, war klein.} (Préposition \all{in} + datif \all{dem}).
\end{itemize}
\end{exemplebox}

\begin{attention}[Pièges fréquents pour francophones]
\begin{enumerate}
\item \textbf{Oubli du verbe à la fin} : \all{*Die Riten, die sind streng, ...} $\rightarrow$ Correct : \all{Die Riten, die \textbf{streng sind}, ...}
\item \textbf{Mauvais cas} : \all{*Der Mann, der ich sah} (français « que j'ai vu ») $\rightarrow$ Correct : \all{Der Mann, \textbf{den} ich sah} (Accusatif).
\item \textbf{Préposition séparée} : En français « la maison où j'habite ». En allemand, la préposition \textbf{doit} précéder le pronom : \all{Das Haus, \textbf{in dem} ich wohne.} (Pas \all{*wo ich wohne in}).
\item \textbf{wo / wohin / woher} : Utilisés seulement pour des \textbf{lieux} sans préposition explicite (ou remplaçant \all{in dem}, \all{in das}, \all{aus dem}). Exemple texte : \all{... ein Haus, \textbf{wo} ich bleiben musste.} (Correct, familier/standard pour \all{in dem}).
\end{enumerate}
\end{attention}

\courssub{Le prétérit (Präteritum) des verbes forts et mixtes}

Le prétérit (narratif) est le temps du récit au passé. Les verbes forts changent de voyelle radicale (Ablaut) et n'ajoutent pas de \all{-te}.

\begin{propriete}[Tableau des verbes forts / mixtes essentiels pour la leçon]
\begin{center}
\begin{tabularx}{\linewidth}{|l|l|l|l|X|}
\hline
\rowcolor{popblueL} \textbf{Infinitif} & \textbf{Präteritum (ich/er/sie/es)} & \textbf{Partizip II} & \textbf{Auxiliaire} & \textbf{Français} \\
\hline
sterben & starb & \textbf{ist} gestorben & sein & mourir \\
\hline
schneiden & schnitt & \textbf{hat} geschnitten & haben & couper \\
\hline
tragen & trug & \textbf{hat} getragen & haben & porter \\
\hline
geben & gab & \textbf{hat} gegeben & haben & donner \\
\hline
kommen & kam & \textbf{ist} gekommen & sein & venir \\
\hline
sprechen & sprach & \textbf{hat} gesprochen & haben & parler \\
\hline
bleiben & blieb & \textbf{ist} geblieben & sein & rester \\
\hline
verlassen & verließ & \textbf{hat} verlassen & haben & quitter \\
\hline
finden & fand & \textbf{hat} gefunden & haben & trouver \\
\hline
halten & hielt & \textbf{hat} gehalten & haben & tenir, s'arrêter \\
\hline
lassen & ließ & \textbf{hat} gelassen & haben & laisser \\
\hline
heißen & hieß & \textbf{hat} geheißen & haben & s'appeler, ordonner \\
\hline
\end{tabularx}
\end{center}

\emph{Rappel :} Verbes à auxiliaire \all{sein} = mouvement vers un lieu / changement d'état (\all{sterben, kommen, bleiben, gehen, fallen, werden}). Les autres \all{haben}.
\end{propriete}

\begin{attention}[Verbes mixtes : voyelle changée + -te]
Certains verbes changent de voyelle \textbf{et} prennent \all{-te} : \all{denken $\rightarrow$ dachte, hat gedacht} ; \all{kennen $\rightarrow$ kannte, hat gekannt} ; \all{brennen $\rightarrow$ brannte, hat gebrannt} ; \all{bringen $\rightarrow$ brachte, hat gebracht} ; \all{wissen $\rightarrow$ wusste, hat gewusst}. Apprenez-les par cœur.
\end{attention}

\courssub{Le participe passé (Partizip II) : formation complète}

Le Partizip II est indispensable pour le Passif (\all{werden} + Partizip II), le Parfait (\all{haben/sein} + Partizip II) et le Plus-que-parfait.

\begin{propriete}[Formation du Partizip II]
\begin{center}
\begin{tabularx}{\linewidth}{|l|l|X|}
\hline
\rowcolor{popblueL} \textbf{Type} & \textbf{Règle} & \textbf{Exemples} \\
\hline
Faibles (réguliers) & \textbf{ge-} + radical + \textbf{-t} & \all{kochen $\rightarrow$ gekocht}, \all{betreuen $\rightarrow$ betreut}, \all{veröffentlichen $\rightarrow$ veröffentlicht} \\
\hline
Forts (irréguliers) & \textbf{ge-} + radical modifié + \textbf{-en} & \all{schneiden $\rightarrow$ geschnitten}, \all{geben $\rightarrow$ gegeben}, \all{sprechen $\rightarrow$ gesprochen}, \all{sterben $\rightarrow$ gestorben} \\
\hline
Mixtes & \textbf{ge-} + radical modifié + \textbf{-t} & \all{denken $\rightarrow$ gedacht}, \all{kennen $\rightarrow$ gekannt}, \all{bringen $\rightarrow$ gebracht} \\
\hline
Verbes en \textbf{-ieren} & \textbf{sans ge-} + \textbf{-t} & \all{studieren $\rightarrow$ studiert}, \all{telefonieren $\rightarrow$ telefoniert} \\
\hline
Verbes à préfixe \textbf{inséparable} (be-, er-, ver-, ent-, zer-, ge-, miss-) & \textbf{sans ge-} & \all{besuchen $\rightarrow$ besucht}, \all{verlassen $\rightarrow$ verlassen}, \all{erzählen $\rightarrow$ erzählt} \\
\hline
Verbes à particule \textbf{séparable} (ab-, an-, auf-, aus-, bei-, ein-, los-, mit-, nach-, vor-, weg-, zu-, zurück-) & \textbf{ge-} entre préfixe et radical & \all{ablegen $\rightarrow$ abgelegt}, \all{anfangen $\rightarrow$ angefangen}, \all{aufhören $\rightarrow$ aufgehört} \\
\hline
\end{tabularx}
\end{center}
\end{propriete}

\begin{exemplebox}[Participes passés du texte — Vérification]
\begin{itemize}
\item \all{empfangen} (séparable \all{emp-}? Non, \all{empfangen} est verbe fort à préfixe \all{emp-} non séparable $\rightarrow$ \textbf{empfangen}) — \checkmark
\item \all{geschnitten} (fort \all{schneiden} $\rightarrow$ \textbf{ge-schnitt-en}) — \checkmark
\item \all{gegeben} (fort \all{geben} $\rightarrow$ \textbf{ge-gab-en}) — \checkmark
\item \all{getragen} (fort \all{tragen} $\rightarrow$ \textbf{ge-trug-en}) — \checkmark
\item \all{gekocht} (faible \all{kochen} $\rightarrow$ \textbf{ge-koch-t}) — \checkmark
\item \all{betreut} (faible \all{betreuen} $\rightarrow$ \textbf{betreu-t}, pas de \all{ge-} car \all{be-} inséparable) — \checkmark
\item \all{beobachtet} (faible \all{beobachten}, \all{be-} inséparable $\rightarrow$ \textbf{beobachtet}) — \checkmark
\item \all{gesprochen} (fort \all{sprechen} $\rightarrow$ \textbf{ge-sproch-en}) — \checkmark
\item \all{beendet} (faible \all{beenden}, \all{be-} inséparable $\rightarrow$ \textbf{beendet}) — \checkmark
\item \all{abgelegt} (séparable \all{ab|legen} $\rightarrow$ \textbf{ab-ge-leg-t}) — \checkmark
\item \all{geteilt} (faible \all{teilen} $\rightarrow$ \textbf{ge-teilt}) — \checkmark
\item \all{veröffentlicht} (faible \all{veröffentlichen}, \all{ver-} inséparable $\rightarrow$ \textbf{veröffentlicht}) — \checkmark
\end{itemize}
\end{exemplebox}

\courssub{Exercices d'application : Grammaire 2}

\begin{exercice}[Propositions relatives : cas et prépositions]
Complétez par le pronom relatif correct (avec préposition si nécessaire).
\begin{enumerate}
\item Die Witwe, \_\_\_\_\_\_\_\_\_\_ der Mann starb, trägt Weiß. (Génitif)
\item Das Tuch, \_\_\_\_\_\_\_\_\_\_ sie trägt, ist weiß. (Accusatif)
\item Die Frauen, \_\_\_\_\_\_\_\_\_\_ sie empfangen wurde, sind alt. (Datif après \all{von} : \all{von \_\_\_\_\_\_\_\_\_\_})
\item Das Haus, \_\_\_\_\_\_\_\_\_\_ sie blieb, ist klein. (Préposition \all{in} + Datif)
\item Die Riten, \_\_\_\_\_\_\_\_\_\_ man beobachtet, sind alt. (Nominatif)
\end{enumerate}
\end{exercice}

\begin{corrige}[Exercice « Propositions relatives »]
\begin{enumerate}
\item \all{deren} (féminin singulier, génitif : \all{die Witwe, \textbf{deren} Mann ...})
\item \all{das} (neutre singulier, accusatif : \all{das Tuch, \textbf{das} sie trägt})
\item \all{denen} (pluriel, datif après \all{von} : \all{von \textbf{denen}})
\item \all{in dem} (ou \all{wo}) : \all{das Haus, \textbf{in dem} (/\textbf{wo}) sie blieb}
\item \all{die} (pluriel, nominatif : \all{die Riten, \textbf{die} man beobachtet})
\end{enumerate}
\end{corrige}

\begin{exercice}[Prétérit et Partizip II : formes fortes et mixtes]
Donnez le Prétérit (3e pers. sg.) et le Partizip II des verbes suivants.
\begin{enumerate}
\item finden
\item halten
\item bringen
\item heißen
\item sterben
\item verlassen
\end{enumerate}
\end{exercice}

\begin{corrige}[Exercice « Prétérit et Partizip II »]
\begin{enumerate}
\item \all{fand / hat gefunden}
\item \all{hielt / hat gehalten}
\item \all{brachte / hat gebracht}
\item \all{hieß / hat geheißen}
\item \all{starb / ist gestorben}
\item \all{verließ / hat verlassen}
\end{enumerate}
\end{corrige}

\coursec{Méthode du commentaire et production guidée}

\courssub{Méthode : Commenter un texte documentaire (Comparaison culturelle)}

\begin{methode}[Étapes du commentaire structuré (Écrit / Oral)]
\begin{enumerate}
\item \textbf{Introduction} (2--3 phrases) :
      \begin{itemize}
      \item Présenter le document : titre, auteur (ici : témoin), thème, pays concernés.
      \item Annoncer la problématique : \all{Der Text vergleicht die Witwenrituale in Kamerun und Deutschland.}
      \end{itemize}
\item \textbf{Résumé structuré} (4--5 phrases) :
      \begin{itemize}
      \item Utilisez les marqueurs chronologiques : \all{Zuerst ...}, \all{Danach ...}, \all{Nach einem Jahr ...}, \all{Schließlich ...}.
      \item Employez le \textbf{passif} pour les rites subis : \all{Die Haare wurden geschnitten.}
      \item Séparez clairement les deux parties : Cameroun (paragraphe 1--3) vs Allemagne (paragraphe 4).
      \end{itemize}
\item \textbf{Analyse thématique et linguistique} :
      \begin{itemize}
      \item \textbf{Thème} : Rigueur vs souplesse, protection communautaire vs individualisme, rôle des femmes/voisines.
      \item \textbf{Langue} : Repérez le vocabulaire du deuil, les verbes forts au prétérit, les propositions relatives (\all{die/wo/das}), le passif.
      \end{itemize}
\item \textbf{Conclusion / Ouverture} (1--2 phrases) :
      \begin{itemize}
      \item Avis personnel ou comparaison avec sa propre culture : \all{Meiner Meinung nach sind die kamerunischen Riten wichtig für den sozialen Zusammenhalt.}
      \end{itemize}
\end{enumerate}
\end{methode}

\courssub{Commentaire modèle (niveau Baccalauréat)}

\begin{textebox}[Commentaire modèle : « Witwenrituale — Ein Vergleich »]
Der Text „Witwenrituale in Kamerun und in Deutschland“ ist ein Erfahrungsbericht einer Witwe aus Bafoussam. Er vergleicht die traditionellen Riten in Kamerun mit der modernen Praxis in Deutschland.

In Kamerun ist die Trauerzeit sehr streng und dauert ein Jahr. Zuerst wurde die Witwe von den älteren Frauen empfangen und in ein separates Haus gebracht. Ihre Haare wurden geschnitten, und sie erhielt ein weißes Tuch. Während der Trauerzeit durfte sie das Haus nicht verlassen. Die Nachbarinnen kochten, die Verwandten besuchten sie täglich. Am Morgen wurde geweint, am Abend gebetet. Diese Riten, die von der Gemeinschaft beobachtet werden, zeigen Respekt vor dem Verstorbenen und schützen die Familie. Erst nach einem Jahr, bei einer großen Zeremonie, wurde die Trauerkleidung abgelegt und das Erbe geteilt.

In Deutschland hingegen ist alles anders. Die Witwe trägt Schwarz nur zur Beerdigung. Die Todesanzeige wird in der Zeitung veröffentlicht, und nach wenigen Wochen normalisiert sich der Alltag. Die Riten sind weniger streng, aber die Witwe verliert auch den Schutz der Gemeinschaft.

Meiner Meinung nach haben die kamerunischen Riten einen großen sozialen Wert: sie geben der Witwe Halt und Würde in einer schweren Zeit. Der deutsche Weg ist freier, aber auch einsamer.
\end{textebox}

\courssub{Production guidée : Écrit et Oral}

\begin{experience}[Production orale : « Rites de chez moi » — Jeu de rôle]
\textbf{Consigne :} En binôme. L'un est un journaliste allemand (\all{Reporter/in}), l'autre un témoin camerounais (\all{Zeuge/Zeugin}).
\begin{enumerate}
\item Le journaliste pose des questions sur un rite de veuvage (ou autre rite de passage : naissance, mariage, funérailles) de la région du témoin.
\item Le témoin répond en utilisant : \textbf{passif} (\all{wurden ... gemacht}), \textbf{marqueurs chronologiques} (\all{zuerst, danach, schließlich}), \textbf{relatives} (\all{die/der/das/wo}).
\item Exemple de début :
      \all{Reporter: „Können Sie mir ein Ritual aus Ihrer Region beschreiben?“} \\
      \all{Zeuge: „Bei uns im Norden wird zuerst das Grab vom Onkel gegraben. Danach wird der Leichnam ...“}
\end{enumerate}
\textbf{Critères d'évaluation :} Exactitude du passif (forme + ordre), richesse du vocabulaire (deuil, famille), fluidité (verbe en position 2, verbe à la fin en relative).
\end{experience}

\begin{exercice}[Production écrite : « Mon rite — Description au passif »]
Rédigez un paragraphe (80--100 mots) en allemand pour décrire un rite traditionnel de votre région (veuvage, naissance, initiation, funérailles).
\textbf{Contraintes obligatoires :}
\begin{enumerate}
\item Au moins \textbf{4 formes du passif} (Prétérit ou Présent).
\item Au moins \textbf{3 marqueurs chronologiques} (\all{zuerst, danach, dann, schließlich, nach ...}).
\item Au moins \textbf{2 propositions relatives} (avec \all{der/die/das}, \all{wo}, ou préposition+\all{dem/denen}).
\item Vocabulaire thématique : \all{der Brauch, die Trauerzeit, die Zeremonie, die Gemeinschaft, der Respekt} (ou adaptés à votre rite).
\end{enumerate}
\end{exercice}

\begin{corrige}[Exercice « Production écrite » — Proposition de correction type]
\textbf{Sujet :} Rite de funérailles chez les Bamiléké (exemple).
\all{Wenn ein Älterer stirbt, wird die Familie sofort informiert. Zuerst wird der Leichnam gewaschen und in ein schönes Tuch gewickelt. Danach wird er in einem speziellen Haus aufgebahrt, wo die Verwandten ihn besuchen kommen. Während der Trauerzeit werden Trommeln geschlagen und Opfertiere geschlachtet. Schließlich wird der Verstorbene auf dem Familiengrundstück beerdigt. Die Riten, die dabei beobachtet werden, stärken den Zusammenhalt der Gemeinschaft.}
\begin{itemize}
\item Passifs : \all{wird informiert, wird gewaschen, wird gewickelt, wird aufgebahrt, werden geschlagen, werden geschlachtet, wird beerdigt, werden beobachtet} (8 formes).
\item Marqueurs : \all{Zuerst, Danach, Während, Schließlich}.
\item Relatives : \all{wo die Verwandten ...}, \all{die dabei beobachtet werden}.
\end{itemize}
\end{corrige}

\begin{savaistu}[La Trauerzeit en Allemagne : hier et aujourd'hui]
En Allemagne, la \all{Trauerzeit} traditionnelle durait elle aussi un an : la veuve portait le \all{Witwenkleid} (la robe de veuve) noir et ne participait à aucune fête. On distinguait même la \all{tiefe Trauer} (deuil profond, six premiers mois) de la \all{halbe Trauer} (demi-deuil). Aujourd'hui, ces rites se sont fortement assouplis, surtout en ville : le noir n'est souvent porté que le jour de la \all{Beerdigung}. En revanche, la \all{Todesanzeige} dans le journal reste très courante, et en Autriche catholique, on célèbre encore la messe du \all{Todestag} (anniversaire du décès). Comparer ces traditions avec les rites camerounais est un excellent sujet de production orale à l'examen !
\end{savaistu}

\begin{aretenir}[Bilan final de la leçon AL03]
\begin{itemize}
\item \textbf{Vocabulaire clé} : \all{die Witwe / der Witwer}, \all{der Brauch (die Bräuche)}, \all{die Trauer}, \all{das Ritual}, \all{die Trauerzeit}, \all{die Beerdigung}, \all{der Verstorbene}, \all{das Erbe}, \all{die Todesanzeige}, \all{beerdigen}, \all{erben}, \all{die Haare schneiden}, \all{Trauerkleidung tragen}, \all{das Erbe teilen}.
\item \textbf{Passif (Vorgangspassiv)} : \all{werden} + Partizip II. Prétérit : \all{wurden} + Partizip II. Présent : \all{werden} + Partizip II. Agent : \all{von} + Datif. Ordre : Marqueur — Verbe — Sujet — ... — Partizip II.
\item \textbf{Propositions relatives} : Pronom relatif accordé en genre/nombre, cas selon fonction. Verbe à la fin. Préposition + pronom (\all{in dem, mit denen}). \all{wo} pour lieu.
\item \textbf{Prétérit forts} : \all{starb, schnitt, trug, gab, kam, sprach, blieb, verließ, fand, hielt, brachte, hieß}. Auxiliaires \all{sein/haben} selon verbe.
\item \textbf{Partizip II} : Faibles \all{ge-...-t}, Forts \all{ge-...-en}, Mixtes \all{ge-...-t}, \all{-ieren} \all{...-t}, Insep. \all{ohne ge-}, Séparables \all{Prefix-ge-Stamm-t/en}.
\item \textbf{Production} : Décrire un rite = Passif + Chronologie + Relatives + Vocabulaire précis.
\end{itemize}
\end{aretenir}

\coursec{Compréhension du texte}

Après la lecture approfondie du texte support \emph{« Witwenrituale in Kamerun und in Deutschland »} (section 1), nous passons maintenant à l'exploitation méthodique de ce document. La compréhension de l'écrit au baccalauréat ne se limite pas à répondre « oui » ou « non » : elle exige de \cle{repérer} l'information, de la \cle{reformuler} avec ses propres mots en allemand correct, et de \cle{justifier} chaque réponse par une citation précise du texte. Cette section vous guide pas à pas, de la compréhension globale à la production d'un résumé structuré, en passant par l'analyse lexicale et l'interprétation culturelle.

\courssub{Compréhension globale (Gesamtverständnis)}

Avant d'entrer dans le détail, il faut saisir la \cle{macrostructure} du texte : de quoi parle-t-il ? Qui s'exprime ? Quels espaces culturels sont mis en relation ?

\begin{exemplebox}[Questions types de compréhension globale]
\begin{enumerate}
    \item \textbf{Thema :} \all{Worum geht es in dem Text?} \textit{(De quoi parle le texte ?)}
    \item \textbf{Textsorte :} \all{Um was für einen Text handelt es sich? (Bericht, Kommentar, Interview, Erzählung?)} \textit{(De quel type de texte s'agit-il ?)}
    \item \textbf{Akteure :} \all{Wer sind die Hauptpersonen oder Gruppen, die genannt werden?} \textit{(Qui sont les personnes ou groupes principaux mentionnés ?)}
    \item \textbf{Vergleichender Aspekt :} \all{Welche zwei Länder oder Kulturräume werden miteinander verglichen?} \textit{(Quels deux pays ou espaces culturels sont comparés ?)}
    \item \textbf{These :} \all{Welche Hauptthese vertritt der Text?} \textit{(Quelle thèse principale le texte défend-il ?)}
\end{enumerate}
\end{exemplebox}

\begin{methode}[Répondre à une question de compréhension en phrase complète]
\textbf{Consigne :} au Bac, on ne répond jamais par un seul mot. On reprend les mots de la question pour construire une phrase allemande grammaticalement correcte (verbe conjugué en position 2, sujet explicite).

\textbf{Étapes :}
\begin{enumerate}
    \item \textbf{Analyser la question :} repérer le mot interrogatif (\all{Wer? Was? Wo? Wie lange? Warum?}) et le sujet.
    \item \textbf{Trouver l'information dans le texte :} souligner la phrase correspondante.
    \item \textbf{Transformer en réponse :} reprendre le sujet de la question, conjuguer le verbe, compléter avec l'information trouvée (souvent au passif ou avec un pronom relatif).
    \item \textbf{Vérifier :} majuscule en début de phrase, point final, ordre des mots (sujet -- verbe -- compléments en proposition principale).
\end{enumerate}

\textbf{Modèle :}
\begin{itemize}
    \item \textit{Frage :} \all{Was wird mit den Haaren der Witwe gemacht?}
    \item \textit{Textstelle :} \all{„Die Haare der Witwe werden abrasiert.“}
    \item \textit{Antwort :} \all{Die Haare der Witwe werden abrasiert.} (passif !) — « Les cheveux de la veuve sont rasés. »
\end{itemize}
\end{methode}

\begin{exoresolu}[Application : compréhension globale]
\textbf{Énoncé :} Répondez en allemand aux questions suivantes en vous basant sur le texte « Witwenrituale in Kamerun und in Deutschland ».
\begin{enumerate}
    \item \all{Worum geht es in dem Text?}
    \item \all{Welche zwei Länder werden verglichen?}
    \item \all{Wer führt die Riten in Kamerun durch?}
\end{enumerate}
\tcblower
\textbf{Solution détaillée :}
\begin{enumerate}
    \item \all{Der Text handelt von den unterschiedlichen Witwenritualen in Kamerun und in Deutschland.} \textit{(Le texte traite des différents rites de veuvage au Cameroun et en Allemagne.)} \textbf{Analyse :} sujet = \cle{der Text} ; verbe = \cle{handelt} (\all{handeln von} + datif) ; complément = \cle{den unterschiedlichen Witwenritualen} (datif pluriel avec article et adjectif déclinés).
    \item \all{Es werden Kamerun und Deutschland verglichen.} \textit{(Le Cameroun et l'Allemagne sont comparés.)} \textbf{Analyse :} passif \cle{Vorgangspassiv} (\all{werden} + participe II), car l'accent est mis sur l'action et non sur l'auteur ; \cle{es} est ici un sujet de remplissage (expletif) en tête de phrase.
    \item \all{In Kamerun werden die Riten von den älteren Frauen der Gemeinschaft durchgeführt.} \textit{(Au Cameroun, les rites sont accomplis par les femmes âgées de la communauté.)} \textbf{Analyse :} complément d'agent introduit par \cle{von} + datif pluriel (\all{von den älteren Frauen}) ; participe II du verbe à particule \all{durchführen} : \cle{durchgeführt} (particule inséparable, donc sans \all{ge-}).
\end{enumerate}
\end{exoresolu}

\courssub{Compréhension détaillée (Detailverständnis)}

Ici, l'examinateur vérifie si vous avez saisi des informations précises, chiffrées ou techniques. Les réponses doivent être \cle{précises} et \cle{justifiées}.

\begin{exemplebox}[Questions types de compréhension détaillée sur le texte]
\begin{enumerate}
    \item \all{Was geschieht mit den Haaren der Witwe in Kamerun?} \textit{(Que fait-on des cheveux de la veuve au Cameroun ?)}
    \item \all{Wie lange dauert die Trauerzeit traditionell in Kamerun?} \textit{(Combien de temps dure le deuil traditionnellement au Cameroun ?)}
    \item \all{Was muss die Witwe in Kamerun tragen?} \textit{(Que doit porter la veuve au Cameroun ?)}
    \item \all{Wie äußert sich die Trauer in Deutschland meistens?} \textit{(Comment le deuil s'exprime-t-il généralement en Allemagne ?)}
    \item \all{Welche Rolle spielen die älteren Frauen in Kamerun?} \textit{(Quel rôle jouent les femmes âgées au Cameroun ?)}
\end{enumerate}
\end{exemplebox}

\begin{propriete}[Le piège du copier-coller : il faut reformuler]
Au Bac, \textbf{recopier mot pour mot} une longue phrase du texte est pénalisé (sauf lorsque la citation est explicitement demandée, par exemple dans les items Vrai/Faux). Vous devez \cle{reformuler} : changer la structure syntaxique (actif $\leftrightarrow$ passif), remplacer le vocabulaire par des synonymes, utiliser des pronoms relatifs pour relier les idées.
\begin{center}
\begin{tabularx}{\linewidth}{|X|X|X|}
\hline
\rowcolor{popblueL} \textbf{Information du texte} & \textbf{Mauvaise réponse (copier-coller)} & \textbf{Bonne réponse (reformulée)} \\
\hline
\all{„Die Haare der Witwe werden von den alten Frauen abrasiert.“} & \all{Die Haare der Witwe werden von den alten Frauen abrasiert.} & \all{Die alten Frauen rasieren der Witwe die Haare ab.} (actif) \\
\hline
\all{„Die Trauerzeit dauert ein ganzes Jahr.“} & \all{Die Trauerzeit dauert ein ganzes Jahr.} & \all{Ein Jahr lang muss die Witwe trauern.} (verbe modal + infinitif) \\
\hline
\all{„Sie trägt schwarze Kleidung.“} & \all{Sie trägt schwarze Kleidung.} & \all{Ihre Kleidung ist schwarz.} (adjectif attribut) \\
\hline
\end{tabularx}
\end{center}
\end{propriete}

\begin{exoresolu}[Application : compréhension détaillée]
\textbf{Énoncé :} Répondez en phrases complètes en allemand.
\begin{enumerate}
    \item \all{Was passiert mit den Haaren der Witwe?}
    \item \all{Wie lange muss die Witwe in Kamerun Trauerkleidung tragen?}
    \item \all{Nennen Sie zwei Unterschiede zwischen der Trauer in Kamerun und in Deutschland.}
\end{enumerate}
\tcblower
\textbf{Solution détaillée :}
\begin{enumerate}
    \item \all{Die Haare der Witwe werden abgeschnitten.} (passif) \textit{ou} \all{Man schneidet der Witwe die Haare ab.} (tournure active avec \cle{man}). \textbf{Traduction :} « Les cheveux de la veuve sont coupés » / « On coupe les cheveux de la veuve. » \textbf{Attention :} au passif, le sujet est \all{die Haare} (nominatif pluriel) ; on ne peut pas écrire \all{den Haaren ... werden}, qui mélangerait datif et nominatif.
    \item \all{Sie muss ein Jahr lang schwarze Trauerkleidung tragen.} \textbf{Traduction :} « Elle doit porter des vêtements de deuil noirs pendant un an. » \textbf{Analyse :} \all{ein Jahr lang} = complément de durée à l'accusatif ; \all{schwarze Trauerkleidung} = accusatif féminin singulier.
    \item \all{In Kamerun sind die Riten streng vorgeschrieben und öffentlich, während die Trauer in Deutschland meist privat und individuell abläuft.} \textbf{Traduction :} « Au Cameroun, les rites sont strictement prescrits et publics, tandis qu'en Allemagne le deuil se déroule le plus souvent de façon privée et individuelle. » \textbf{Analyse :} \cle{während} introduit une subordonnée de contraste avec le verbe conjugué en fin de proposition (\all{abläuft}, verbe à particule \all{ablaufen}).
\end{enumerate}
\end{exoresolu}

\courssub{Vrai ou Faux ? Justifiez par une citation du texte}

C'est un exercice classique du Bac. La règle d'or : \textbf{aucune réponse sans citation}. La citation doit être \cle{exacte} (orthographe, majuscules, ponctuation) et \cle{courte} (le segment pertinent seulement).

\begin{methode}[Technique du Vrai/Faux justifié]
\begin{enumerate}
    \item Lire l'affirmation et repérer les \cle{mots-clés} (noms, verbes, chiffres, négations).
    \item Parcourir le texte pour localiser le passage concerné.
    \item Comparer : l'affirmation correspond-elle \textbf{exactement} au texte ? Attention aux pièges : \all{alle} contre \all{manche}, \all{immer} contre \all{oft}, \all{müssen} contre \all{können}.
    \item Rédiger : \all{Richtig: „...“} ou \all{Falsch: „...“} suivi de la citation exacte.
    \item \textbf{Ne pas traduire la citation} : la laisser en allemand.
\end{enumerate}
\end{methode}

\begin{exoresolu}[Exercice Vrai / Faux]
\textbf{Énoncé :} \all{Entscheiden Sie: Richtig oder falsch? Begründen Sie mit einem Zitat aus dem Text.}
\begin{enumerate}
    \item \all{In Kamerun dürfen die Witwen selbst entscheiden, ob sie ihre Haare schneiden.}
    \item \all{Die Trauerzeit in Kamerun dauert genau sechs Monate.}
    \item \all{In Deutschland gibt es keine staatlich vorgeschriebenen Trauerrituale.}
    \item \all{Die älteren Frauen in Kamerun unterstützen die Witwe psychologisch.}
\end{enumerate}
\tcblower
\textbf{Solution détaillée :}
\begin{enumerate}
    \item \textbf{Falsch.} \all{„Die Haare der Witwe werden abrasiert.“} \textit{(Le texte indique une obligation exprimée au passif, pas un choix personnel.)}
    \item \textbf{Falsch.} \all{„Die Trauerzeit dauert ein ganzes Jahr.“} \textit{(Le texte dit un an, pas six mois.)}
    \item \textbf{Richtig.} \all{„In Deutschland gibt es keine vorgeschriebenen Riten; die Trauer ist Privatsache.“} \textit{(Le texte souligne l'absence de rites obligatoires en Allemagne.)}
    \item \textbf{Falsch.} \all{„Die älteren Frauen überwachen die Einhaltung der Riten.“} \textit{(Le texte dit qu'elles surveillent le respect des rites, non qu'elles apportent un soutien psychologique. Attention aux interprétations excessives : ne répondez que ce que le texte dit.)}
\end{enumerate}
\end{exoresolu}

\courssub{Vocabulaire en contexte (Wortschatz im Kontext)}

Le Bac teste votre capacité à \cle{déduire le sens} d'un mot par le contexte (champ lexical, synonymes, antonymes, structure du mot) et à \cle{réemployer} le vocabulaire thématique.

\begin{definition}[Familles de mots : der Tod -- die Trauer -- das Ritual]
\begin{itemize}
    \item \cle{der Tod} (m., pluriel rare : \all{die Tode}) : la mort. \textbf{Verbe :} \all{sterben} (\all{stirbt, starb, ist gestorben}) : mourir. \textbf{Adjectif :} \all{tödlich} : mortel.
    \item \cle{die Trauer} (f., en général sans pluriel) : le deuil, la tristesse. \textbf{Verbe :} \all{trauern} (\all{trauern um} + accusatif : pleurer quelqu'un). \textbf{Adjectif :} \all{traurig} : triste. \textbf{Composés :} \all{die Trauerzeit, -en} (période de deuil), \all{die Trauerkleidung} (vêtements de deuil), \all{das Beileid} (condoléances), \all{sein Beileid aussprechen} (présenter ses condoléances).
    \item \cle{das Ritual} (n., pl. \all{die Rituale}) : le rite, le rituel. \textbf{Synonymes :} \cle{der Brauch} (m., pl. \all{die Bräuche}) : la coutume ; \cle{die Zeremonie} (f., pl. \all{-n}) : la cérémonie ; \cle{die Tradition} (f., pl. \all{-en}) : la tradition.
    \item \cle{die Witwe} (f., pl. \all{-n}) : la veuve. \textbf{Masculin :} \cle{der Witwer} (m., pl. \all{-}) : le veuf.
\end{itemize}
\end{definition}

\begin{exemplebox}[Recherche de synonymes dans le texte]
\textbf{Consigne :} trouvez dans le texte un mot ou une expression correspondant à la définition.
\begin{enumerate}
    \item Synonyme de \all{der Tod} (au sens de « décès ») $\rightarrow$ \all{der Verlust} (la perte), \all{das Ableben} (le décès, soutenu).
    \item Synonyme de \all{die Zeremonie} (action solennelle et codifiée) $\rightarrow$ \all{das Ritual}, \all{der Brauch}.
    \item Verbe signifiant « respecter, observer une règle ou une tradition » $\rightarrow$ \all{befolgen}, \all{einhalten} (particule séparable : \all{hält ein, hielt ein, hat eingehalten}).
    \item Adjectif signifiant « obligatoire, imposé par la règle » $\rightarrow$ \all{vorgeschrieben}, \all{verbindlich}.
\end{enumerate}
\end{exemplebox}

\begin{exercice}[Entraînement vocabulaire]
\textbf{Énoncé :} Complétez les phrases suivantes avec un mot du champ lexical, à la forme correcte (genre, nombre, cas, temps).
\begin{enumerate}
    \item Nach dem \_\_\_\_\_\_\_\_ (Tod) ihres Mannes muss die Frau schwarze Kleidung tragen.
    \item Die \_\_\_\_\_\_\_\_ (Ritual, pluriel) werden streng \_\_\_\_\_\_\_\_ (befolgen, passif).
    \item Eine \_\_\_\_\_\_\_\_ (Witwe) hat oft eine schwere \_\_\_\_\_\_\_\_ (Trauerzeit).
    \item In Deutschland \_\_\_\_\_\_\_\_ (trauern, présent) die Menschen oft still für sich.
\end{enumerate}
\end{exercice}

\begin{corrige}[Exercice vocabulaire]
\begin{enumerate}
    \item Nach dem \textbf{Tod} ihres Mannes muss die Frau schwarze Kleidung tragen. \textit{(\all{nach} + datif : \all{dem Tod} ; la forme \all{dem Tode} est soutenue mais correcte.)}
    \item Die \textbf{Rituale} werden streng \textbf{befolgt}. \textit{(Pluriel de \all{das Ritual} : \all{die Rituale} ; participe II de \all{befolgen} : \all{befolgt}, sans \all{ge-} car le verbe commence par le préfixe inséparable \all{be-}.)}
    \item Eine \textbf{Witwe} hat oft eine schwere \textbf{Trauerzeit}. \textit{(Nominatif féminin puis accusatif féminin : \all{eine schwere Trauerzeit}.)}
    \item In Deutschland \textbf{trauern} die Menschen oft still für sich. \textit{(Présent, 3\textsuperscript{e} personne du pluriel ; le verbe reste en position 2, le sujet suit.)}
\end{enumerate}
\end{corrige}

\courssub{Question d'interprétation (Interpretation)}

Cette question demande de dépasser le texte pour expliquer \textbf{pourquoi} ces rites existent, en mobilisant vos connaissances culturelles (\all{Landeskunde}) et votre logique. Le texte donne des indices : \all{Respekt}, \all{Tradition}, \all{Gemeinschaft}, \all{Identität}.

\begin{exemplebox}[Question type Bac : interprétation]
\all{Warum gibt es diese strengen Witwenrituale in der kamerunischen Tradition? Nennen Sie drei Gründe.}

\textbf{Modèle de réponse :}

\all{Erstens dienen die Riten dem Respekt vor dem Verstorbenen und zeigen die Verbundenheit der Witwe mit ihrem Mann. Zweitens stärken sie die Gemeinschaft, weil alle Beteiligten die gleichen Regeln befolgen. Drittens geben sie der Witwe eine klare Struktur in einer schwierigen Lebensphase.}

\textit{(Premièrement, les rites servent le respect envers le défunt et montrent le lien de la veuve avec son mari. Deuxièmement, ils renforcent la communauté, car tous les participants suivent les mêmes règles. Troisièmement, ils donnent à la veuve une structure claire dans une phase difficile de sa vie.)}
\end{exemplebox}

\begin{aretenir}[Outils d'argumentation (Redemittel)]
\begin{itemize}
    \item \all{Ein Grund ist, dass ...} (Une raison est que...)
    \item \all{Zudem / Außerdem ...} (De plus / En outre...)
    \item \all{Man kann sagen, dass ...} (On peut dire que...)
    \item \all{Die Riten dienen dazu, ...} (Les rites servent à...)
    \item \all{Sie signalisieren der Gesellschaft, dass ...} (Ils signalent à la société que...)
\end{itemize}
\end{aretenir}

\courssub{Technique du résumé (Zusammenfassung schreiben)}

Le résumé (\cle{die Zusammenfassung}) est une compétence majeure du Bac. Il ne s'agit pas de « faire court » n'importe comment, mais de \cle{restituer l'essentiel} (idées principales + articulation logique) avec ses \textbf{propres mots} (\all{mit eigenen Worten}), dans un allemand correct, en \textbf{4 à 5 phrases} (environ 80 à 100 mots).

\begin{popfigure}
\begin{tikzpicture}[
    node distance=0.9cm and 0.7cm,
    box/.style={draw=popblue, thick, fill=popblueL, rounded corners, minimum width=2.5cm, minimum height=1.1cm, align=center, font=\sffamily\small},
    arrow/.style={-Stealth, thick, popdark}
]
\node[box, align=center] (lire){\textbf{1. Lire}\\\textit{verstehen}};
\node[box, right=of lire, align=center] (rep){\textbf{2. Repérer}\\\textit{Stichpunkte}};
\node[box, right=of rep, align=center] (ref){\textbf{3. Reformuler}\\\textit{eigene Sätze}};
\node[box, right=of ref, align=center] (rel){\textbf{4. Relier}\\\textit{Konnektoren}};
\node[box, right=of rel, align=center] (ver){\textbf{5. Vérifier}\\\textit{Länge, Grammatik}};
\draw[arrow] (lire) -- (rep);
\draw[arrow] (rep) -- (ref);
\draw[arrow] (ref) -- (rel);
\draw[arrow] (rel) -- (ver);
\end{tikzpicture}
\legende{Méthode en cinq étapes pour rédiger un résumé réussi (Zusammenfassung).}
\end{popfigure}

\begin{methode}[Rédiger une Zusammenfassung : mode d'emploi]
\begin{enumerate}
    \item \textbf{Lisez deux fois :} une première lecture globale, puis une seconde lecture en soulignant les mots-clés (\all{Wer? Was? Wo? Wann? Wie? Warum?}). Ignorez les exemples, les adjectifs et les répétitions.
    \item \textbf{Notez des mots-clés en allemand :} \all{Kamerun: strenge Riten, Haare abrasiert, schwarze Kleidung, ein Jahr, ältere Frauen überwachen. -- Deutschland: privat, Blumen, Kerzen, keine Zwänge.}
    \item \textbf{Groupez et ordonnez :} introduction (sujet, pays) $\rightarrow$ Cameroun (détails) $\rightarrow$ Allemagne (contraste) $\rightarrow$ conclusion (fonction des rites).
    \item \textbf{Rédigez des phrases complètes} reliées par des \cle{connecteurs logiques} : \all{zunächst} (d'abord), \all{danach} (ensuite), \all{im Gegensatz dazu} (en revanche), \all{schließlich / zusammenfassend} (finalement / en résumé).
    \item \textbf{Contrôle final :} comptez les mots (viser 80-100). Vérifiez : le \cle{passif} pour décrire les rites, des \cle{propositions relatives} pour fluidifier, aucun « ich », aucune opinion personnelle.
\end{enumerate}
\end{methode}

\begin{attention}[Erreurs fatales au résumé, sanctionnées au Bac]
\begin{itemize}
    \item \textbf{Copier-coller :} recopier des phrases entières du texte entraîne une note très lourde de pénalité.
    \item \textbf{Opinion personnelle :} \all{Ich finde das gut.} est hors sujet dans un résumé.
    \item \textbf{Trop de détails :} inutile de citer la couleur exacte des fleurs ou le nom d'un village.
    \item \textbf{Phrases télégraphiques :} \all{Kamerun streng. Deutschland frei.} n'est pas de l'allemand correct (pas de verbe conjugué).
    \item \textbf{Oubli du contraste :} le texte \textbf{compare} deux pays ; le résumé doit rendre cette comparaison.
\end{itemize}
\end{attention}

\begin{exoresolu}[Modèle de résumé : corrigé type Bac]
\textbf{Consigne :} \all{Fassen Sie den Text in vier bis fünf Sätzen (ca. 90 Wörter) auf Deutsch zusammen.}
\tcblower
\textbf{Proposition de résumé (Musterlösung) :}

\all{Der Text vergleicht die Witwenrituale in Kamerun und Deutschland. In Kamerun sind die Riten streng vorgeschrieben: Die Witwe muss sich die Haare abrasieren lassen, ein Jahr lang schwarze Kleidung tragen, und sie wird von älteren Frauen überwacht. Diese Bräuche zeigen Respekt vor dem Toten und stärken den Gemeinschaftssinn. Im Gegensatz dazu ist die Trauer in Deutschland Privatsache ohne zwingende Rituale; man drückt sein Beileid durch Blumen oder Kerzen aus. Insgesamt spiegeln die unterschiedlichen Praktiken kulturelle Vorstellungen von Tod und sozialer Zugehörigkeit wider.}

\textbf{Analyse du modèle :}
\begin{itemize}
    \item \textbf{Phrase 1 (introduction) :} \all{Der Text vergleicht...} annonce le thème et les deux pays.
    \item \textbf{Phrase 2 (Cameroun) :} passif \all{sind vorgeschrieben} et \all{wird ... überwacht} ; modalité \all{muss ... tragen} ; tournure factitive \all{sich die Haare abrasieren lassen} (se faire raser les cheveux).
    \item \textbf{Phrase 3 (fonction des rites) :} \all{Diese Bräuche zeigen...} relie les pratiques à leur sens.
    \item \textbf{Phrase 4 (Allemagne, contraste) :} connecteur obligatoire \all{im Gegensatz dazu} ; verbe à particule \all{ausdrücken} (\all{drückt ... aus}).
    \item \textbf{Phrase 5 (conclusion) :} verbe à particule \all{widerspiegeln} (\all{spiegeln ... wider}), sujet pluriel \all{die unterschiedlichen Praktiken}.
    \item \textbf{Longueur :} environ 90 mots ; allemand standard ; présent de vérité générale, temps cohérent pour un résumé de texte.
\end{itemize}
\end{exoresolu}

\courssub{Entraînement guidé : à vous de jouer !}

\begin{exercice}[Compréhension et résumé : sujet d'entraînement Bac]
\textbf{Texte de référence :} « Witwenrituale in Kamerun und in Deutschland » (section 1).

\textbf{Partie A : questions de compréhension} (répondez en allemand, en phrases complètes).
\begin{enumerate}
    \item \all{Welche Kleidung trägt die Witwe in Kamerun während der Trauerzeit, und wie lange?}
    \item \all{Wer überwacht die Einhaltung der Riten in Kamerun?}
    \item \all{Wie drücken die Menschen in Deutschland ihr Beileid aus? Nennen Sie zwei Beispiele.}
    \item \all{Richtig oder falsch? „In Kamerun ist die Trauerzeit kürzer als in Deutschland.“ Begründen Sie mit einem Zitat.}
    \item Trouvez dans le texte un synonyme de \all{die Vorschrift} (la règle) et un synonyme de \all{privat}.
\end{enumerate}

\textbf{Partie B : production écrite : résumé.}

\all{Schreiben Sie eine Zusammenfassung des Textes auf Deutsch (vier bis fünf Sätze, ca. 80 bis 100 Wörter).} Utilisez la méthode en cinq étapes, employez le passif (\all{Vorgangspassiv}) pour décrire les rites au Cameroun et un connecteur de contraste (\all{im Gegensatz dazu}) pour l'Allemagne.
\end{exercice}

\begin{corrige}[Exercice Compréhension et résumé]
\textbf{Partie A :}
\begin{enumerate}
    \item \all{Die Witwe in Kamerun muss ein Jahr lang schwarze Trauerkleidung tragen.}
    \item \all{Die Einhaltung der Riten wird von den älteren Frauen der Gemeinschaft überwacht.}
    \item \all{In Deutschland drücken die Menschen ihr Beileid durch Blumen und Kerzen aus.}
    \item \textbf{Falsch.} \all{„Die Trauerzeit dauert ein ganzes Jahr.“} \textit{(Le texte fixe un an au Cameroun et n'indique aucune durée contraignante pour l'Allemagne ; l'affirmation est donc fausse.)}
    \item \all{die Vorschrift} $\rightarrow$ \all{die Regel} / \all{das Gebot} ; \quad \all{privat} $\rightarrow$ \all{persönlich} / \all{individuell}.
\end{enumerate}

\textbf{Partie B : proposition de résumé (une version possible parmi d'autres) :}

\all{Der Text vergleicht die Witwenrituale in Kamerun und Deutschland. In Kamerun unterliegt die Witwe strengen Vorschriften: Ihre Haare werden abrasiert, sie muss ein Jahr lang Schwarz tragen, und sie wird von älteren Frauen kontrolliert. Diese Rituale sichern den Respekt vor dem Verstorbenen und den Zusammenhalt der Gemeinschaft. Im Gegensatz dazu gestaltet sich die Trauer in Deutschland individuell und ohne Zwang; Blumen und Kerzen sind übliche Zeichen des Beileids. Die Gegenüberstellung zeigt, wie Tradition und Moderne den Umgang mit dem Tod prägen.}

\textit{(Environ 90 mots : structure claire, passifs employés, contraste marqué, vocabulaire précis.)}
\end{corrige}

\begin{savaistu}[Le veuvage au Cameroun et en Allemagne : regards croisés]
Au Cameroun, comme dans de nombreuses sociétés africaines (Bamiléké, Bassa, Bamoun, etc.), le veuvage n'est pas une affaire privée. La \cle{communauté} (\all{die Gemeinschaft}) prend en charge la veuve, mais exige aussi sa \cle{soumission} à des normes collectives. Le rasage des cheveux (\all{die Haare abrasieren}) symbolise la \cle{rupture} avec le statut d'épouse et rend le deuil \cle{visible}. La durée d'un an (\all{ein Jahr}) marque un cycle social complet.

En Allemagne, la sécularisation et l'individualisme ont transformé le deuil en \cle{Privatsache} (affaire privée) : l'État ne prescrit aucun rite de veuvage. Les funérailles sont souvent confiées à des professionnels, les \all{Bestatter} (entreprises de pompes funèbres). Le port du noir pendant une longue période subsiste dans certaines régions catholiques, mais il n'a plus rien d'obligatoire.

\textbf{Pour le Bac :} si le sujet porte sur « Tradition et modernité », ce texte est une mine d'arguments pour illustrer les oppositions \all{Tradition} contre \all{Individualismus}, \all{Gemeinschaft} contre \all{Individuum}, \all{Zwang} contre \all{Freiheit}.
\end{savaistu}

\coursec{Lexique thématique : deuil, veuvage et traditions}

Après avoir lu et compris le texte \all{Witwenrituale in Kamerun und in Deutschland}, nous disposons maintenant d'un contexte riche. Mais pour parler avec précision du veuvage, du deuil et des traditions, il faut posséder un vocabulaire solide et bien organisé. Cette section construit ce lexique pas à pas : d'abord les personnes concernées, puis le deuil lui-même, ensuite les rites et traditions, enfin les gestes concrets du rite. Nous terminerons par les familles de mots et les expressions utiles, avant de consolider le tout par des exercices.

\begin{methode}[La règle d'or du vocabulaire allemand : le triplet]
En allemand, on n'apprend jamais un nom seul. On l'apprend toujours sous forme de \cle{triplet} :
\begin{enumerate}
\item le \textbf{nom} avec sa majuscule : \all{Brauch} ;
\item son \textbf{article} (le genre) : \all{der} ;
\item son \textbf{pluriel} : \all{die Bräuche}.
\end{enumerate}
On retient donc : \all{der Brauch, die Bräuche}. Pourquoi ? Parce que l'article conditionne toute la grammaire de la phrase : la déclinaison de l'adjectif, le pronom relatif, le pronom personnel qui reprend le nom. Un élève qui dit \all{die Brauch} au lieu de \all{der Brauch} fera ensuite des fautes en cascade. Mémorisez chaque mot de cette leçon avec son article et son pluriel.
\end{methode}

\courssub{Champ 1 : les personnes}

\begin{definition}[Les personnes autour du veuvage]
\begin{itemize}
\item \all{die Witwe (-n)} : la veuve. \all{Die Witwe trägt schwarze Kleidung.} « La veuve porte des vêtements noirs. »
\item \all{der Witwer (-)} : le veuf. Attention : ce mot n'a pas de pluriel en \all{-n} ; son pluriel est identique au singulier : \all{die Witwer}.
\item \all{der Verstorbene} (n-Deklination : \all{des Verstorbenen, dem Verstorbenen, die Verstorbenen}) / \all{die Verstorbene} (féminin : \all{der Verstorbenen, die Verstorbene, die Verstorbenen}) : le défunt / la défunte. C'est un adjectif nominalisé : il prend la majuscule mais garde la déclinaison de l'adjectif.
\item \all{die Verwandten} (pluriel de \all{der/die Verwandte, -n}) : les parents, les membres de la famille. \all{Alle Verwandten kommen zur Beerdigung.} « Tous les parents viennent à l'enterrement. »
\item \all{der Erbe (-n)} : l'héritier ; \all{die Erbin (-nen)} : l'héritière.
\end{itemize}
\end{definition}

\begin{attention}[Verstorbener : un adjectif devenu nom]
\all{der Verstorbene} vient du participe passé de \all{sterben} (mourir). Comme tous les adjectifs nominalisés, il prend la \textbf{majuscule} mais garde la \textbf{déclinaison de l'adjectif} : \all{die Witwe des Verstorbenen} « la veuve du défunt », \all{wir gedenken der Verstorbenen} « nous honorons la mémoire des défunts ». Ne dites jamais \all{der Tote} dans un contexte solennel : \all{der Verstorbene} est le terme respectueux.
\end{attention}

\courssub{Champ 2 : le deuil}

\begin{definition}[Le vocabulaire du deuil]
\begin{itemize}
\item \all{die Trauer} (pas de pluriel usuel) : le deuil, le chagrin. Expression clé : \all{in Trauer sein} « être en deuil ».
\item \all{die Trauerzeit (-en)} : la période de deuil. \all{Die Trauerzeit dauert vierzig Tage.} « La période de deuil dure quarante jours. »
\item \all{die Beerdigung (-en)} : l'enterrement, les obsèques.
\item \all{das Grab (\"er)} : la tombe. Pluriel avec Umlaut : \all{die Gräber}.
\item \all{der Schmerz (-en)} : la douleur (physique ou morale). \all{Der Schmerz ist groß.} « La douleur est grande. »
\item \all{weinen (weinte, hat geweint)} : pleurer. \all{Die Kinder weinen am Grab.} « Les enfants pleurent à la tombe. »
\item \all{trauern (um + Akkusativ)} : porter le deuil de, être en deuil de. La préposition \all{um} est obligatoire pour nommer la personne pleurée.
\end{itemize}
\end{definition}

\begin{exemplebox}[Exemples traduits]
\all{Sie trauert um ihren Mann.} « Elle porte le deuil de son mari. »\par
\all{Nach der Beerdigung besucht die Familie das Grab.} « Après les obsèques, la famille visite la tombe. »\par
\all{Die ganze Familie ist in Trauer.} « Toute la famille est en deuil. »\par
\all{Die Witwe weint oft in der Nacht.} « La veuve pleure souvent la nuit. »
\end{exemplebox}

\begin{attention}[\all{begraben} $\neq$ \all{beerdigen}]
Deux verbes pour « enterrer », mais ils ne se conjuguent pas pareil :
\begin{itemize}
\item \all{begraben} : verbe \textbf{fort} et irrégulier — \all{begrub, hat begraben}. \all{Man begrub ihn auf dem Dorffriedhof.} « On l'enterra au cimetière du village. »
\item \all{beerdigen} : verbe \textbf{faible} et régulier — \all{beerdigte, hat beerdigt}. \all{Er wurde gestern beerdigt.} « Il a été inhumé hier. »
\end{itemize}
Les deux sont corrects ; \all{beerdigen} est un peu plus administratif, \all{begraben} plus imagé. Mais attention au prétérit : on ne dit jamais \all{begrabte} !
\end{attention}

\courssub{Champ 3 : les rites et les traditions}

\begin{definition}[\all{das Ritual}]
Un \all{Ritual} est une \textbf{suite codifiée de gestes et de paroles} accomplis selon un ordre précis lors d'une cérémonie religieuse ou traditionnelle. Pluriel : \all{die Rituale}. Exemple : \all{Das Ritual des Haareschneidens ist sehr alt.} « Le rituel de la coupe des cheveux est très ancien. »
\end{definition}

\begin{definition}[Traditions et croyances]
\begin{itemize}
\item \all{der Brauch (\"e)} : la coutume, l'usage. Pluriel : \all{die Bräuche}.
\item \all{die Tradition (-en)} : la tradition. Expression : \all{eine Tradition pflegen} « entretenir une tradition ».
\item \all{die Zeremonie (-n)} : la cérémonie.
\item \all{der Glaube (-n)} : la croyance, la foi. Attention : le \all{-n} du génitif et du pluriel, mais le datif singulier reste \all{dem Glauben}.
\item \all{heilig} : sacré, saint. \all{Ein heiliger Ort} « un lieu sacré ».
\item \all{befolgen} : observer, suivre (une règle, un rite). \all{Die Witwe befolgt alle Regeln.} « La veuve observe toutes les règles. »
\item \all{respektieren} : respecter. \all{Man respektiert die Traditionen der Ahnen.} « On respecte les traditions des ancêtres. »
\end{itemize}
\end{definition}

\begin{exemplebox}[Exemples traduits]
\all{Der Brauch wird seit Generationen gepflegt.} « La coutume est entretenue depuis des générations. »\par
\all{Mit diesem Ritual zeigt die Familie den Respekt vor dem Verstorbenen.} « Par ce rituel, la famille montre le respect envers le défunt. »\par
\all{Der Glaube an die Ahnen spielt eine wichtige Rolle.} « La croyance aux ancêtres joue un rôle important. »
\end{exemplebox}

\begin{savaistu}[Le deuil dans les pays germanophones]
En \textbf{Autriche} catholique, on célèbre la \all{Todestagsmesse}, une messe dite chaque année à l'anniversaire du décès d'un proche : la famille se réunit à l'église puis partage un repas. En \textbf{Suisse} alémanique, les avis de décès (\all{Todesanzeigen}) paraissent encore aujourd'hui dans les journaux locaux, souvent avec une croix, une phrase de remerciement et l'annonce des obsèques. En Allemagne, la période de deuil traditionnelle des veuves s'appelait autrefois \all{das Trauerjahr} : pendant un an, la veuve portait du noir. Comme au Cameroun, le deuil y est donc aussi une affaire de communauté et de gestes visibles.
\end{savaistu}

\courssub{Champ 4 : les gestes du rite}

Les rites de veuvage se décrivent par des \textbf{actions}. Voici les verbes et expressions essentiels :

\begin{center}
\begin{tabularx}{\linewidth}{|X|X|X|}\hline
\rowcolor{popblueL} Deutsch & Français & Exemple \\ \hline
\all{sich die Haare schneiden} & se couper les cheveux & \all{Die Witwe schneidet sich die Haare.} \\ \hline
\all{schwarze Kleidung tragen} & porter des vêtements noirs & \all{Sie trägt ein Jahr lang schwarze Kleidung.} \\ \hline
\all{fasten (fastete, hat gefastet)} & jeûner & \all{Die Familie fastet drei Tage.} \\ \hline
\all{beten (betete, hat gebetet)} & prier & \all{Man betet für den Verstorbenen.} \\ \hline
\all{erben (erbte, hat geerbt)} & hériter & \all{Der Sohn erbt das Haus.} \\ \hline
\all{verbieten (verbot, hat verboten)} & interdire & \all{Der Brauch verbietet der Witwe, das Haus zu verlassen.} \\ \hline
\end{tabularx}
\end{center}

\begin{attention}[Verbes à régime particulier]
\begin{itemize}
\item \all{beten für + Akkusativ} : prier \textbf{pour} quelqu'un. \all{Wir beten für die Verstorbenen.} « Nous prions pour les défunts. »
\item \all{trauern um + Akkusativ} : porter le deuil \textbf{de}. Ne traduisez jamais mot à mot par \all{für} !
\item \all{verbieten} est un verbe \textbf{fort} : \all{verbot, hat verboten}. \all{Man verbot ihr, zu arbeiten.} « On lui interdit de travailler. » Il se construit avec l'infinitif en \all{zu} : \all{jemandem verbieten, etwas zu tun}.
\item \all{sich die Haare schneiden} : en allemand, on utilise l'article défini et non le possessif pour les parties du corps. Pas \all{ihre Haare}, mais \all{sich die Haare}.
\end{itemize}
\end{attention}

\courssub{Familles de mots : un mot, trois formes}

L'allemand construit des familles de mots très régulières autour d'une même racine. Les connaître multiplie votre vocabulaire par trois sans effort supplémentaire. La dérivation suit souvent le sens : \textbf{Verbe $\rightarrow$ Nom $\rightarrow$ Adjectif}.

\begin{popfigure}
\begin{center}
\begin{tikzpicture}[node distance=1.5cm and 2.5cm]
\node[draw=popblue, fill=popblueL, rounded corners, align=center, font=\bfseries] (v1) {Verb};
\node[draw=poporange, fill=poporangeL, rounded corners, align=center, font=\bfseries, right=of v1] (n1) {Substantiv};
\node[draw=popgreen, fill=popgreenL, rounded corners, align=center, font=\bfseries, right=of n1] (a1) {Adjektiv};

\node[align=center] (v1a) at (v1.south) [anchor=north, yshift=-0.5cm] {\all{trauern}};
\node[align=center] (n1a) at (n1.south) [anchor=north, yshift=-0.5cm] {\all{die Trauer}};
\node[align=center] (a1a) at (a1.south) [anchor=north, yshift=-0.5cm] {\all{traurig}};

\node[align=center] (v2a) at (v1a.south) [anchor=north, yshift=-0.5cm] {\all{sterben}};
\node[align=center] (n2a) at (n1a.south) [anchor=north, yshift=-0.5cm] {\all{der Tod}};
\node[align=center] (a2a) at (a1a.south) [anchor=north, yshift=-0.5cm] {\all{tot}};
\node[align=center] (n2b) at (n2a.south) [anchor=north, yshift=-0.3cm] {\all{der Verstorbene}};

\node[align=center] (v3a) at (v2a.south) [anchor=north, yshift=-0.5cm] {\all{überliefern}};
\node[align=center] (n3a) at (n2a.south) [anchor=north, yshift=-0.5cm] {\all{die Tradition}};
\node[align=center] (a3a) at (a2a.south) [anchor=north, yshift=-0.5cm] {\all{traditionell}};

\draw[-{Stealth}, popmuted, thick] (v1a) -- (n1a);
\draw[-{Stealth}, popmuted, thick] (n1a) -- (a1a);
\draw[-{Stealth}, popmuted, thick] (v2a) -- (n2a);
\draw[-{Stealth}, popmuted, thick] (n2a) -- (a2a);
\draw[-{Stealth}, popmuted, thick, dashed] (n2a) -- (n2b);
\draw[-{Stealth}, popmuted, thick] (v3a) -- (n3a);
\draw[-{Stealth}, popmuted, thick] (n3a) -- (a3a);
\end{tikzpicture}
\end{center}
\legende{Trois familles de mots : du verbe au nom, du nom à l'adjectif (flèche pointillée = dérivation secondaire).}
\end{popfigure}

\begin{exemplebox}[Les familles en phrases]
\begin{itemize}
\item \all{trauern $\rightarrow$ die Trauer $\rightarrow$ traurig} : \all{Die Witwe trauert. Ihre Trauer ist groß. Sie ist sehr traurig.} « La veuve porte le deuil. Son chagrin est grand. Elle est très triste. »
\item \all{sterben $\rightarrow$ der Tod $\rightarrow$ tot / der Verstorbene} : \all{Der Mann starb. Sein Tod erschüttert die Familie. Er ist tot. Man ehrt den Verstorbenen.} « L'homme mourut. Sa mort bouleverse la famille. Il est mort. On honore le défunt. »
\item \all{überliefern $\rightarrow$ die Tradition $\rightarrow$ traditionell} : \all{Man überlieferte den Brauch. Die Tradition ist alt. Es ist ein traditionelles Ritual.} « On a transmis la coutume. La tradition est ancienne. C'est un rituel traditionnel. »
\end{itemize}
\end{exemplebox}

\begin{aretenir}[Expressions utiles à réemployer]
\begin{itemize}
\item \all{den Respekt vor dem Verstorbenen zeigen} : « montrer le respect envers le défunt ». Attention au datif : \all{vor dem Verstorbenen}.
\item \all{eine Tradition pflegen} : « entretenir, perpétuer une tradition ».
\item \all{in Trauer sein} : « être en deuil ».
\item \all{um jemanden trauern} : « porter le deuil de quelqu'un ».
\item \all{jemandem etwas verbieten} : « interdire quelque chose à quelqu'un » (Datif + Accusatif).
\end{itemize}
\end{aretenir}

\courssub{Texte d'application : le lexique en contexte}

\begin{textebox}[Ein altes Ritual in einem Dorf bei Bafoussam]
In einem Dorf bei Bafoussam starb vor einem Monat der Lehrer Mbappe. Seine Frau, die Witwe Nathalie, trauert um ihren Mann. Die ganze Familie ist in Trauer, und der Schmerz ist groß.

Nach der Beerdigung besuchten alle Verwandten das Grab. Dann begann die Trauerzeit. Nach einem alten Brauch des Dorfes schneidet sich die Witwe die Haare und trägt vierzig Tage lang schwarze Kleidung. Sie fastet am Morgen und betet jeden Abend für den Verstorbenen. Das Ritual verbietet ihr, das Haus allein zu verlassen. Eine Schwester ihres Mannes bleibt bei ihr und hilft ihr.

„Diese Tradition pflegen wir seit Generationen“, erklärt der Älteste der Familie. „Mit diesen Ritualen zeigen wir den Respekt vor dem Verstorbenen. Der Glaube an die Ahnen ist für uns heilig.“

Am Ende der Trauerzeit gibt es eine große Zeremonie. Die Familie kocht, die Nachbarn kommen, und man dankt Gott. Der älteste Sohn erbt das Haus des Vaters, denn er ist der erste Erbe. Aber die Witwe darf dort weiter wohnen — das respektiert die ganze Familie.

Nathalie ist noch traurig, aber sie sagt: „Die Rituale helfen mir. Ich bin nicht allein. Die Familie und das Dorf tragen meinen Schmerz mit mir.“
\end{textebox}

\textbf{Glossaire.} \all{der Älteste (-n)} (n-Deklination) : l'aîné, le doyen ; \all{die Ahnen} (pluriel) : les ancêtres ; \all{die Nachbarn} (pluriel de \all{der Nachbar, -n}) : les voisins ; \all{weiter wohnen} : continuer à habiter ; \all{tragen (trug, hat getragen)} : ici « porter, supporter » au figuré.

\textbf{Questions de compréhension.}
\begin{enumerate}
\item \all{Warum schneidet sich Nathalie die Haare?}
\item \all{Was verbietet das Ritual der Witwe?}
\item \all{Wer erbt das Haus des Verstorbenen?}
\item \all{Wie lange dauert die Trauerzeit?}
\end{enumerate}

\courssub{Expression orale guidée}

\begin{experience}[Jeu de rôle : « Au chevet de la veuve »]
\textbf{Situation :} Nous sommes dans la cour de la famille Mbappe, à Bafoussam. La période de deuil de 40 jours vient de s'achever.
\textbf{Rôles (groupes de 3-4) :}
\begin{itemize}
\item \textbf{La veuve Nathalie} : elle explique ce qu'elle a vécu (cheveux coupés, jeûne, prières, interdiction de sortir), ses sentiments (douleur, mais aussi soutien).
\item \textbf{L'aîné de la famille (der Älteste)} : il justifie le rite (respect des ancêtres, tradition, cohésion du clan).
\item \textbf{Un(e) voisin(e) / ami(e)} : il/elle pose des questions, exprime sa compassion, compare avec d'autres coutumes (ou avec l'Allemagne).
\item \textbf{(Optionnel) Un journaliste de la CRTV} : il mène l'interview pour un reportage sur « Tradition et modernité ».
\end{itemize}
\textbf{Contraintes linguistiques :} Utiliser au moins 5 mots du lexique (\all{trauern, der Brauch, das Ritual, verbieten, der Verstorbene, die Tradition pflegen, in Trauer sein, der Respekt vor...}). Utiliser le prétérit pour raconter le passé (\all{starb, begann, verbot, trug, fastete}) et le présent pour les sentiments actuels.
\textbf{Déroulement :} 5 min de préparation (notes en allemand), 3-4 min de jeu par groupe. Quelques groupes jouent devant la classe.
\end{experience}

\courssub{Exercice résolu et entraînement}

\begin{exoresolu}[Complétez avec le mot qui convient]
Complétez avec : \all{die Witwe, das Grab, trauert, der Brauch, erbt, verbietet}.
\begin{enumerate}
\item Nach dem Tod ihres Mannes ist Frau Kamga eine \dots
\item Sie \dots\ um ihren Mann und weint oft.
\item Die Familie besucht jeden Sonntag das \dots\ des Verstorbenen.
\item Es ist ein alter \dots, dass die Witwe schwarze Kleidung trägt.
\item Das Ritual \dots\ ihr, in den ersten Wochen zu arbeiten.
\item Der älteste Sohn \dots\ das Feld seines Vaters.
\end{enumerate}
\tcblower
\textbf{Solution détaillée.}
\begin{enumerate}
\item \all{eine Witwe} : après la mort de son mari, une femme devient veuve. Article indéfini féminin à l'accusatif : \all{eine}.
\item \all{trauert} : verbe \all{trauern um + Akkusativ}, conjugué à la 3\up{e} personne du singulier : \all{sie trauert}.
\item \all{das Grab} : la tombe ; neutre, accusatif identique au nominatif.
\item \all{ein alter Brauch} : la coutume ; masculin, nominatif après \all{es ist}, d'où l'adjectif en \all{-er}.
\item \all{verbietet} : verbe fort \all{verbieten}, au présent 3\up{e} personne : \all{das Ritual verbietet}. Construction : \all{jemandem verbieten, etwas zu tun}.
\item \all{erbt} : verbe \all{erben}, 3\up{e} personne du singulier : \all{der Sohn erbt}.
\end{enumerate}
\end{exoresolu}

\begin{exercice}[À vous]
\begin{enumerate}
\item Traduisez en allemand : « La veuve porte le deuil de son mari et respecte toutes les traditions. »
\item Formez la famille de mots de \all{sterben} (nom, adjectif, nom de la personne défunte) et écrivez une phrase avec chaque forme.
\item Complétez avec la bonne préposition : \all{Wir beten \dots\ den Verstorbenen. Sie trauert \dots\ ihren Vater.}
\item Mettez au pluriel : \all{der Brauch, das Grab, die Zeremonie, das Ritual, die Witwe}.
\end{enumerate}
\end{exercice}

\begin{corrige}[Exercice « À vous »]
\begin{enumerate}
\item \all{Die Witwe trauert um ihren Mann und respektiert alle Traditionen.} (Ne pas oublier \all{um} après \all{trauern} et la majuscule aux noms.)
\item \all{der Tod, tot, der Verstorbene}. Exemples : \all{Sein Tod war plötzlich. Der Mann ist tot. Wir ehren den Verstorbenen.}
\item \all{für} (prier pour) ; \all{um} (porter le deuil de).
\item \all{die Bräuche, die Gräber, die Zeremonien, die Rituale, die Witwen}.
\end{enumerate}
\end{corrige}

\begin{aretenir}[L'essentiel du lexique]
Retenez les triplets fondamentaux : \all{die Witwe (-n)}, \all{der Witwer (-)}, \all{der Verstorbene} (n-Dekl.), \all{die Trauer}, \all{die Trauerzeit (-en)}, \all{die Beerdigung (-en)}, \all{das Grab (\"er)}, \all{der Brauch (\"e)}, \all{die Tradition (-en)}, \all{das Ritual (-e)}, \all{die Zeremonie (-n)}. Maîtrisez les constructions à préposition : \all{trauern um + A}, \all{beten für + A}, \all{der Respekt vor + D}. Et rappelez-vous : un mot allemand s'apprend toujours avec son article et son pluriel — c'est la clé de toute la grammaire.
\end{aretenir}

\coursec{Grammaire 1 : le passif (Vorgangspassiv) pour décrire un rite}

Dans la section précédente, nous avons constitué le lexique du deuil et du veuvage : \all{die Witwe, die Witwen}, \all{der Witwer, die Witwer}, \all{der Brauch, die Bräuche}, \all{die Trauer}, \all{das Ritual, die Rituale}. Mais un lexique ne suffit pas : pour \emph{décrire} un rite, pour raconter objectivement ce qui se fait, étape par étape, lors d'une cérémonie traditionnelle, l'allemand dispose d'un outil grammatical privilégié : le \cle{passif}. C'est lui que nous allons maintenant observer, comprendre et manipuler.

\courssub{Observation : que remarque-t-on dans le texte ?}

Reprenons deux phrases tirées de notre texte support \all{„Witwenrituale in Kamerun und in Deutschland“} :

\all{Die Haare werden geschnitten.} « Les cheveux sont coupés. »

\all{Die Riten wurden befolgt.} « Les rites ont été suivis. »

Observons attentivement. Dans les deux phrases, on trouve :

\begin{itemize}
\item une forme du verbe \all{werden} conjuguée (\all{werden} au présent, \all{wurden} au prétérit) ;
\item un \cle{participe passé} (\all{geschnitten}, \all{befolgt}) placé \textbf{en toute fin de phrase} ;
\item un sujet qui \emph{subit} l'action : ce ne sont pas les cheveux qui coupent, ce ne sont pas les rites qui suivent ! Les cheveux \emph{sont coupés}, les rites \emph{sont suivis}.
\end{itemize}

Qui coupe les cheveux ? Qui suit les rites ? La phrase ne le dit pas — et c'est précisément l'intérêt du passif : l'accent est mis sur \textbf{l'action} et sur \textbf{celui qui la subit}, pas sur celui qui la fait.

\begin{definition}[Le passif (das Passiv)]
Le \cle{passif} est une forme verbale dans laquelle le sujet grammatical \emph{subit} l'action au lieu de la faire. En allemand, on distingue le \cle{Vorgangspassiv} (passif d'action, formé avec \all{werden} + participe passé), qui décrit une action ou un processus, et le \cle{Zustandspassiv} (passif d'état, formé avec \all{sein} + participe passé : \all{Die Tür ist geschlossen.} « La porte est fermée. »). Au programme de la Terminale : le \textbf{Vorgangspassiv}, simplement appelé « passif ».
\end{definition}

\begin{propriete}[Formation du Vorgangspassiv]
\textbf{Forme :} sujet + \all{werden} conjugué (à la 2\textsuperscript{e} position) + participe passé \textbf{en dernière position}.

\begin{center}
\begin{tabular}{|l|l|l|}
\hline
\rowcolor{popblueL} Sujet (celui qui subit) & \all{werden} conjugué & Participe passé (fin) \\
\hline
\all{Die Haare} & \all{werden} & \all{geschnitten.} \\
\hline
\all{Das Ritual} & \all{wird} & \all{wiederholt.} \\
\hline
\all{Die Riten} & \all{wurden} & \all{befolgt.} \\
\hline
\end{tabular}
\end{center}

Règles essentielles :
\begin{itemize}
\item \all{werden} porte \textbf{seul} la conjugaison (personne, nombre, temps) ;
\item le participe passé est \textbf{invariable} : il ne s'accorde jamais, contrairement au français (« les cheveux sont coupé\emph{s} », mais \all{die Haare werden geschnitten}) ;
\item le participe passé occupe la \textbf{dernière place} de la phrase : c'est la fameuse « parenthèse » verbale allemande (\all{Satzklammer}).
\end{itemize}
\end{propriete}

\begin{attention}[Le participe passé ne s'accorde pas]
Le francophone a le réflexe d'accorder le participe : « la veuve est soutenu\emph{e} », « les rites sont suivi\emph{s} ». En allemand, \textbf{jamais d'accord} : \all{Die Witwe wird unterstützt.} / \all{Die Riten werden befolgt.} Une seule forme, en toutes circonstances.
\end{attention}

\courssub{Les temps du passif}

Le passif peut se conjuguer à tous les temps : c'est \all{werden} qui change de forme, le participe passé reste fixe en fin de phrase. Voici le tableau complet avec le verbe \all{schneiden} (couper), verbe fort : \all{schneiden – schnitt – geschnitten}.

\begin{center}
\begin{tabularx}{\linewidth}{|l|X|X|}
\hline
\rowcolor{popblueL} Temps & Forme (3\textsuperscript{e} pers. sing.) & Exemple traduit \\
\hline
Präsens & \all{wird geschnitten} & \all{Die Haare werden geschnitten.} « Les cheveux sont coupés. » \\
\hline
Präteritum & \all{wurde geschnitten} & \all{Die Haare wurden geschnitten.} « Les cheveux furent coupés. » \\
\hline
Perfekt & \all{ist geschnitten worden} & \all{Die Haare sind geschnitten worden.} « Les cheveux ont été coupés. » \\
\hline
Futur & \all{wird geschnitten werden} & \all{Die Haare werden geschnitten werden.} « Les cheveux seront coupés. » \\
\hline
\end{tabularx}
\end{center}

\begin{propriete}[Conjugaison de \all{werden} au passif]
\begin{center}
\begin{tabular}{|l|l|l|}
\hline
\rowcolor{popblueL} Personne & Präsens & Präteritum \\
\hline
ich & werde & wurde \\
\hline
du & wirst & wurdest \\
\hline
er/sie/es & wird & wurde \\
\hline
wir & werden & wurden \\
\hline
ihr & werdet & wurdet \\
\hline
sie/Sie & werden & wurden \\
\hline
\end{tabular}
\end{center}
Attention : au passif, \all{werden} est un \textbf{auxiliaire} et ne signifie pas « devenir ». \all{Die Haare werden geschnitten} ne veut pas dire « les cheveux deviennent coupés » mais « les cheveux sont coupés ».
\end{propriete}

\begin{attention}[Perfekt passif : \all{worden}, pas \all{geworden} !]
C'est \textbf{l'erreur classique} des francophones. Au Perfekt passif, le participe de \all{werden} perd son \all{ge-} : on dit \all{ist gemacht \textbf{worden}} et jamais « ist gemacht geworden ». Comparez :
\begin{itemize}
\item \all{Sie ist Lehrerin geworden.} « Elle est devenue professeure. » (\all{werden} verbe plein → \all{geworden})
\item \all{Das Grab ist geschmückt worden.} « La tombe a été décorée. » (passif → \all{worden})
\end{itemize}
Moyen mnémotechnique : au passif, \all{werden} a déjà assez de travail comme auxiliaire — il perd son \all{ge-} !
\end{attention}

\begin{popfigure}
\begin{tikzpicture}[font=\small]
\draw[thick, popline] (0,0) -- (10.5,0);
\foreach \x/\temps/\forme in {0.8/{Präsens}/{wird gemacht}, 5.2/{Präteritum}/{wurde gemacht}, 9.6/{Perfekt}/{ist gemacht worden}}{
  \fill[poporange] (\x,0) circle (0.12);
  \node[above=0.25cm, align=center] at (\x,0){ \textbf{\temps}\\ \all{\forme} };
}
\draw[-{Stealth}, thick, popblue] (0,-0.6) -- (10.5,-0.6) node[midway, below, text=popmuted] {passé $\longrightarrow$};
\end{tikzpicture}
\legende{Frise des temps du passif : seul l'auxiliaire \all{werden} change ; le participe reste en fin de phrase.}
\end{popfigure}

\courssub{Le complément d'agent et l'instrument}

Si l'on veut tout de même nommer celui qui fait l'action, l'allemand utilise \cle{von + Datif} — exactement comme le français « par » :

\all{Die Witwe wurde von ihrer Familie unterstützt.} « La veuve a été soutenue par sa famille. »

Pour l'objet avec lequel l'action est faite (l'instrument), on utilise \cle{mit + Datif} :

\all{Die Haare werden mit einer Schere geschnitten.} « Les cheveux sont coupés avec des ciseaux. »

\begin{center}
\begin{tabularx}{\linewidth}{|X|X|X|}
\hline
\rowcolor{popblueL} Fonction & Préposition & Exemple \\
\hline
Agent (personne) & \all{von} + Datif & \all{von den Frauen} « par les femmes » \\
\hline
Instrument (chose) & \all{mit} + Datif & \all{mit einer Schere} « avec des ciseaux » \\
\hline
Cause, moyen & \all{durch} + Akkusativ & \all{durch die Tradition} « par la tradition » \\
\hline
\end{tabularx}
\end{center}

\begin{attention}
Ne confondez pas \all{von} (agent) et \all{mit} (instrument) : \all{Die Witwe wird von den Frauen mit Wasser gewaschen.} « La veuve est lavée par les femmes avec de l'eau. » La personne → \all{von} ; la chose → \all{mit}.
\end{attention}

\courssub{Transformation actif → passif}

Le passage de l'actif au passif obéit à une mécanique précise. Observons :

\all{Die Frauen schneiden die Haare.} (actif) → \all{Die Haare werden (von den Frauen) geschnitten.} (passif)

\begin{enumerate}
\item L'\textbf{objet accusatif} de l'actif (\all{die Haare}) devient le \textbf{sujet nominatif} du passif ;
\item le verbe actif (\all{schneiden}) devient \all{werden} + participe passé (\all{werden geschnitten}) ;
\item l'ancien sujet (\all{die Frauen}) devient complément d'agent \all{von + Datif} (\all{von den Frauen}) — ou disparaît s'il est sans importance ;
\item le verbe \all{werden} s'accorde avec le \textbf{nouveau} sujet : \all{die Haare} (pluriel) → \all{werden}.
\end{enumerate}

\begin{popfigure}
\begin{tikzpicture}[font=\small, node distance=0.35cm]
\node[draw=popblue, fill=popblueL, rounded corners, align=center, text width=4.6cm] (aktiv) {\textbf{Actif}\\ \all{Die Frauen} (Subjekt)\\ \all{schneiden} (Verb)\\ \all{die Haare} (Akkusativ)};
\node[draw=poporange, fill=poporangeL, rounded corners, align=center, text width=5.4cm, right=2.2cm of aktiv] (passiv) {\textbf{Passiv}\\ \all{Die Haare} (neues Subjekt)\\ \all{werden geschnitten} (werden + Partizip II)\\ \all{(von den Frauen)} (von + Dativ)};
\draw[-{Stealth}, very thick, poporange] (aktiv) -- (passiv) node[midway, above, align=center, text=popdark] {Objekt $\rightarrow$ Subjekt};
\draw[-{Stealth}, thick, poppurple, bend right=35] (aktiv.north) to node[above, align=center, text=popdark] {Subjekt $\rightarrow$ \all{von} + Dativ} (passiv.north);
\end{tikzpicture}
\legende{Schéma de la transformation actif → passif : l'objet direct devient sujet, l'ancien sujet devient complément d'agent.}
\end{popfigure}

\begin{propriete}[Condition : un objet direct]
Seul un verbe construit avec un \textbf{objet direct (accusatif)} peut former un passif « normal » (personnel). \all{Die Frauen schneiden die Haare} → passif possible. Mais \all{Die Witwe weint} (« la veuve pleure », sans objet) → pas de passif personnel possible. C'est une différence majeure avec le français « on ».
\end{propriete}

\begin{attention}[Le français « on » ≠ le passif allemand]
Le français dit volontiers « \emph{on} coupe les cheveux ». L'allemand a deux solutions :
\begin{itemize}
\item l'actif avec \all{man} : \all{Man schneidet die Haare.} (tournure la plus fréquente à l'oral) ;
\item le passif, qui exige un objet direct : \all{Die Haare werden geschnitten.}
\end{itemize}
Mais on ne peut \textbf{jamais} traduire « on pleure » par un passif personnel : \all{Man weint.} uniquement. Le passif allemand est donc \emph{moins} libre que le « on » français.
\end{attention}

\courssub{Le passif avec un verbe modal}

Très utile pour décrire les obligations d'un rite : modal conjugué + participe passé + \all{werden} à l'infinitif, \textbf{tout à la fin}.

\all{Die Kleider müssen getragen werden.} « Les vêtements doivent être portés. »

\all{Die Trauerzeit muss respektiert werden.} « La période de deuil doit être respectée. »

\begin{center}
\begin{tabularx}{\linewidth}{|X|X|}
\hline
\rowcolor{popblueL} Deutsch & Français \\
\hline
\all{Das Ritual muss befolgt werden.} & Le rituel doit être suivi. \\
\hline
\all{Die Regeln können geändert werden.} & Les règles peuvent être modifiées. \\
\hline
\all{Der Brauch soll bewahrt werden.} & La coutume doit être préservée. \\
\hline
\all{Die Witwe darf nicht allein gelassen werden.} & La veuve ne doit pas être laissée seule. \\
\hline
\end{tabularx}
\end{center}

Structure à mémoriser : sujet + \textbf{modal conjugué} + … + \textbf{participe passé} + \all{werden} (infinitif, dernière place absolue).

\courssub{Emploi privilégié : décrire un rite objectivement}

Le passif est la forme idéale pour décrire un procédé, une cérémonie, une recette ou un rite, car il met l'action au premier plan et donne un ton objectif, presque documentaire. C'est exactement le ton de notre texte support.

\begin{textebox}[„Die Schritte eines traditionellen Ritus“ — texte d'application]
\all{Zuerst wird die Witwe von den ältesten Frauen der Familie empfangen. Dann werden ihre Haare geschnitten, und sie wird mit schwarzen Kleidern bedeckt. Danach wird ein traditionelles Essen vorbereitet: Das Essen wird von den Nachbarn gebracht und mit der ganzen Familie geteilt. Während der Trauerzeit werden Gebete gesprochen, und jeden Abend wird eine Kerze angezündet. Das Grab des Verstorbenen wird mit Blumen geschmückt. Früher wurden die Riten sehr streng befolgt; heute werden sie in den Städten oft vereinfacht. Trotzdem muss die Trauerzeit respektiert werden, denn die Tradition soll an die jungen Generationen weitergegeben werden. Am Ende der Zeremonie wird die Witwe offiziell wieder in die Gemeinschaft aufgenommen.}

\textbf{Glossaire :} \all{empfangen} : accueillir ; \all{bedecken} : couvrir ; \all{die Kerze, die Kerzen} : la bougie ; \all{anzünden} : allumer ; \all{das Grab, die Gräber} : la tombe ; \all{der Verstorbene, die Verstorbenen} : le défunt ; \all{streng} : strictement ; \all{vereinfachen} : simplifier ; \all{weitergeben} : transmettre ; \all{aufnehmen} : accueillir, réintégrer ; \all{die Gemeinschaft, die Gemeinschaften} : la communauté.

\textbf{Questions :}
1) \all{Wer empfängt die Witwe?}
2) \all{Was wird mit den Haaren der Witwe gemacht?}
3) \all{Wie werden die Riten heute in den Städten verändert?}
\end{textebox}

Comptez les passifs dans ce texte : il y en a plus de dix ! C'est bien la preuve que le passif est \emph{la} forme de la description rituelle.

\begin{methode}[Décrire un rite étape par étape au passif]
\begin{enumerate}
\item Repérer les étapes du rite (accueil, coupe des cheveux, repas, prières, fin du deuil).
\item Pour chaque étape, identifier l'objet qui subit l'action : \all{die Haare}, \all{das Essen}, \all{das Grab}…
\item Construire le passif : objet-sujet + \all{werden} conjugué + participe passé en fin.
\item Enchaîner avec des connecteurs chronologiques : \all{zuerst wird …, dann wird …, danach wird …, schließlich wird …}
\item Ajouter si besoin un agent (\all{von + Dativ}) ou un modal (\all{muss … werden}).
\end{enumerate}
\end{methode}

\begin{exemplebox}[Exemples traduits]
\all{Das Ritual wird jedes Jahr wiederholt.} « Le rituel est répété chaque année. »

\all{Die Witwe wurde von ihrer Familie unterstützt.} « La veuve a été soutenue par sa famille. »

\all{Das Grab ist geschmückt worden.} « La tombe a été décorée. »

\all{In Kamerun werden viele Traditionen noch gepflegt.} « Au Cameroun, beaucoup de traditions sont encore entretenues. »
\end{exemplebox}

\begin{exoresolu}[Transformation actif → passif]
Mettez au passif (même temps) : \all{Die Frauen bereiteten das Essen vor.} Puis : \all{Man hat die Witwe unterstützt.}

\tcblower
\textbf{Solution détaillée.}

Phrase 1 : le verbe est au Präteritum (\all{bereiteten … vor}, verbe à particule séparable \all{vorbereiten}). L'objet accusatif \all{das Essen} devient sujet ; \all{werden} se met au Prétérit à la 3\textsuperscript{e} personne du singulier : \all{wurde} ; le participe \all{vorbereitet} va en fin de phrase.

\all{Das Essen wurde (von den Frauen) vorbereitet.} « Le repas fut préparé (par les femmes). »

Phrase 2 : le verbe est au Perfekt (\all{hat … unterstützt}). Le passif au Perfekt exige \all{sein} + participe + \all{worden} (sans \all{ge-} !). L'objet \all{die Witwe} devient sujet.

\all{Die Witwe ist unterstützt worden.} « La veuve a été soutenue. »

Piège évité : on n'écrit pas « ist unterstützt geworden ».
\end{exoresolu}

\begin{exercice}[À vous]
1. Mettez au passif : a) \all{Die Familie schmückt das Grab.} b) \all{Die Nachbarn brachten das Essen.} c) \all{Man muss die Tradition bewahren.}

2. Traduisez : « Les rites doivent être respectés. » / « Les cheveux ont été coupés. »
\end{exercice}

\begin{corrige}[Exercice « À vous »]
1. a) \all{Das Grab wird (von der Familie) geschmückt.} b) \all{Das Essen wurde (von den Nachbarn) gebracht.} c) \all{Die Tradition muss bewahrt werden.}

2. \all{Die Riten müssen respektiert werden.} / \all{Die Haare sind geschnitten worden.}
\end{corrige}

\begin{aretenir}[L'essentiel du Vorgangspassiv]
\begin{itemize}
\item Forme : \all{werden} conjugué + participe passé invariable en dernière position.
\item Temps : \all{wird gemacht} / \all{wurde gemacht} / \all{ist gemacht worden} / \all{wird gemacht werden}.
\item Agent : \all{von + Datif} ; instrument : \all{mit + Datif}.
\item Avec modal : modal + participe + \all{werden} à l'infinitif final.
\item Emploi : description objective des étapes d'un rite ; le français « on » se rend souvent par \all{man} + actif.
\end{itemize}
\end{aretenir}

\coursec{Grammaire 2 : propositions relatives, prétérit des verbes forts et participes}

Dans la section précédente, nous avons appris à décrire un rite au moyen du passif : \all{Die Riten werden von den Frauen gepflegt.} Mais comment insérer une information supplémentaire sur ces rites sans répéter le nom ? Comment raconter ensuite ce qui s'est passé autrefois ? C'est exactement le rôle des \cle{propositions relatives} et du \cle{prétérit des verbes forts}, que nous étudions maintenant. Ces trois outils — passif, relatives, prétérit — forment ensemble la grammaire du récit traditionnel.

\courssub{Observation : le pronom relatif dans le texte}

Reprenons une phrase de notre texte support :

\all{Die Riten, die bei uns beobachtet werden, sind alt.} « Les rites qui sont observés chez nous sont anciens. »

Et trois autres exemples tirés du thème du veuvage :

\all{Die Witwe, die im Dorf lebt, trägt Schwarz.} « La veuve qui vit au village porte du noir. »

\all{Der Brauch, den die Frauen pflegen, ist sehr alt.} « La coutume que les femmes entretiennent est très ancienne. »

\all{Die Frauen, von denen die Riten befolgt werden, sind alt.} « Les femmes par qui les rites sont observés sont âgées. »

Observons. Le mot \all{die} dans la première phrase reprend \all{die Riten} : il a le même \cle{genre} et le même \cle{nombre} que le nom auquel il se réfère. En revanche, dans la deuxième phrase, on trouve \all{den} et non \all{der} : pourquoi ? Parce que dans la relative, ce pronom est \emph{complément d'objet direct} de \all{pflegen} (les femmes entretiennent \emph{quoi ?} — la coutume). Le \cle{cas} du relatif dépend donc de sa fonction \emph{à l'intérieur de la subordonnée}, et non de la fonction du nom dans la phrase principale.

\begin{propriete}[La double règle du pronom relatif]
Le pronom relatif obéit à deux règles distinctes :
\begin{enumerate}
\item Son \cle{genre} et son \cle{nombre} sont ceux du nom auquel il se réfère (son antécédent) : \all{der Mann, der …} ; \all{die Frau, die …} ; \all{das Kind, das …} ; \all{die Leute, die …}.
\item Son \cle{cas} dépend de son rôle dans la proposition relative : sujet (nominatif), COD (accusatif), COI (datif), complément de préposition.
\end{enumerate}
Exemple clé : \all{der Mann, dem ich helfe} « l'homme que j'aide » — datif, parce que le verbe \all{helfen} exige le datif. De même : \all{die Witwe, der wir Geld geben} « la veuve à qui nous donnons de l'argent ».
\end{propriete}

\begin{attention}[Le français ne décline pas, l'allemand si]
En français, « qui » et « que » ne marquent ni le genre ni le cas : « la veuve \emph{qui} vit », « la veuve \emph{que} je vois », « la veuve \emph{à qui} je parle ». En allemand, il faut choisir à chaque fois la forme exacte : \all{die Witwe, die lebt} / \all{die Witwe, die ich sehe} / \all{die Witwe, mit der ich spreche}. Ne traduis jamais « que » automatiquement par \all{das} ou \all{die} : demande-toi d'abord quelle est la fonction du relatif dans sa propre proposition.
\end{attention}

\courssub{Tableau complet des pronoms relatifs}

Les formes du pronom relatif sont presque identiques à celles de l'article défini ; seules trois formes du datif et du génitif diffèrent (\all{denen}, \all{dessen}, \all{deren}).

\begin{center}
\begin{tabularx}{\linewidth}{|l|X|X|X|X|}
\hline
\rowcolor{popblueL} Cas & Masculin & Féminin & Neutre & Pluriel \\
\hline
Nominatif & der & die & das & die \\
\hline
Accusatif & den & die & das & die \\
\hline
Datif & dem & der & dem & denen \\
\hline
Génitif & dessen & deren & dessen & deren \\
\hline
\end{tabularx}
\end{center}

\begin{exemplebox}[Les quatre cas en action]
\begin{itemize}
\item Nominatif : \all{Der Witwer, der nebenan wohnt, heißt Herr Mbarga.} « Le veuf qui habite à côté s'appelle M. Mbarga. »
\item Accusatif : \all{Das Ritual, das wir gesehen haben, ist sehr alt.} « Le rituel que nous avons vu est très ancien. »
\item Datif : \all{Die Frau, der die Kinder helfen, ist Witwe.} « La femme que les enfants aident est veuve. »
\item Génitif : \all{Der Mann, dessen Frau gestorben ist, trauert.} « L'homme dont la femme est morte est en deuil. »
\end{itemize}
\end{exemplebox}

Avec une préposition, celle-ci se place \emph{devant} le relatif et impose son cas :

\begin{center}
\begin{tabularx}{\linewidth}{|X|X|X|}
\hline
\rowcolor{popblueL} Français & Deutsch & Remarque \\
\hline
avec lequel & mit dem (m./n.) & \all{mit} + datif \\
\hline
pour laquelle & für die (f.) & \all{für} + accusatif \\
\hline
par qui / desquels & von denen (pl.) & \all{von} + datif \\
\hline
dont (possession) & dessen / deren & génitif, sans préposition \\
\hline
\end{tabularx}
\end{center}

\all{Die Tradition, für die die Alten kämpfen, verschwindet.} « La tradition pour laquelle les anciens se battent disparaît. »

\all{Die Frauen, von denen die Riten befolgt werden, sind alt.} « Les femmes par qui les rites sont observés sont âgées. »

\begin{propriete}[Virgule obligatoire et verbe en fin]
Deux règles typographiques et syntaxiques absolues :
\begin{enumerate}
\item La proposition relative est \cle{toujours} précédée d'une virgule (et suivie d'une virgule si elle est insérée au milieu de la phrase).
\item Le verbe conjugué de la relative se place \cle{à la toute fin} de la subordonnée : \all{der Brauch, der gepflegt wird} ; au parfait, l'auxiliaire passe après le participe : \all{das Ritual, das wir gesehen haben}.
\end{enumerate}
\end{propriete}

\begin{popfigure}
\begin{tikzpicture}[
  box/.style={draw=popblue, thick, rounded corners, fill=popblueL, align=center, inner sep=6pt, font=\small},
  rbox/.style={draw=poporange, thick, rounded corners, fill=poporangeL, align=center, inner sep=6pt, font=\small},
  lbl/.style={font=\footnotesize\itshape, text=popmuted}
]
\node[box, align=center] (h) at (0,0){Hauptsatz :\\ Der Brauch \ldots{} ist alt.};
\node[rbox, align=center] (r) at (7.5,0){Relativsatz :\\ , den die Frauen \textbf{pflegen},};
\draw[-{Stealth}, very thick, poppurple] (r.north) to[bend left=25] node[above, font=\footnotesize, text=poppurple]{relatif : genre + nombre du nom, cas = fonction} (h.north);
\draw[-{Stealth}, very thick, popgreen] (8.6,-0.9) to[bend right=20] node[below, font=\footnotesize, text=popgreen]{verbe conjugué à la fin} (6.4,-0.9);
\node[lbl] at (0,-1.3) {phrase principale};
\node[lbl] at (7.5,-1.3) {subordonnée relative (virgule !)};
\end{tikzpicture}
\legende{Structure de la phrase avec proposition relative : le relatif renvoie au nom, le verbe conjugué se place en fin de subordonnée.}
\end{popfigure}

\begin{attention}[Les deux oublis classiques]
Ne jamais oublier la virgule avant la relative ni le verbe conjugué à la toute fin. Faux : \all{*Der Brauch der die Frauen pflegen ist alt.} Juste : \all{Der Brauch, den die Frauen pflegen, ist alt.} Autre piège : ne pas confondre le relatif avec l'article — dans \all{die Frauen, die …}, le premier \all{die} est article, le second est relatif.
\end{attention}

\courssub{Prétérit des verbes forts du texte}

Pour raconter un rite accompli autrefois — « le mari mourut, on coupa les cheveux de la veuve » — l'allemand écrit utilise le \cle{prétérit} (Präteritum). Les verbes forts changent de voyelle radicale et ne prennent aucune terminaison à la 1\iere{} et à la 3\ie{} personne du singulier.

\begin{textebox}[Aus der Geschichte eines Dorfes]
Als der Lehrer Nkoulou starb, kam die ganze Familie ins Dorf. Die Witwe trug ein weißes Tuch, und die Schwestern schnitten ihr die Haare. Man nahm ihr den Ehering und gab ihr ein einfaches Armband. Sieben Tage blieb sie im Haus und sprach mit niemandem. Die Alten verboten ihr, das Haus vor Sonnenaufgang zu verlassen. Viele Nachbarn kamen, brachten Essen und halfen den Kindern. Heute, sagen die Jungen, sind diese Riten zu streng; die Großmütter aber antworten: „Was wir von unseren Müttern nahmen, geben wir an unsere Töchter weiter.“ So blieb der Brauch lebendig, auch wenn sich das Leben im Dorf verändert hat.
\end{textebox}

Glossaire : \all{das Tuch, „er} (le tissu, le foulard) ; \all{der Ehering, -e} (l'alliance) ; \all{das Armband, „er} (le bracelet) ; \all{der Sonnenaufgang} (le lever du soleil) ; \all{weitergeben} (transmettre) ; \all{lebendig} (vivant).

Les verbes forts de ce récit forment exactement la liste à connaître :

\begin{center}
\begin{tabularx}{\linewidth}{|X|X|X|X|}
\hline
\rowcolor{popblueL} Infinitif & Prétérit (3\textsuperscript{e} pers.) & Participe passé & Français \\
\hline
sterben & starb & ist gestorben & mourir \\
\hline
tragen & trug & getragen & porter \\
\hline
schneiden & schnitt & geschnitten & couper \\
\hline
nehmen & nahm & genommen & prendre \\
\hline
geben & gab & gegeben & donner \\
\hline
kommen & kam & ist gekommen & venir \\
\hline
bleiben & blieb & ist geblieben & rester \\
\hline
verbieten & verbot & verboten & interdire \\
\hline
\end{tabularx}
\end{center}

\begin{propriete}[Conjugaison au prétérit : modèle tragen]
\all{tragen} (porter), radical du prétérit : \all{trug-}.
\begin{center}
\begin{tabular}{|l|l|l|}
\hline
\rowcolor{popblueL} Personne & Forme & Terminaison \\
\hline
ich & trug & --- \\
\hline
du & trugst & -st \\
\hline
er/sie/es & trug & --- \\
\hline
wir & trugen & -en \\
\hline
ihr & trugt & -t \\
\hline
sie/Sie & trugen & -en \\
\hline
\end{tabular}
\end{center}
Point capital : \cle{ich} et \cle{er/sie/es} n'ont aucune terminaison : \all{ich trug}, \all{sie trug} — comme \all{ich war}, \all{sie war}.
\end{propriete}

\begin{methode}[Apprendre un verbe fort : le triplet]
Pour trouver le prétérit d'un verbe fort, il faut connaître par cœur son \cle{triplet} infinitif – prétérit – participe : \all{tragen – trug – getragen}. Méthode en trois étapes :
\begin{enumerate}
\item Apprendre les triplets par familles de voyelles :
  \begin{itemize}
  \item \all{ei – ie – ie} : \all{bleiben – blieb – geblieben}, \all{schreiben – schrieb – geschrieben}, \all{treiben – trieb – getrieben}.
  \item \all{ei – i – i} : \all{schneiden – schnitt – geschnitten}, \all{beißen – biss – gebissen}, \all{leiden – litt – gelitten}.
  \item \all{e – a – o} : \all{nehmen – nahm – genommen}, \all{sterben – starb – gestorben}, \all{helfen – half – geholfen}. (Attention : \all{geben – gab – gegeben} fait \all{e – a – e}).
  \item \all{a – u – a} : \all{tragen – trug – getragen}, \all{fahren – fuhr – gefahren}, \all{schlagen – schlug – geschlagen}.
  \end{itemize}
\item Réciter chaque triplet à voix haute avec une phrase du thème : \all{Die Witwe trug Schwarz.}
\item Tester régulièrement les 15 verbes forts les plus fréquents au Bac : \all{sein, haben, werden, kommen, gehen, geben, nehmen, sehen, tragen, schneiden, bleiben, sterben, schreiben, sprechen, finden}.
\end{enumerate}
\end{methode}

\courssub{Participes passés et choix de l'auxiliaire}

\begin{propriete}[Formation du participe passé]
\begin{itemize}
\item Verbes forts : \cle{ge- + radical (souvent modifié) + -en} : \all{getragen, geschnitten, gegeben, genommen, geblieben}.
\item Verbes faibles : \cle{ge- + radical + -t} : \all{gepflegt, gemacht, getrauert}.
\item Verbes à préfixe inséparable (be-, ver-, er-, ent-, ge-, miss-, zer-) : \cle{sans ge-} : \all{begraben} (enterrer), \all{verboten} (interdit), \all{verändert} (changé), \all{zerstört} (détruit).
\item Exception utile : \all{erben} (hériter) est un verbe faible sans préfixe : \all{geerbt}.
\end{itemize}
\end{propriete}

\begin{propriete}[sein ou haben au parfait ?]
Le parfait se forme avec \all{haben} + participe, sauf pour :
\begin{itemize}
\item les verbes de \cle{mouvement} : \all{ist gekommen, ist gegangen, ist gefahren} ;
\item les verbes de \cle{changement d'état} : \all{ist gestorben, ist aufgewacht, ist geworden} ;
\item \all{bleiben} et \all{sein} : \all{ist geblieben, ist gewesen}.
\end{itemize}
Tous les autres prennent \all{haben} : \all{hat getragen, hat geschnitten, hat gegeben, hat verboten}.
\end{propriete}

\begin{attention}[« ist gestorben », pas « hat gestorben »]
Mourir est un changement d'état : l'allemand dit \all{Er ist gestorben.} Les francophones écrivent souvent \all{*hat gestorben} par calque du français « il a mouru » (lui-même incorrect : on dit « il est mort »). Retiens : \all{sterben} et \all{bleiben} prennent toujours \all{sein}, même si le français utilise « avoir » : \all{Sie ist im Dorf geblieben.} « Elle est restée au village. »
\end{attention}

\begin{exoresolu}[Relatives, prétérit et participes]
Énoncé. 1) Combine les deux phrases avec un pronom relatif : \all{Die Witwe trägt Schwarz. Die Witwe lebt im Dorf.} 2) Mets au prétérit : \all{Die Schwestern schneiden ihr die Haare.} 3) Mets au parfait : \all{Der Mann stirbt im Krankenhaus.} 4) Complète avec le bon relatif : \all{Der Brauch, \ldots{} die Frauen pflegen, ist alt.} 5) Traduis : « Les femmes avec lesquelles j'ai parlé sont très âgées. »
\tcblower
Solution détaillée. 1) Le nom répété est \all{die Witwe}, sujet dans les deux phrases, donc relatif au nominatif féminin : \all{Die Witwe, die im Dorf lebt, trägt Schwarz.} (virgules autour de la relative, verbe \all{lebt} en fin). 2) \all{schneiden – schnitt – geschnitten} ; 3\ie{} personne du pluriel : \all{Die Schwestern schnitten ihr die Haare.} 3) \all{sterben} est un changement d'état, donc auxiliaire \all{sein} : \all{Der Mann ist im Krankenhaus gestorben.} 4) \all{der Brauch} est masculin ; dans la relative, il est COD de \all{pflegen} (les femmes entretiennent quoi ?), donc accusatif masculin : \all{den}. Phrase : \all{Der Brauch, den die Frauen pflegen, ist alt.} 5) « avec » = \all{mit} + datif ; « les femmes » est pluriel, datif pluriel du relatif : \all{denen} ; au parfait, \all{sprechen} prend \all{haben} : \all{Die Frauen, mit denen ich gesprochen habe, sind sehr alt.}
\end{exoresolu}

\begin{exercice}[À toi de jouer]
1) Complète avec le pronom relatif correct : a) \all{Der Witwer, \ldots{} ich gestern besucht habe, heißt Etoundi.} b) \all{Die Tradition, von \ldots{} die Alten sprechen, ist uralt.} c) \all{Die Frau, \ldots{} Mann gestorben ist, trauert.} 2) Conjugue au prétérit : \all{kommen, geben, bleiben, tragen} (3\ie{} pers. sing.). 3) Forme le parfait : \all{verbieten, sterben, schneiden, erben}. 4) Traduis : « Le rituel que nous avons observé au village est très ancien. »
\end{exercice}

\begin{corrige}[Exercice : relatives et verbes forts]
1) a) \all{den} (masculin, COD de \all{besuchen}) ; b) \all{der} (féminin, datif après \all{von}) ; c) \all{deren} (génitif féminin, « dont le mari »). 2) \all{kam, gab, blieb, trug}. 3) \all{hat verboten} (préfixe inséparable, pas de \all{ge-}) ; \all{ist gestorben} (changement d'état) ; \all{hat geschnitten} ; \all{hat geerbt}. 4) \all{Das Ritual, das wir im Dorf beobachtet haben, ist sehr alt.} — relatif accusatif neutre \all{das}, auxiliaire \all{haben} en fin de relative.
\end{corrige}

\begin{aretenir}[L'essentiel de la section]
\begin{itemize}
\item Le pronom relatif prend le \cle{genre} et le \cle{nombre} de son antécédent, mais son \cle{cas} dépend de sa fonction dans la subordonnée.
\item Virgule obligatoire ; verbe conjugué à la fin de la relative.
\item Verbes forts : apprendre le triplet \all{tragen – trug – getragen}.
\item Participe fort : \all{ge-…-en} ; préfixe inséparable (be-, ver-, er-, ent-, ge-, miss-, zer-) : pas de \all{ge-} (\all{verboten, verändert}).
\item Parfait avec \all{sein} pour mouvement et changement d'état : \all{ist gestorben, ist geblieben, ist gekommen}.
\end{itemize}
\end{aretenir}

\coursec{Méthode du commentaire, commentaire modèle et production guidée}

Dans la section précédente, nous avons appris à construire les propositions relatives, à conjuguer les verbes forts au prétérit et à former les participes passés. Ces outils grammaticaux ne sont pas des fins en soi : ils servent à \emph{dire} et à \emph{écrire}. C'est exactement ce que l'on attend de vous au baccalauréat : lire un texte, le commenter, puis produire un court texte personnel. Cette dernière section vous donne la méthode complète, un modèle rédigé et des productions guidées, écrite et orale.

\courssub{La méthode du commentaire de texte en trois temps}

Commenter un texte allemand, ce n'est pas le traduire mot à mot ni le résumer sans ordre. C'est répondre à trois questions, dans un ordre fixe : \emph{de quoi s'agit-il ?} (présentation), \emph{que dit le texte et comment le dit-il ?} (analyse), \emph{que vaut ce texte et que nous apprend-il ?} (bilan).

\begin{methode}[Commenter un texte en trois temps]
\begin{enumerate}
\item \textbf{Einleitung — la présentation.} Présentez le texte en une ou deux phrases : l'auteur ou l'autrice (si connu), la source (journal, manuel, site), le type de texte (article, récit, reportage) et le thème général. Formulation type : \all{Der Text handelt von...} (le texte traite de...).
\item \textbf{Hauptteil — l'analyse.} Dégagez les idées principales dans l'ordre du texte, relevez le lexique du thème (ici : \all{die Witwe, die Trauer, das Ritual}) et signalez les structures grammaticales remarquables (ici : le passif et les propositions relatives). Formulation type : \all{Die Autorin beschreibt, wie...} (l'autrice décrit comment...).
\item \textbf{Schluss — le bilan.} Dites quel est l'intérêt du texte, donnez votre avis et ouvrez sur une comparaison interculturelle (Cameroun / pays germanophones). Formulation type : \all{Man kann feststellen, dass...} (on peut constater que...).
\end{enumerate}
\end{methode}

\begin{popfigure}
\begin{tikzpicture}[
  node distance=0.9cm,
  box/.style={draw=popblue, thick, rounded corners, fill=popblueL, align=center, text width=5.2cm, minimum height=1.1cm},
  arr/.style={-{Stealth[length=3mm]}, thick, poporange}
]
\node[box, align=center] (ein){\textbf{Einleitung}\\ Thema, Quelle, Textsorte};
\node[box, below=of ein, align=center] (haupt){\textbf{Hauptteil}\\ Inhalt (idées) + Sprache (lexique,\\ Passiv, Relativsätze)};
\node[box, below=of haupt, align=center] (schluss){\textbf{Schluss}\\ eigene Meinung + Vergleich\\ Kamerun / Deutschland};
\draw[arr] (ein) -- (haupt);
\draw[arr] (haupt) -- (schluss);
\end{tikzpicture}
\legende{Organigramme du commentaire de texte : trois étapes obligatoires}
\end{popfigure}

\begin{aretenir}[Les formulations types à mémoriser]
\begin{itemize}
\item \all{Der Text handelt von den Witwenritualen in Kamerun.} — Le texte traite des rites de veuvage au Cameroun.
\item \all{Die Autorin beschreibt die Schritte des Rituals.} — L'autrice décrit les étapes du rituel.
\item \all{Im Text werden viele Passivsätze benutzt.} — Dans le texte, beaucoup de phrases au passif sont utilisées.
\item \all{Man kann feststellen, dass die Traditionen streng befolgt werden.} — On peut constater que les traditions sont strictement observées.
\item \all{Der Text ist interessant, weil er zwei Kulturen vergleicht.} — Le texte est intéressant parce qu'il compare deux cultures.
\end{itemize}
Remarquez que \all{dass} et \all{weil} renvoient le verbe conjugué \textbf{en fin de proposition}.
\end{aretenir}

\courssub{Commentaire modèle rédigé}

Voici un commentaire complet du texte support de la leçon, \all{Witwenrituale in Kamerun und in Deutschland}. Lisez-le d'abord en repérant les trois temps de la méthode.

\begin{textebox}[Commentaire modèle]
Der Text handelt von den Witwenritualen in Kamerun und in Deutschland. Er ist ein informativer Artikel aus einem Schulbuch. Die Autorin beschreibt, wie die Haare der Witwe geschnitten werden und wie lange die Trauerzeit dauert. Sie erklärt auch, dass in Deutschland der Abschied meist kürzer und privater ist. Man kann feststellen, dass die Traditionen bei uns strenger befolgt werden als in Europa. Besonders interessant sind die vielen Passivsätze, zum Beispiel „Die Haare werden geschnitten.“, und die Relativsätze, die die Bräuche genauer erklären. Der Text ist wichtig, weil er zwei Kulturen respektvoll vergleicht. Meiner Meinung nach muss man die Rituale kennen, um die Familien zu verstehen, die sie praktizieren.
\end{textebox}

\textbf{Glossaire :} \all{der Abschied, -e} (les adieux, la séparation) ; \all{der Brauch, die Bräuche} (la coutume) ; \all{respektvoll} (avec respect) ; \all{praktizieren} (pratiquer) ; \all{meiner Meinung nach} (à mon avis).

\textbf{Analyse du modèle.} La première phrase présente (thème + type de texte) : c'est l'Einleitung. Les phrases avec \all{beschreibt} et \all{erklärt} analysent le contenu, puis une phrase relève la langue (passif + relatives) : c'est le Hauptteil. Les trois dernières phrases donnent l'intérêt du texte, la comparaison et l'avis personnel : c'est le Schluss. Comptez : environ neuf lignes, la longueur idéale au Bac.

\courssub{Production guidée écrite : décrire un rite de veuvage de ta région}

\textbf{Sujet.} Décris un rite de veuvage de ta région en 6 à 8 phrases : utilise le passif, au moins 2 propositions relatives et au moins 3 verbes forts au prétérit.

\begin{methode}[Plan phrase par phrase]
\begin{enumerate}
\item \textbf{Satz 1 :} allgemeine Einleitung — présente le rite en général (\all{Bei uns in der Region... gibt es ein Ritual, das...}).
\item \textbf{Sätze 2–5 :} die Schritte des Rituals im Passiv — décris les étapes au passif, dans l'ordre chronologique (\all{zuerst..., dann..., danach..., schließlich...}).
\item \textbf{Satz 6 :} ein Relativsatz — ajoute une précision avec une proposition relative (\all{die Witwe, die...}).
\item \textbf{Satz 7 :} Vergleich mit Deutschland — compare avec l'Allemagne (\all{Im Gegensatz zu Deutschland...}).
\item \textbf{Satz 8 :} eigene Meinung — donne ton avis (\all{Ich finde dieses Ritual wichtig, weil...}).
\end{enumerate}
\end{methode}

\begin{exemplebox}[Modèle de production rédigée]
\all{Bei uns in der Westregion gibt es ein Witwenritual, das sehr ernst genommen wird. Zuerst werden die Haare der Witwe geschnitten, dann wird sie von den älteren Frauen gewaschen. Danach wird weiße Kleidung getragen, und schließlich wird ein Fest für die Familie vorbereitet. Die Witwe, die oft sehr jung ist, bleibt mehrere Monate im Haus. Früher blieb die Witwe ein ganzes Jahr im Haus, und die Frau trug nur Weiß; man sprach nicht laut in ihrer Nähe. Im Gegensatz zu Deutschland, wo der Abschied kürzer ist, wird bei uns die Trauer streng befolgt. Ich finde dieses Ritual wichtig, weil es der Familie hilft, den Verlust gemeinsam zu tragen.}

« Dans notre région de l'Ouest, il existe un rite de veuvage qui est pris très au sérieux. D'abord, les cheveux de la veuve sont coupés, puis elle est lavée par les femmes âgées. Ensuite, des vêtements blancs sont portés, et finalement une fête est préparée pour la famille. La veuve, qui est souvent très jeune, reste plusieurs mois à la maison. Autrefois, la veuve restait une année entière à la maison, et la femme ne portait que du blanc ; on ne parlait pas fort près d'elle. Contrairement à l'Allemagne, où les adieux sont plus courts, chez nous le deuil est strictement observé. Je trouve ce rituel important parce qu'il aide la famille à porter la perte ensemble. »
\end{exemplebox}

Vérifions le contrat : passif (\all{werden geschnitten, wird gewaschen, wird getragen, wird vorbereitet, wird befolgt}) ; deux relatives (\all{das sehr ernst genommen wird} ; \all{die oft sehr jung ist}) ; trois verbes forts au prétérit (\all{blieb} de \all{bleiben}, \all{trug} de \all{tragen}, \all{sprach} de \all{sprechen}).

\begin{exoresolu}[Corriger une phrase de production]
Énoncé : un élève a écrit \all{Während die Trauerzeit in Deutschland ist kürzer, wird sie bei uns streng befolgt.} Trouvez et corrigez les deux erreurs.
\tcblower
Solution : (1) \all{während} introduit une subordonnée : le verbe conjugué va en fin — \all{während die Trauerzeit in Deutschland kürzer \textbf{ist}}. (2) La subordonnée \all{während...} étant en position 1, la principale qui suit doit commencer par le verbe : \all{\textbf{wird} sie bei uns streng befolgt}. La phrase corrigée (avec la subordonnée en tête) est : \all{Während die Trauerzeit in Deutschland kürzer ist, wird sie bei uns streng befolgt.} — « Alors que la période de deuil est plus courte en Allemagne, elle est strictement observée chez nous. » Variante (principale en tête) : \all{Die Trauerzeit wird bei uns streng befolgt, während sie in Deutschland kürzer ist.}
\end{exoresolu}

\begin{attention}[Le piège de « während »]
En français, « pendant que / alors que » est suivi d'une phrase au verbe en deuxième position mentale : « alors qu'elle est plus courte ». En allemand, \all{während} est une \textbf{conjonction de subordination} : le verbe conjugué part \textbf{à la fin} : \all{während sie in Deutschland kürzer \textbf{ist}}. Ne dites jamais \all{während sie ist kürzer}. Même vigilance avec \all{dass}, \all{weil} et \all{obwohl}.
\end{attention}

\textbf{Grille d'auto-évaluation.} Avant de rendre votre production, cochez chaque critère :

\begin{center}
\begin{tabularx}{\linewidth}{|X|l|}
\hline
\rowcolor{popblueL} \textbf{Critère} & \textbf{Points} \\
\hline
Au moins 5 mots du lexique du thème (Witwe, Trauer, Ritual, Brauch, Tradition...) & 2 \\
\hline
Passif correct : \all{werden} + participe II à la fin & 3 \\
\hline
Au moins 2 propositions relatives avec verbe en fin & 3 \\
\hline
Au moins 3 verbes forts au prétérit, correctement conjugués & 3 \\
\hline
Connecteurs chronologiques : \all{zuerst, dann, danach, schließlich} & 2 \\
\hline
Comparaison Cameroun/Allemagne + avis personnel & 2 \\
\hline
\end{tabularx}
\end{center}

\courssub{Production orale guidée : comparer en une minute}

Préparez une prise de parole d'une minute en suivant ce plan ; parlez lentement, en phrases complètes, sans lire.

\begin{experience}[Présentation orale : Kamerun und Deutschland im Vergleich]
\begin{enumerate}
\item \textbf{Einleitung (10 s) :} \all{Ich spreche über die Witwenrituale in Kamerun und in Deutschland.}
\item \textbf{Kamerun (20 s) :} deux étapes du rite au passif : \all{Bei uns werden die Haare der Witwe geschnitten. Die Trauerzeit wird streng befolgt.}
\item \textbf{Deutschland (15 s) :} \all{In Deutschland ist der Abschied kürzer und privater.}
\item \textbf{Schluss (15 s) :} avis personnel : \all{Ich finde beide Traditionen respektvoll, aber unsere Rituale helfen der Familie mehr.}
\end{enumerate}
Travaillez à deux : l'un parle, l'autre vérifie avec la grille (passif, relatives, prétérit), puis échangez les rôles.
\end{experience}

\courssub{Complément Bac : les structures de comparaison}

Pour viser l'excellence, enrichissez vos productions avec ces trois structures de comparaison :

\begin{center}
\begin{tabularx}{\linewidth}{|X|X|X|}
\hline
\rowcolor{popblueL} \textbf{Structure} & \textbf{Beispiel} & \textbf{Traduction} \\
\hline
\all{Im Gegensatz zu + Dativ} & \all{Im Gegensatz zu Deutschland dauert die Trauer bei uns länger.} & Contrairement à l'Allemagne, le deuil dure plus longtemps chez nous. \\
\hline
\all{während + verbe final} & \all{Bei uns wird getrauert, während in Deutschland das Leben weitergeht.} & Chez nous on est en deuil, alors qu'en Allemagne la vie continue. \\
\hline
\all{sowohl... als auch...} & \all{Sowohl in Kamerun als auch in Deutschland wird der Tod respektiert.} & Tant au Cameroun qu'en Allemagne, la mort est respectée. \\
\hline
\end{tabularx}
\end{center}

\begin{attention}[Deux pièges de la comparaison]
(1) Après \all{im Gegensatz zu}, le nom est au \textbf{datif} : \all{im Gegensatz zu \textbf{Deutschland}}, \all{im Gegensatz zu \textbf{den} deutschen Bräuchen}. (2) \all{sowohl... als auch...} encadre deux éléments \emph{parallèles} (deux lieux, deux actions) ; ne l'utilisez pas pour opposer, mais pour réunir.
\end{attention}

\begin{savaistu}[Comment la production écrite est notée au Bac]
Au baccalauréat, la production écrite est évaluée sur trois critères : la richesse du lexique, la correction grammaticale et la cohérence du texte. Les correcteurs récompensent explicitement la réutilisation des structures de la leçon : une phrase au passif bien construite et une proposition relative correcte rapportent des points de maîtrise de la langue. Inutile d'écrire long : huit phrases propres, variées et correctes valent mieux que quinze phrases risquées. Relisez toujours en vérifiant trois choses : la place du verbe, la majuscule aux noms, la forme correcte du participe (ge-...-t/en) et l'accord de l'auxiliaire \all{werden} au sujet.
\end{savaistu}

\begin{exercice}[À vous de jouer]
Rédigez votre propre production sur un rite de veuvage de votre région (6 à 8 phrases) en suivant le plan phrase par phrase. Utilisez obligatoirement : 5 mots du lexique, 3 formes de passif, 2 propositions relatives, 3 verbes forts au prétérit et une structure de comparaison du complément Bac. Auto-évaluez-vous avec la grille, puis présentez votre texte oralement en une minute à votre voisin.
\end{exercice}

\newpage

\coursec{Activité d'intégration}

\courssub{Objectifs de l'activité}

À l'issue de cette activité d'intégration, l'élève sera capable de :

\begin{itemize}
\item réinvestir dans une production authentique le \cle{vocabulaire} du deuil et du veuvage (\all{die Witwe}, \all{der Witwer}, \all{der Brauch}, \all{die Trauer}, \all{das Ritual}, \all{die Tradition}) ;
\item décrire un rite au \cle{passif} (\all{Vorgangspassiv}) avec des marqueurs chronologiques (\all{zuerst}, \all{dann}, \all{danach}, \all{schließlich}) ;
\item employer au moins deux \cle{propositions relatives} correctement accordées et trois \cle{verbes forts au prétérit} ;
\item formuler une \cle{comparaison interculturelle} entre le Cameroun et l'Allemagne à l'aide de \all{während} et \all{im Gegensatz zu} ;
\item présenter oralement un texte court en deux minutes, avec une prononciation claire, et écouter activement pour compléter un tableau collectif.
\end{itemize}

\courssub{Situation}

Le lycée de Yaoundé participe à un échange scolaire avec un Gymnasium de Hambourg. Dans le cadre de la semaine culturelle commune intitulée „Traditionen in unseren Ländern“, les partenaires allemands ont envoyé un courriel : ils préparent une petite exposition sur les rites du deuil en Allemagne et demandent aux élèves camerounais de leur faire découvrir, en retour, les rites de veuvage pratiqués dans leurs régions. Les meilleures fiches seront traduites, affichées à Hambourg et publiées sur le blog de l'échange.

\begin{textebox}[E-Mail der Partnerschule in Hamburg]
Liebe Freunde in Yaoundé,\\
für unsere Kulturwoche „Traditionen in unseren Ländern“ sammeln wir Texte über Trauerrituale. In Deutschland trägt die Witwe oder der Witwer oft schwarze Kleidung, und nach der Beerdigung wird ein Trauerbrief an die Familie geschickt. Aber wie ist das in Kamerun? Welche Rituale werden in euren Regionen befolgt? Wer organisiert sie? Wie lange dauern sie?\\
Schreibt uns kurze Steckbriefe (8--10 Sätze) auf Deutsch! Wir freuen uns auf eure Antworten.\\
Viele Grüße\\
Anna und Jonas, Klasse 11, Gymnasium Hamburg
\end{textebox}

\textbf{Glossaire :} \all{die Beerdigung, -en} : l'enterrement, les obsèques ; \all{der Trauerbrief, -e} : la lettre de condoléances ; \all{befolgen} : suivre, observer (un rite) ; \all{der Steckbrief, -e} : la fiche descriptive ; \all{dauern} : durer ; \all{sich freuen auf} + accusatif : se réjouir de (quelque chose à venir).

Votre mission : répondre à Anna et Jonas en rédigeant, par binômes, une fiche intitulée „Witwenrituale in meiner Region“, puis la présenter oralement à la classe.

\courssub{Consignes}

\begin{enumerate}
\item \textbf{Vorbereitung (10 min).} En binôme, choisissez une région du Cameroun (Ouest, Nord-Ouest, Centre, Littoral, Nord, etc.) et notez en allemand cinq mots-clés du lexique du deuil qui correspondent au rite que vous connaissez : \all{die Witwe}, \all{die Trauer}, \all{das Ritual}, \all{der Brauch}, \all{die Familie}, \all{die Kleidung}...

\item \textbf{Rédaction de la fiche (25 min).} Rédigez un texte de 8 à 10 phrases intitulé „Witwenrituale in meiner Region“. Contraintes obligatoires :
\begin{itemize}
\item au moins \textbf{quatre verbes au passif} pour décrire les étapes du rite, introduites par \all{zuerst}, \all{dann}, \all{danach}, \all{schließlich} ;
\item au moins \textbf{deux propositions relatives} (avec \all{die}, \all{der}, \all{das}, \all{denen}...) ;
\item au moins \textbf{trois verbes forts au prétérit} pour raconter un souvenir ou un exemple concret (\all{kam}, \all{trug}, \all{blieb}, \all{sang}, \all{sprach}...) ;
\item \textbf{une phrase de comparaison} avec l'Allemagne commençant par \all{Während...} ou \all{Im Gegensatz zu Deutschland...}.
\end{itemize}

\item \textbf{Relecture croisée (10 min).} Échangez votre fiche avec un autre binôme. Vérifiez à l'aide de la grille : le verbe conjugué est-il en deuxième position ? Le participe passé est-il à la fin au passif ? Le pronom relatif est-il au bon cas et au bon genre ? Y a-t-il une majuscule à tous les noms ?

\item \textbf{Présentation orale (2 min par binôme).} Présentez votre fiche à la classe. Répartition : le premier élève décrit le rite, le second raconte l'exemple au prétérit et formule la comparaison. Parlez lentement, sans lire mot à mot.

\item \textbf{Tableau collectif (10 min).} Pendant les exposés, la classe complète au tableau le tableau „Kamerun / Deutschland“ : qui porte le deuil ? quelle couleur ? quelle durée ? quel rôle de la famille ? Notez au moins deux différences et un point commun.

\item \textbf{Affichage.} Les meilleures fiches, corrigées et illustrées, sont affichées au mur de la salle d'allemand et envoyées par courriel aux partenaires de Hambourg.
\end{enumerate}

\courssub{Ressources à mobiliser}

\begin{aretenir}[Boîte à outils linguistique]
\textbf{1. Le passif pour décrire un rite} : \all{werden} + participe passé à la fin.\\
\all{Zuerst wird das Haus der Familie vorbereitet.} « D'abord, la maison de la famille est préparée. »\\
Au prétérit : \all{wurde} + participe. \all{Die Witwe wurde von den Tanten begleitet.} « La veuve était accompagnée par les tantes. »

\textbf{2. Les propositions relatives} : le pronom relatif s'accorde en genre et en nombre avec le nom auquel il se rapporte ; son cas dépend de sa fonction dans la relative ; le verbe conjugué va à la fin.\\
\all{Die Witwe, die schwarze Kleidung trägt, bleibt zu Hause.} « La veuve, qui porte des vêtements noirs, reste à la maison. »\\
\all{Das Ritual, das am Abend beginnt, dauert drei Tage.} « Le rituel, qui commence le soir, dure trois jours. »

\textbf{3. Prétérit des verbes forts} (à revoir dans la liste) : \all{kommen -- kam -- ist gekommen}, \all{tragen -- trug -- hat getragen}, \all{bleiben -- blieb -- ist geblieben}, \all{singen -- sang -- hat gesungen}, \all{sprechen -- sprach -- hat gesprochen}, \all{gehen -- ging -- ist gegangen}, \all{geben -- gab -- hat gegeben}.

\textbf{4. La comparaison} : \all{während} + verbe conjugué à la fin ; \all{im Gegensatz zu} + nom.\\
\all{Während in Deutschland die Trauer oft privat ist, wird sie in Kamerun mit der ganzen Gemeinde geteilt.} « Alors qu'en Allemagne le deuil est souvent privé, il est partagé au Cameroun avec toute la communauté. »
\end{aretenir}

\begin{attention}[Pièges à éviter]
\begin{itemize}
\item Ne confondez pas \all{werden} (auxiliaire du passif : « être ») et \all{bekommen} : \all{bekommen} signifie « recevoir », jamais « devenir ». Faux ami classique avec l'anglais \emph{to become}.
\item Au passif, l'auxiliaire \all{wird} est en deuxième position et le participe passé à la toute fin : on ne dit pas \all{wird vorbereitet das Haus}, mais \all{das Haus wird vorbereitet}.
\item Le pronom relatif n'est jamais omis en allemand, contrairement au français : « la femme que j'ai vue » = \all{die Frau, die ich gesehen habe}.
\item Attention à l'agent du passif : \all{von} + datif (\all{von den älteren Frauen}), et non \all{durch}, réservé au moyen ou à la cause.
\item Majuscule à tous les noms : \all{die Trauer}, \all{das Ritual}, \all{die Gemeinde}.
\end{itemize}
\end{attention}

\courssub{Proposition de production et corrigé commenté}

\begin{exoresolu}[Fiche modèle : „Witwenrituale in meiner Region“ (Ouest du Cameroun)]
Rédigez une fiche de 8 à 10 phrases respectant toutes les contraintes de la consigne 2.
\tcblower
\textbf{Production modèle :}

„Witwenrituale in meiner Region\\
In meiner Region, im Westen Kameruns, gibt es Rituale, die nach dem Tod des Mannes befolgt werden. Zuerst wird die Witwe, die meistens schwarze oder weiße Kleidung trägt, von den älteren Frauen der Familie begleitet. Dann wird das Haus gereinigt, und danach werden traditionelle Lieder gesungen. Während dieser Zeit wird der Witwe das Essen von den Nachbarn gebracht, damit sie nicht kochen muss. Schließlich wird nach einigen Monaten eine große Zeremonie organisiert, die das Ende der Trauer markiert.\\
Ich erinnere mich an meine Tante Marie: sie blieb vierzig Tage zu Hause, sie trug jeden Tag schwarze Kleidung, und ihre Freundinnen kamen jeden Abend und sangen mit ihr. Im Gegensatz zu Deutschland, wo die Trauer oft privat bleibt, wird sie bei uns mit der ganzen Gemeinde geteilt. Diese Tradition, die sehr alt ist, gibt der Witwe Kraft und Trost.“

\textbf{Commentaire des choix de langue :}
\begin{itemize}
\item \textbf{Passif} : sept occurrences (\all{werden befolgt}, \all{wird begleitet}, \all{wird gereinigt}, \all{werden gesungen}, \all{wird gebracht}, \all{wird organisiert}, \all{wird geteilt}), toujours avec le participe en fin de proposition et les marqueurs chronologiques demandés (\all{zuerst}, \all{dann}, \all{danach}, \all{schließlich}).
\item \textbf{Relatives} : \all{Rituale, die... befolgt werden} (nominatif pluriel), \all{die Witwe, die... trägt} (nominatif féminin singulier), \all{eine große Zeremonie, die das Ende der Trauer markiert}, \all{diese Tradition, die sehr alt ist} : quatre relatives, verbe conjugué à la fin.
\item \textbf{Prétérit des verbes forts} : \all{blieb} (\all{bleiben}), \all{trug} (\all{tragen}), \all{kamen} (\all{kommen}), \all{sangen} (\all{singen}) pour le souvenir personnel de la tante.
\item \textbf{Comparaison} : \all{Im Gegensatz zu Deutschland, wo...} avec la relative de lieu \all{wo} et le passif \all{wird geteilt}.
\item \textbf{Lexique du thème} : \all{die Witwe}, \all{die Trauer}, \all{das Ritual}, \all{die Zeremonie}, \all{die Gemeinde}, \all{der Trost}.
\end{itemize}
\end{exoresolu}

\begin{methode}[Grille d'évaluation de la fiche (sur 20 points)]
\begin{enumerate}
\item \textbf{Respect de la tâche} (4 pts) : 8 à 10 phrases, titre, réponse aux questions du courriel (quels rites ? qui organise ? quelle durée ?).
\item \textbf{Passif} (4 pts) : au moins quatre formes correctes, participe passé en fin de proposition.
\item \textbf{Propositions relatives} (4 pts) : au moins deux, pronom relatif au bon genre et au bon cas, verbe à la fin.
\item \textbf{Prétérit des verbes forts} (4 pts) : au moins trois formes exactes dans le récit du souvenir.
\item \textbf{Comparaison et lexique} (4 pts) : une phrase avec \all{während} ou \all{im Gegensatz zu}, cinq mots du champ lexical du deuil correctement employés, majuscules et orthographe soignées.
\end{enumerate}
\end{methode}

\courssub{Bilan de l'activité}

Cette activité a permis de mobiliser simultanément les quatre savoir-faire de la leçon : la compréhension (lecture du courriel des partenaires et écoute des exposés), le lexique thématique du deuil, la grammaire (passif descriptif, propositions relatives, prétérit des verbes forts) et la production écrite et orale dans une situation de communication authentique. Le tableau collectif „Kamerun / Deutschland“ a mis en évidence une leçon culturelle essentielle : dans les deux pays, le deuil obéit à des rites précis, mais la dimension communautaire est souvent plus forte au Cameroun, tandis qu'en Allemagne la \all{Trauer} est davantage vécue dans le cadre privé. Pour aller plus loin, les élèves pourront transformer leur fiche en article pour le blog de l'échange ou préparer une version orale enregistrée à envoyer à Hambourg — une étape concrète vers le dialogue interculturel, compétence clé du programme de Terminale.

\newpage

\coursec{Méthodes et exercices résolus}

\courssub{Méthodes et exercices résolus}

\begin{methode}[Comprendre un texte de société : lecture globale et détaillée]
\textbf{Objectif :} Répondre à des questions de compréhension (type Bac) sur un texte authentique traitant de faits de société (ici : les rites de veuvage au Cameroun et en Allemagne).\par
\textbf{Démarche en 5 étapes :}
\begin{enumerate}
    \item \textbf{Lecture globale (1\textsuperscript{re} lecture) :} Identifier le \cle{type de texte} (article de presse, rapport, témoignage), la \cle{source}, le \cle{thème général} et la \cle{thèse} ou l'\cle{intention} de l'auteur. Ne pas s'arrêter au vocabulaire inconnu.
    \item \textbf{Repérage des informations clés :} Souligner les \cle{mots-clés} (noms propres, chiffres, dates, connecteurs logiques : \all{aber}, \all{denn}, \all{weil}, \all{jedoch}) et les \cle{parties du texte} (introduction, développement, conclusion).
    \item \textbf{Lecture détaillée (2\textsuperscript{e} lecture) :} Lire par paragraphes. Pour chaque question : localiser la zone de réponse dans le texte, reformuler l'idée en français (brouillon), puis rédiger la réponse \textbf{en allemand} en reprenant le vocabulaire du texte (pas de copier-coller brut : adapter la grammaire — pronom, temps, cas).
    \item \textbf{Vérification grammaticale :} Contrôler l'accord sujet/verbe, le \cle{cas} (souvent accusatif après des prépositions comme \all{durch}, \all{für}, \all{gegen} ; datif après \all{mit}, \all{nach}, \all{von}), l'\cle{ordre des mots} (verbe conjugué en 2\textsuperscript{e} position en principale, à la fin en subordonnée).
    \item \textbf{Synthèse :} Pour une question du type \all{Fassen Sie zusammen...}, structurer en 3 phrases : Contexte $\rightarrow$ Problème/Description $\rightarrow$ Évolution/Comparaison/Opinion.
\end{enumerate}
\textbf{Réflexes Bac :} Citer la ligne (\all{In Zeile 5...}), utiliser le discours indirect (\all{Der Autor schreibt, dass...}), ne pas inventer d'informations extérieures au texte.
\end{methode}

\begin{exoresolu}[Compréhension du texte : « Witwenrituale in Kamerun und in Deutschland »]
\textbf{Énoncé :} Lisez le texte ci-après et répondez \textbf{en allemand} aux questions suivantes.
\begin{textebox}[Texte support : Witwenrituale in Kamerun und in Deutschland]
In vielen kamerunischen Kulturen ist der Tod eines Ehemannes nicht nur ein privater Verlust, sondern ein sozialer Umbruch. Die Witwe (die Witwe, die Witwen) muss oft eine mehrmonatige Trauerzeit (die Trauerzeit, die Trauerzeiten) einhalten. Während dieser Zeit trägt sie schwarze Kleidung, rasiert sich den Kopf und darf das Haus nicht verlassen. Sie wird von den Verwandten des Verstorbenen (der Verstorbene, die Verstorbenen) beobachtet, ob sie „schuldig“ am Tod ist. In manchen Regionen muss sie das Wasser trinken, mit dem der Leichnam gewaschen wurde (der Leichnam, die Leichname), um ihre Unschuld zu beweisen.

In Deutschland hingegen ist die Trauer Privatsache. Es gibt keine vorgeschriebenen Rituale. Die Witwe oder der Witwer (der Witwer, die Witwer) entscheidet selbst, wie lange er oder sie Schwarz trägt. Psychologische Beratung (die Beratung, die Beratungen) und Selbsthilfegruppen (die Selbsthilfegruppe, die Selbsthilfegruppen) helfen bei der Bewältigung. Der Staat sichert die finanzielle Absicherung durch die Witwenrente (die Witwenrente, die Witwenrenten) zu.

Dennoch gibt es Parallelen: In beiden Gesellschaften wird der Verstorbene geehrt. In Kamerun durch große Beerdigungsfeiern (die Beerdigungsfeier, die Beerdigungsfeiern) mit Hunderten von Gästen, in Deutschland durch stille Trauerfeiern im engsten Kreis. Der soziale Druck auf die Witwe ist in Kamerun jedoch deutlich höher.
\end{textebox}

\begin{enumerate}
    \item Laut Text: Welche drei Handlungen muss die Witwe in Kamerun während der Trauerzeit ausführen? (Nennen Sie drei Details.)
    \item Wie wird die „Unschuld“ der Witwe in manchen Regionen Kameruns geprüft? Beschreiben Sie das Ritual.
    \item Nennen Sie zwei Unterschiede zwischen der Trauer in Kamerun und in Deutschland (sozialer Aspekt / finanzielle Absicherung).
    \item Erklären Sie den Satz: „Der soziale Druck auf die Witwe ist in Kamerun jedoch deutlich höher.“ Begründen Sie mit zwei Beispielen aus dem Text.
    \item \textbf{Expression personnelle (3-4 Sätze) :} Was halten Sie von der Aussage: „Traditionen schützen die Witwe, aber sie beschränken auch ihre Freiheit“? Beziehen Sie sich auf den Text.
\end{enumerate}
\tcblower
\textbf{Solution détaillée :}

\textbf{Question 1 (Repérage explicite — lignes 3-5) :}
\textbf{Antwort :} Die Witwe muss schwarze Kleidung tragen, sich den Kopf rasieren und das Haus nicht verlassen.
\emph{Justification :} Le texte liste trois obligations : \all{trägt sie schwarze Kleidung, rasiert sich den Kopf und darf das Haus nicht verlassen}. Réponse en phrase complète. Attention à la conjugaison (\all{muss ... tragen / rasieren / verlassen}, infinitifs à la fin) et aux articles (\all{die Kleidung}, \all{der Kopf}, \all{das Haus}).

\textbf{Question 2 (Détail précis — lignes 5-7) :}
\textbf{Antwort :} In manchen Regionen muss die Witwe das Wasser trinken, mit dem der Leichnam gewaschen wurde, um ihre Unschuld zu beweisen.
\emph{Justification :} Reprise de la structure \all{muss ... trinken} + proposition relative \all{mit dem der Leichnam gewaschen wurde} (passif \all{Vorgangspassiv} au Präteritum : \all{wurde ... gewaschen}). Le pronom relatif \all{dem} est au datif (préposition \all{mit} + datif).

\textbf{Question 3 (Comparaison — structuration) :}
\textbf{Antwort :}
\begin{enumerate}
    \item \textbf{Sozialer Aspekt:} In Kamerun ist die Trauer eine gemeinschaftliche Pflicht mit strengen Ritualen (Beobachtung durch die Verwandten); in Deutschland ist die Trauer Privatsache ohne vorgeschriebene Rituale.
    \item \textbf{Finanzielle Absicherung:} In Deutschland sichert der Staat die Witwe durch die Witwenrente ab; im Text wird für Kamerun keine staatliche Rente erwähnt — die Absicherung erfolgt traditionell über die Familie und die Verwandten.
\end{enumerate}
\emph{Justification :} Utiliser des connecteurs de comparaison (\all{während}, \all{hingegen}, \all{im Gegensatz dazu}). Vocabulaire clé : \all{die Pflicht}, \all{die Selbsthilfegruppe}, \all{die Witwenrente}, \all{die Absicherung}. Attention : \all{absichern} est un verbe à particule séparable (\all{sichert ... ab}).

\textbf{Question 4 (Inférence / justification — « sozialer Druck ») :}
\textbf{Antwort :} Der Druck ist höher, weil erstens die Witwe von den Verwandten des Verstorbenen beobachtet wird, ob sie „schuldig“ am Tod ist (Verdacht der Mitschuld), und weil sie zweitens strenge Verhaltensregeln einhalten muss (Kleidung, Haare, Hausarrest), während in Deutschland keine solchen Regeln existieren.
\emph{Justification :} Lier l'adjectif \all{höher} à des preuves textuelles. \all{weil} renvoie le verbe conjugué à la fin (\all{existieren}). \all{des Verstorbenen} : participe II substantivé, génitif singulier, décliné comme un adjectif.

\textbf{Question 5 (Expression guidée — opinion + référence au texte) :}
\textbf{Antwort :} Ich stimme der Aussage zu. Der Text zeigt, dass die Tradition in Kamerun die Witwe sozial einbindet (Schutz durch die Gemeinschaft bei der Beerdigung), aber ihre Freiheit stark einschränkt (Hausarrest, Zwangsrituale wie das Trinken des Leichenwassers). In Deutschland hat die Witwe mehr Freiheit, aber weniger traditionellen sozialen Halt. Ein Gleichgewicht wäre ideal.
\emph{Justification :} Structure : \all{Ich stimme zu / nicht zu} + \all{Der Text zeigt, dass...} (subordonnée en \all{dass}, verbe à la fin) + arguments contrastés (\all{aber}, \all{jedoch}) + conclusion. Vocabulaire : \all{einschränken} (restreindre), \all{das Zwangsrual} (pluriel \all{die Zwangsrituale}), \all{der soziale Halt} (le soutien social).
\end{exoresolu}

\begin{methode}[Maîtriser le lexique thématique : familles de mots, articles, pluriels]
\textbf{Objectif :} Construire et utiliser le vocabulaire du deuil, du veuvage et des traditions (\cle{Wortschatz}) avec exactitude morphologique (genre, pluriel, Umlaut) et sémantique (faux amis, registres).\par
\textbf{Démarche en 4 étapes :}
\begin{enumerate}
    \item \textbf{Fiche lexicale catégorisée :} Classer par champs : \textit{Personnes} (\all{die Witwe}, \all{der Witwer}, \all{der/die Verstorbene}), \textit{Actions/Rites} (\all{das Ritual}, \all{der Brauch}, \all{die Beerdigung}, \all{rasieren}, \all{bewältigen}), \textit{Sentiments/États} (\all{die Trauer}, \all{der Schmerz}, \all{die Schuld}, \all{unschuldig}), \textit{Institutions} (\all{die Witwenrente}, \all{die Selbsthilfegruppe}, \all{die Beratung}).
    \item \textbf{Systématiser Article + Pluriel + Umlaut :} Apprendre \textbf{toujours} le triplet : \all{die Witwe, die Witwen} (pas d'Umlaut) ; \all{der Brauch, die Bräuche} (Umlaut \cle{au $\rightarrow$ äu}) ; \all{das Ritual, die Rituale} (pluriel en \cle{-e}) ; \all{der Verstorbene, die Verstorbenen} (adjectif substantivé, décliné comme un adjectif).
    \item \textbf{Familles de mots (dérivation) :} Identifier la racine et les suffixes : \all{trauern} (v) $\rightarrow$ \all{die Trauer} (n) $\rightarrow$ \all{traurig} (adj) $\rightarrow$ \all{die Traurigkeit} (n). \all{sterben} $\rightarrow$ \all{der Tod}, \all{der Verstorbene}, \all{sterblich}. \all{beweisen} $\rightarrow$ \all{der Beweis}, \all{bewiesen}.
    \item \textbf{Collocations (mots « amis ») :} Mémoriser les couples inséparables : \all{die Trauerzeit einhalten}, \all{eine Beerdigung abhalten}, \all{die Unschuld beweisen}, \all{eine Rente beziehen}, \all{die Trauer bewältigen}.
\end{enumerate}
\textbf{Pièges francophones :} \all{der Brauch} (la coutume) $\neq$ \all{brauchen} (avoir besoin) ; \all{die Leiche} (le cadavre, registre familier/péjoratif) $\neq$ \all{der Leichnam} (la dépouille, registre formel et respectueux) ; \all{beerdigen} (enterrer) $\rightarrow$ Partizip II \all{beerdigt} (verbe faible, sans \all{ge-} à cause du préfixe inséparable \all{be-}), mais \all{sterben} $\rightarrow$ \all{gestorben} (verbe fort).
\end{methode}

\begin{exoresolu}[Entraînement lexical : dérivation, collocations, articles/pluriels]
\textbf{Énoncé :}
\begin{enumerate}
    \item Complétez le tableau de dérivation. Donnez l'article, le pluriel et la traduction française. Corrigez le pluriel erroné signalé par *.
    \begin{center}
    \begin{tabular}{|l|l|l|l|}\hline
    \rowcolor{popblueL} \textbf{Verbe} & \textbf{Nom (Art. + Pl.)} & \textbf{Adjectif / Participe} & \textbf{Traduction FR} \\ \hline
    \all{trauern} & \all{die Trauer, die Traueren*} & \all{traurig} & chagrin / triste \\ \hline
    \all{sterben} & \all{der Tod, die Tode} & \all{gestorben / tödlich} & mort / mortel \\ \hline
    \all{bewältigen} & \all{...} & \all{bewältigt} & surmonter / surmonté \\ \hline
    \all{beweisen} & \all{der Beweis, die Beweise} & \all{...} & preuve / prouvé \\ \hline
    \all{rasieren} & \all{die Rasur, die Rasuren} & \all{rasiert} & rasage / rasé \\ \hline
    \end{tabular}
    \end{center}
    \item Choisissez le mot juste (collocation) et conjuguez au temps indiqué (Présent ou Prétérit).
    \begin{enumerate}
        \item Die Witwe \_\_\_\_\_\_\_\_\_\_ (die Trauerzeit / einhalten) seit drei Monaten. (Présent)
        \item Gestern \_\_\_\_\_\_\_\_\_\_ der Verstorbene \_\_\_\_\_\_\_\_\_\_ (beerdigen / werden). (Prétérit, Passif)
        \item Die Verwandten \_\_\_\_\_\_\_\_\_\_ (beobachten) sie, ob sie schuldig \_\_\_\_\_\_\_\_\_\_ (sein) am Tod. (Présent)
        \item In Deutschland \_\_\_\_\_\_\_\_\_\_ die Witwe die Witwenrente \_\_\_\_\_\_\_\_\_\_ (beziehen). (Présent)
    \end{enumerate}
    \item Traduisez en allemand (thème) en utilisant le vocabulaire de la leçon :
    \begin{enumerate}
        \item La veuve doit respecter les coutumes traditionnelles.
        \item Le deuil est une affaire privée en Allemagne.
        \item Elle a prouvé son innocence en buvant l'eau.
    \end{enumerate}
\end{enumerate}
\tcblower
\textbf{Solution détaillée :}

\textbf{1. Tableau de dérivation (corrections en gras) :}
\begin{center}
\begin{tabular}{|l|l|l|l|}\hline
\rowcolor{popblueL} \textbf{Verbe} & \textbf{Nom (Art. + Pl.)} & \textbf{Adjectif / Participe} & \textbf{Traduction} \\ \hline
\all{trauern} & \textbf{\all{die Trauer}} (\textbf{sans pluriel usuel}) & \all{traurig} & le chagrin / triste \\ \hline
\all{sterben} & \all{der Tod, die Tode} (pluriel rare) & \all{gestorben / tödlich} & la mort / mortel \\ \hline
\all{bewältigen} & \textbf{\all{die Bewältigung, -en}} (suffixe \cle{-ung} $\rightarrow$ féminin) & \all{bewältigt} & la maîtrise / maîtrisé \\ \hline
\all{beweisen} & \all{der Beweis, die Beweise} & \textbf{\all{bewiesen}} (Partizip II) & la preuve / prouvé \\ \hline
\all{rasieren} & \all{die Rasur, die Rasuren} & \all{rasiert} & le rasage / rasé \\ \hline
\end{tabular}
\end{center}
\emph{Points de vigilance :} \all{die Trauer} est féminin et n'a pas de pluriel usuel (on dira \all{Trauerfälle} selon le contexte) ; \all{die Bewältigung} : les nominalisations en \cle{-ung} sont toujours féminines, pluriel en \cle{-en} ; \all{bewiesen} : \all{beweisen} est un verbe fort (\all{beweist — bewies — hat bewiesen}).

\textbf{2. Collocations + conjugaison :}
\begin{enumerate}
    \item Die Witwe \textbf{\all{hält die Trauerzeit seit drei Monaten ein}}. (Verbe à particule séparable \all{einhalten} $\rightarrow$ \all{hält ... ein} ; \all{halten} est fort : \all{hält}.)
    \item Gestern \textbf{\all{wurde}} der Verstorbene \textbf{\all{beerdigt}}. (Passif \all{Vorgangspassiv} au Prétérit : \all{werden} $\rightarrow$ \all{wurde} + Partizip II \all{beerdigt} à la fin. Sujet \all{der Verstorbene} au nominatif.)
    \item Die Verwandten \textbf{\all{beobachten}} sie, ob sie schuldig \textbf{\all{ist}}. (\all{beobachten} au présent ; subordonnée en \all{ob}, verbe conjugué \all{ist} à la fin.)
    \item In Deutschland \textbf{\all{bezieht}} die Witwe die Witwenrente. (Verbe \all{beziehen} au présent, 3\textsuperscript{e} pers. sg. \all{bezieht} ; ordre V2 : \all{In Deutschland bezieht die Witwe...}.)
\end{enumerate}

\textbf{3. Thème (traduction FR $\rightarrow$ DE) :}
\begin{enumerate}
    \item \textbf{\all{Die Witwe muss die traditionellen Bräuche einhalten.}} (\all{die Witwe} féminin ; \all{müssen} + infinitif \all{einhalten} à la fin ; \all{die Bräuche} pluriel avec Umlaut ; adjectif \all{traditionellen} décliné à l'accusatif pluriel.)
    \item \textbf{\all{Die Trauer ist in Deutschland Privatsache.}} (\all{die Trauer} sujet ; \all{Privatsache} attribut sans article.)
    \item \textbf{\all{Sie hat ihre Unschuld bewiesen, indem sie das Wasser trank.}} (Structure \all{indem} + Prétérit : \all{trank}, verbe fort \all{trinken}. Variante : \all{Sie hat ihre Unschuld durch das Trinken des Wassers bewiesen} — préposition \all{durch} + nominalisation \all{das Trinken} + génitif \all{des Wassers}.)
\end{enumerate}
\end{exoresolu}

\begin{methode}[Le passif de procès (Vorgangspassiv) : décrire un rite, une procédure]
\textbf{Objectif :} Utiliser le \cle{Vorgangspassiv} (\all{werden} + \all{Partizip II}) pour mettre l'accent sur l'action subie (le rite) plutôt que sur l'agent (qui l'accomplit) — construction typique des descriptions rituelles, administratives et journalistiques.\par
\textbf{Formation (rappel) :} \textbf{Sujet (patient) + werden (conjugué) + ... + Partizip II (à la fin).} L'agent (optionnel) est introduit par \all{von} (personne) ou \all{durch} (moyen/instrument).\par
\textbf{Tableau des temps (niveau Terminale) :}
\begin{center}
\begin{tabular}{|l|l|l|l|}\hline
\rowcolor{popblueL} \textbf{Temps} & \textbf{Structure} & \textbf{Exemple (die Witwe / beobachten)} & \textbf{Traduction} \\ \hline
Présent & \all{werden} + Part. II & \all{Die Witwe wird beobachtet.} & La veuve est observée. \\ \hline
Prétérit & \all{wurden} + Part. II & \all{Die Witwe wurde beobachtet.} & La veuve fut/était observée. \\ \hline
Parfait & \all{sein} + Part. II + \all{worden} & \all{Die Witwe ist beobachtet worden.} & La veuve a été observée. \\ \hline
Futur I & \all{werden} + Part. II + \all{werden} & \all{Die Witwe wird beobachtet werden.} & La veuve sera observée. \\ \hline
Modal (prés.) & \all{Modal} + Part. II + \all{werden} & \all{Die Witwe muss beobachtet werden.} & La veuve doit être observée. \\ \hline
\end{tabular}
\end{center}
\textbf{Démarche de transformation Actif $\rightarrow$ Passif (3 étapes) :}
\begin{enumerate}
    \item Identifier l'\textbf{objet accusatif} de la phrase active $\rightarrow$ il devient \textbf{sujet nominatif} du passif (accord du verbe !).
    \item Le \textbf{sujet} actif $\rightarrow$ devient \textbf{complément d'agent} (\all{von} + datif) ou est supprimé (s'il est indéfini : \all{man}, \all{jemand}, \all{Leute}).
    \item Mettre le verbe à la forme \textbf{werden + Partizip II} en conservant le temps du verbe actif.
\end{enumerate}
\textbf{Attention :} Les verbes intransitifs (sans accusatif) n'ont pas de passif standard (sauf le passif impersonnel \all{Es wird getanzt}, hors programme). Pour les verbes à double complément (\all{jemandem etwas rasieren}), seul l'accusatif de la chose devient sujet ; la personne reste au datif.
\end{methode}

\begin{exoresolu}[Vorgangspassiv : transformation et rédaction rituelle]
\textbf{Énoncé :}
\begin{enumerate}
    \item Transformez les phrases actives en passif de procès (\all{Vorgangspassiv}). Conservez le temps. Supprimez l'agent si c'est \all{man} ou \all{die Verwandten} (connu par le contexte).
    \begin{enumerate}
        \item Die Verwandten waschen den Leichnam. (Présent)
        \item Man rasiert der Witwe den Kopf. (Présent — \textbf{double complément !})
        \item Die Ältesten haben das Ritual erklärt. (Parfait)
        \item Man wird die Witwe nicht zwingen. (Futur I)
    \end{enumerate}
    \item Complétez le paragraphe descriptif en conjuguant les verbes entre parenthèses au \textbf{passif} (temps au choix selon le sens : présent, prétérit, parfait).
    \begin{textebox}[Lückentext]
    In der Tradition \_\_\_\_\_\_\_\_\_\_ der Leichnam \_\_\_\_\_\_\_\_\_\_ (waschen) und \_\_\_\_\_\_\_\_\_\_ (einbalsamieren). Danach \_\_\_\_\_\_\_\_\_\_ die Witwe in das Trauerhaus \_\_\_\_\_\_\_\_\_\_ (führen). Sie \_\_\_\_\_\_\_\_\_\_ von Männern nicht \_\_\_\_\_\_\_\_\_\_ (sehen). Das Wasser, mit dem der Körper \_\_\_\_\_\_\_\_\_\_ (waschen), \_\_\_\_\_\_\_\_\_\_ von der Witwe \_\_\_\_\_\_\_\_\_\_ (müssen / trinken). Heute \_\_\_\_\_\_\_\_\_\_ diese Rituale von Menschenrechtsorganisationen kritisch \_\_\_\_\_\_\_\_\_\_ (sehen).
    \end{textebox}
    \item \textbf{Thème :} Traduisez en allemand (passif).
    \begin{enumerate}
        \item La veuve est observée par la famille. (Présent)
        \item Ses cheveux ont été rasés. (Parfait — sujet \all{die Haare})
        \item Le rite sera bientôt interdit. (Futur I)
    \end{enumerate}
\end{enumerate}
\tcblower
\textbf{Solution détaillée :}

\textbf{1. Transformation Actif $\rightarrow$ Passif :}
\begin{enumerate}
    \item \textbf{\all{Der Leichnam wird (von den Verwandten) gewaschen.}} (Accusatif \all{den Leichnam} $\rightarrow$ sujet \all{der Leichnam}, masculin singulier $\rightarrow$ \all{wird}. L'agent \all{von den Verwandten} est possible mais omis ici.)
    \item \textbf{\all{Der Witwe wird der Kopf rasiert.}} (\textbf{Piège du double complément} : \all{man rasiert [wem?] der Witwe [was?] den Kopf}. Seul l'accusatif de la chose (\all{der Kopf}) devient sujet nominatif ; la personne (\all{der Witwe}) reste au \textbf{datif} au passif.)
    \item \textbf{\all{Das Ritual ist (von den Ältesten) erklärt worden.}} (Parfait actif \all{haben erklärt} $\rightarrow$ parfait passif \all{sein ... worden}. Partizip II \all{erklärt} + \all{worden} sans \all{ge-}.)
    \item \textbf{\all{Die Witwe wird nicht gezwungen werden.}} (Futur I actif \all{werden zwingen} $\rightarrow$ futur I passif \all{werden gezwungen werden}. Double \all{werden} : auxiliaire de futur conjugué + auxiliaire de passif à l'infinitif.)
\end{enumerate}

\textbf{2. Lückentext (contexte rituel : prétérit majoritaire pour le récit passé, présent pour la vérité générale) :}
\begin{textebox}[Solution]
In der Tradition \textbf{\all{wurde}} der Leichnam \textbf{\all{gewaschen}} und \textbf{\all{einbalsamiert}}. (Prétérit : récit passé.) Danach \textbf{\all{wurde}} die Witwe in das Trauerhaus \textbf{\all{geführt}}. (Prétérit ; \all{führen} est faible : \all{geführt}.) Sie \textbf{\all{wurde}} von Männern nicht \textbf{\all{gesehen}}. (Prétérit ; \all{sehen} fort : \all{gesehen}.) Das Wasser, mit dem der Körper \textbf{\all{gewaschen wurde}}, \textbf{\all{muss}} von der Witwe \textbf{\all{getrunken werden}}. (Relative : prétérit passif \all{wurde gewaschen} ; principale : modal au présent passif \all{muss ... getrunken werden}.) Heute \textbf{\all{werden}} diese Rituale von Menschenrechtsorganisationen kritisch \textbf{\all{gesehen}}. (Présent : actualité.)
\end{textebox}

\textbf{3. Thème :}
\begin{enumerate}
    \item \textbf{\all{Die Witwe wird von der Familie beobachtet.}} (Passif présent standard.)
    \item \textbf{\all{Ihre Haare sind rasiert worden.}} (Parfait passif. \all{Ihre Haare} sujet pluriel $\rightarrow$ \all{sind}. Variante avec datif d'appartenance : \all{Ihr sind die Haare rasiert worden.})
    \item \textbf{\all{Das Ritual wird bald verboten werden.}} (Futur I passif : \all{wird} + \all{verboten} + \all{werden}.)
\end{enumerate}
\end{exoresolu}

\begin{methode}[Propositions relatives, prétérit des verbes forts et participes II : structurer le récit]
\textbf{Objectif :} Enchaîner les informations (relative) et situer l'action dans le passé narratif (prétérit) avec les formes correctes des verbes forts et irréguliers.\par
\textbf{A. La proposition relative (Relativsatz) :}
\begin{enumerate}
    \item \textbf{Antécédent + pronom relatif (décliné selon sa FONCTION dans la relative) + ... + verbe conjugué À LA FIN.}
    \item Tableau des relatifs (identiques aux articles définis, sauf au \textbf{génitif} et au \textbf{datif pluriel}) :
    \begin{center}
    \begin{tabular}{|c|c|c|c|c|}\hline
    \rowcolor{popblueL} Kasus & Maskulinum & Femininum & Neutrum & Plural \\ \hline
    Nominativ & \all{der} & \all{die} & \all{das} & \all{die} \\ \hline
    Akkusativ & \all{den} & \all{die} & \all{das} & \all{die} \\ \hline
    Dativ & \all{dem} & \all{der} & \all{dem} & \textbf{\all{denen}} \\ \hline
    Genitiv & \textbf{\all{dessen}} & \textbf{\all{deren}} & \textbf{\all{dessen}} & \textbf{\all{deren}} \\ \hline
    \end{tabular}
    \end{center}
    \item \textbf{Préposition + relatif :} la préposition précède le relatif (\all{mit dem}, \all{für den}, \all{an die}, \all{wegen dessen}).
\end{enumerate}
\textbf{B. Le prétérit (Präteritum) des verbes forts (unregelmäßige Verben) :}
\begin{itemize}
    \item Pas de \cle{-te} ! La voyelle radicale change (Ablaut) : \all{sterben $\rightarrow$ starb}, \all{trinken $\rightarrow$ trank}, \all{geben $\rightarrow$ gab}, \all{sehen $\rightarrow$ sah}, \all{werden $\rightarrow$ wurde}, \all{sein $\rightarrow$ war}, \all{haben $\rightarrow$ hatte}.
    \item Terminaisons : \cle{-} (ich/er/sie/es), \cle{-st} (du), \cle{-en} (wir/Sie/sie), \cle{-t} (ihr).
    \item \textbf{Liste Bac essentielle à connaître par cœur :} \all{gehen (ging)}, \all{kommen (kam)}, \all{stehen (stand)}, \all{finden (fand)}, \all{geben (gab)}, \all{nehmen (nahm)}, \all{sprechen (sprach)}, \all{schreiben (schrieb)}, \all{lesen (las)}, \all{trinken (trank)}, \all{sterben (starb)}, \all{werden (wurde)}, \all{sein (war)}, \all{haben (hatte)}, \all{können (konnte)}, \all{müssen (musste)}, \all{wollen (wollte)}, \all{sollen (sollte)}, \all{dürfen (durfte)}, \all{mögen (mochte)}.
\end{itemize}
\textbf{C. Le participe II (Partizip II) :}
\begin{itemize}
    \item Verbes faibles : \cle{ge- + radical + -t} (\all{ge-mach-t} ; \all{be-obacht-et} sans \all{ge-} à cause du préfixe inséparable).
    \item Verbes forts : \cle{ge- + radical (souvent modifié) + -en} (\all{ge-seh-en}, \all{ge-trunk-en}, \all{ge-storb-en}, \all{ge-word-en}).
    \item Verbes en \cle{-ieren} et verbes à préfixe inséparable (\all{be-, ver-, ent-, er-, ge-, miss-, zer-}) : \textbf{pas de \all{ge-}} (\all{studiert}, \all{besucht}, \all{verstanden}).
    \item Verbes à particule séparable : \cle{ge- entre le préfixe et le radical} (\all{auf-ge-hört}, \all{ein-ge-halten}).
\end{itemize}
\end{methode}

\begin{exoresolu}[Relatives, prétérit des verbes forts, participes : exercices de synthèse]
\textbf{Énoncé :}
\begin{enumerate}
    \item \textbf{Relatives :} Reliez les deux phrases par une proposition relative (remplacez l'élément en gras par le relatif correctement décliné). Placez la préposition devant le relatif si nécessaire.
    \begin{enumerate}
        \item Die Witwe trinkt das Wasser. Der Leichnam wurde mit \textbf{dem Wasser} gewaschen.
        \item Die Verwandten beobachten die Witwe. \textbf{Die Witwe} trägt die schwarze Kleidung.
        \item Der Brauch ist alt. Der Ethnologe hat \textbf{von dem Brauch} berichtet.
        \item Die Witwenrente sichert das Überleben. \textbf{Ohne die Witwenrente} wäre die Witwe arm.
    \end{enumerate}
    \item \textbf{Prétérit / Participe II :} Complétez le tableau pour les verbes forts et irréguliers du thème.
    \begin{center}
    \begin{tabular}{|l|c|c|c|}\hline
    \rowcolor{popblueL} \textbf{Infinitif} & \textbf{Présent (3e sg)} & \textbf{Prétérit (3e sg)} & \textbf{Participe II} \\ \hline
    \all{sterben} & \all{stirbt} & \_\_\_\_\_\_\_\_\_\_ & \_\_\_\_\_\_\_\_\_\_ \\ \hline
    \all{trinken} & \all{trinkt} & \_\_\_\_\_\_\_\_\_\_ & \_\_\_\_\_\_\_\_\_\_ \\ \hline
    \all{werden} & \all{wird} & \_\_\_\_\_\_\_\_\_\_ & \_\_\_\_\_\_\_\_\_\_ \\ \hline
    \all{sein} & \all{ist} & \_\_\_\_\_\_\_\_\_\_ & \_\_\_\_\_\_\_\_\_\_ \\ \hline
    \all{geben} & \all{gibt} & \_\_\_\_\_\_\_\_\_\_ & \_\_\_\_\_\_\_\_\_\_ \\ \hline
    \all{bewältigen} & \all{bewältigt} & \_\_\_\_\_\_\_\_\_\_ & \_\_\_\_\_\_\_\_\_\_ \\ \hline
    \all{rasieren} & \all{rasiert} & \_\_\_\_\_\_\_\_\_\_ & \_\_\_\_\_\_\_\_\_\_ \\ \hline
    \all{einhalten} & \all{hält ein} & \_\_\_\_\_\_\_\_\_\_ & \_\_\_\_\_\_\_\_\_\_ \\ \hline
    \end{tabular}
    \end{center}
    \item \textbf{Production :} Écrivez 4 phrases au prétérit pour raconter le déroulement d'un rite de veuvage (imaginaire ou tiré du texte). Utilisez : 1 verbe de mouvement, 1 verbe d'action rituelle, 1 verbe modal, 1 verbe d'état. Intégrez une proposition relative dans l'une d'elles.
\end{enumerate}
\tcblower
\textbf{Solution détaillée :}

\textbf{1. Propositions relatives (la déclinaison du relatif dépend de sa fonction DANS la relative) :}
\begin{enumerate}
    \item \textbf{\all{Die Witwe trinkt das Wasser, mit dem der Leichnam gewaschen wurde.}}
    (Antécédent \all{das Wasser}, neutre. Dans la relative : \all{mit} + datif $\rightarrow$ \all{dem}. Verbe conjugué \all{wurde} à la fin.)
    \item \textbf{\all{Die Verwandten beobachten die Witwe, die die schwarze Kleidung trägt.}}
    (Antécédent \all{die Witwe}, féminin. Fonction dans la relative : sujet $\rightarrow$ nominatif \all{die}. Verbe \all{trägt} à la fin.)
    \item \textbf{\all{Der Brauch, von dem der Ethnologe berichtet hat, ist alt.}}
    (Antécédent \all{der Brauch}, masculin. Préposition \all{von} + datif $\rightarrow$ \all{dem}. Verbe \all{hat ... berichtet} à la fin.)
    \item \textbf{\all{Die Witwenrente, ohne die die Witwe arm wäre, sichert das Überleben.}}
    (Antécédent \all{die Witwenrente}, féminin. Préposition \all{ohne} + accusatif $\rightarrow$ \all{die}. Konjunktiv II \all{wäre} car hypothèse irréelle.)
\end{enumerate}

\textbf{2. Tableau de conjugaison (prétérit / participe II) :}
\begin{center}
\begin{tabular}{|l|c|c|c|}\hline
\rowcolor{popblueL} \textbf{Infinitif} & \textbf{Présent (3e sg)} & \textbf{Prétérit (3e sg)} & \textbf{Participe II} \\ \hline
\all{sterben} (fort) & \all{stirbt} & \textbf{\all{starb}} & \textbf{\all{gestorben}} \\ \hline
\all{trinken} (fort) & \all{trinkt} & \textbf{\all{trank}} & \textbf{\all{getrunken}} \\ \hline
\all{werden} (fort/irrégulier) & \all{wird} & \textbf{\all{wurde}} & \textbf{\all{geworden}} \\ \hline
\all{sein} (irrégulier) & \all{ist} & \textbf{\all{war}} & \textbf{\all{gewesen}} \\ \hline
\all{geben} (fort) & \all{gibt} & \textbf{\all{gab}} & \textbf{\all{gegeben}} \\ \hline
\all{bewältigen} (faible, préfixe \all{be-}) & \all{bewältigt} & \textbf{\all{bewältigte}} & \textbf{\all{bewältigt}} (\textbf{pas de ge-}) \\ \hline
\all{rasieren} (faible) & \all{rasiert} & \textbf{\all{rasierte}} & \textbf{\all{rasiert}} \\ \hline
\all{einhalten} (séparable, fort \all{halten}) & \all{hält ein} & \textbf{\all{hielt ein}} & \textbf{\all{eingehalten}} (\textbf{ge- entre préfixe et radical}) \\ \hline
\end{tabular}
\end{center}

\textbf{3. Production (exemple de réponse attendue) :}
\begin{enumerate}
    \item \textbf{\all{Die Witwe ging in das dunkle Haus.}} (Mouvement : \all{gehen $\rightarrow$ ging}.)
    \item \textbf{\all{Sie trank das bittere Wasser, das der Priester ihr gab.}} (Action : \all{trinken $\rightarrow$ trank} ; relative : \all{das} accusatif neutre, \all{gab} prétérit de \all{geben} à la fin.)
    \item \textbf{\all{Sie musste die ganze Nacht wachen.}} (Modal : \all{müssen $\rightarrow$ musste} + infinitif \all{wachen} à la fin.)
    \item \textbf{\all{Sie war sehr traurig, aber stark.}} (État : \all{sein $\rightarrow$ war}.)
\end{enumerate}
\end{exoresolu}

\begin{methode}[Expression écrite guidée : résumé structuré et prise de position argumentée]
\textbf{Objectif :} Produire un texte cohérent (100-150 mots) : \textbf{résumé} (fidèle, neutre, avec son propre vocabulaire) ou \textbf{commentaire/prise de position} (structuré : thèse, arguments, exemples, conclusion).\par
\textbf{Méthode « Résumé » (4 étapes) :}
\begin{enumerate}
    \item \textbf{Découper} le texte en séquences logiques (introduction, aspects Cameroun, aspects Allemagne, comparaison/conclusion).
    \item \textbf{Prise de notes} : mots-clés allemands par séquence (pas de phrases entières). Ex. : \all{Kamerun: strenge Rituale, Hausarrest, Unschuldsbeweis / Deutschland: Privatsache, Witwenrente, Freiheit}.
    \item \textbf{Rédiger en allemand} : phrases courtes, connecteurs (\all{zuerst}, \all{dann}, \all{im Gegensatz dazu}, \all{außerdem}, \all{schließlich}). \textbf{Présent} pour les faits généraux, \textbf{prétérit} pour le récit passé, \textbf{passif} pour les rites, \textbf{relatives} pour lier les idées.
    \item \textbf{Contrôler} : nombre de mots, orthographe (majuscules des noms !), cas, ordre des mots (verbe conjugué à la fin en subordonnée).
\end{enumerate}
\textbf{Méthode « Commentaire / Stellungnahme » (5 étapes) :}
\begin{enumerate}
    \item \textbf{Analyser la citation/le sujet} : définir les termes, identifier le conflit (tradition vs modernité / collectif vs individuel / protection vs oppression).
    \item \textbf{Thèse claire} : \all{Ich stimme zu / nicht zu, weil...} ou \all{Meiner Meinung nach ist... entscheidend.}
    \item \textbf{Arguments (2-3) + exemples (texte + connaissances personnelles / réalité camerounaise) :} structure \all{Erstens... Zweitens... Ein Beispiel dafür ist...}.
    \item \textbf{Contre-argument (nuance) :} \all{Allerdings / Andererseits muss man sagen, dass...}
    \item \textbf{Conclusion} : \all{Zusammenfassend / Abschließend lässt sich sagen, dass...} Ouverture possible.
\end{enumerate}
\textbf{Boîte à outils : connecteurs (Redemittel) :}
\all{einerseits ... andererseits}, \all{nicht nur ... sondern auch}, \all{außerdem}, \all{deshalb / daher}, \all{obwohl / trotzdem}, \all{damit / um ... zu}.
\end{methode}

\begin{exoresolu}[Expression écrite : résumé et prise de position]
\textbf{Énoncé :}
\begin{enumerate}
    \item \textbf{Résumé (env. 80 mots) :} Fassen Sie den Text „Witwenrituale in Kamerun und in Deutschland“ zusammen. Nutzen Sie eigene Wörter (kein wörtliches Zitieren). Struktur: Allgemein $\rightarrow$ Kamerun $\rightarrow$ Deutschland $\rightarrow$ Vergleich.
    \item \textbf{Stellungnahme (env. 100 mots) :} Nehmen Sie Stellung zu folgender These: „Moderne Gesetze schützen die Witwe besser als alte Traditionen.“ Beziehen Sie sich auf den Text und Ihre Kenntnisse der kamerunischen Realität.
\end{enumerate}
\tcblower
\textbf{Solution détaillée :}

\textbf{1. Résumé (Musterlösung — env. 85 mots) :}
\begin{textebox}[Musterlösung Zusammenfassung]
Der Text thematisiert den unterschiedlichen Umgang mit Witwen in Kamerun und Deutschland. In kamerunischen Kulturen unterliegt die Witwe strengen, traditionellen Ritualen: Sie muss eine lange Trauerzeit einhalten, schwarze Kleidung tragen, sich den Kopf rasieren und ihre Unschuld durch riskante Proben beweisen. Der soziale Druck der Verwandten ist enorm. In Deutschland hingegen ist die Trauer Privatsache. Es gibt keine vorgeschriebenen Bräuche; psychologische Hilfe und die staatliche Witwenrente sichern die Existenz. Trotz des gemeinsamen Gedenkens an den Verstorbenen genießt die Witwe in Deutschland deutlich mehr persönliche Freiheit und rechtliche Absicherung.
\end{textebox}
\emph{Analyse :} présent de vérité générale ; passif et constructions modales (\all{unterliegt}, \all{muss ... einhalten}) ; relatives ; comparatif (\all{deutlich mehr}) ; vocabulaire précis (\all{thematisieren}, \all{unterliegen} + datif, \all{die Probe}, \all{die Existenz}, \all{die rechtliche Absicherung}).

\textbf{2. Stellungnahme (Musterlösung — env. 110 mots) :}
\begin{textebox}[Musterlösung Stellungnahme]
Ich stimme der These weitgehend zu. Moderne Gesetze, wie das deutsche Erbrecht oder die Witwenrente, bieten der Witwe eine existenzielle Sicherheit, die traditionelle Rituale nicht garantieren können. In Kamerun schützt die Tradition zwar die soziale Einbettung der Witwe in die Großfamilie, aber sie geht oft mit Menschenrechtsverletzungen einher: Zwangsheirat, Enteignung oder demütigende Unschuldsproben, wie der Text zeigt. Allerdings darf man nicht verallgemeinern: In ländlichen Regionen Kameruns, wo der Staat abwesend ist, bleibt die traditionelle Solidarität der Verwandten oft die einzige Überlebensgarantie. Ein idealer Weg wäre die Verbindung: rechtliche Absicherung durch den Staat bei Respekt vor der kulturellen Identität.
\end{textebox}
\emph{Analyse :} thèse claire (\all{Ich stimme ... zu}) ; arguments (sécurité légale, droits humains / référence au texte) ; nuance (\all{Allerdings} : réalité camerounaise, solidarité familiale) ; conclusion synthétique (\all{die Verbindung}). Grammaire : relatives (\all{die ... garantieren können}, \all{wo der Staat abwesend ist}), Konjunktiv II (\all{wäre}), connecteurs (\all{zwar ... aber}, \all{allerdings}).
\end{exoresolu}

\begin{attention}[Les 5 erreurs qui coûtent des points au Bac]
\begin{enumerate}
    \item \textbf{Majuscules oubliées sur les noms :} \all{die witwe} $\rightarrow$ \textbf{\all{die Witwe}}, \all{der brauch} $\rightarrow$ \textbf{\all{der Brauch}}, \all{das ritual} $\rightarrow$ \textbf{\all{das Ritual}}. \emph{Règle : tout nom commun prend une majuscule en allemand.}
    \item \textbf{Pluriels et Umlauts incorrects :} \all{die Witwens} $\rightarrow$ \textbf{\all{die Witwen}} ; \all{die Brauche} $\rightarrow$ \textbf{\all{die Bräuche}} (Umlaut \cle{au $\rightarrow$ äu}) ; \all{die Rituals} $\rightarrow$ \textbf{\all{die Rituale}} ; pluriel de \all{der Verstorbene} $\rightarrow$ \textbf{\all{die Verstorbenen}}.
    \item \textbf{Faux amis / calques lexicaux :} \all{beweisen} (prouver) $\neq$ \all{versprechen} (promettre) ; \all{die Leiche} (le cadavre, registre péjoratif) à éviter au profit de \textbf{\all{der Leichnam}} (la dépouille, registre respectueux) ; \all{die Tradition} (nom) $\neq$ \all{traditionell} (adjectif) ; \all{die Trauer} (le deuil) $\neq$ \all{traurig} (triste).
    \item \textbf{Ordre des mots (place du verbe) :} principale : verbe conjugué en 2\textsuperscript{e} position (\all{Heute \textbf{wird} das Ritual kritisch gesehen.}) ; subordonnée (\all{dass}, \all{weil}, \all{wenn}, \all{ob}, relative) : verbe conjugué \textbf{à la fin} (\all{..., \textbf{dass} die Witwe beobachtet \textbf{wird}.}) ; passif avec modal : \all{Die Witwe \textbf{muss} beobachtet \textbf{werden}} (double infinitif à la fin).
    \item \textbf{Cas (Kasus) après prépositions et verbes :} \all{mit der Witwe} (datif), pas \all{*mit die Witwe} ; \all{von dem Brauch} (datif) $\rightarrow$ relatif \all{von \textbf{dem}} ; \all{ohne die Rente} (accusatif) ; \all{der Witwe helfen} (\all{helfen} + datif !), pas \all{*die Witwe helfen} ; \all{der Hilfe bedürfen} (\all{bedürfen} + génitif !).
\end{enumerate}
\end{attention}

\newpage

\coursec{Exercices}

\courssub{Je vérifie mes connaissances}

\begin{exercice}[QCM : le passif]
Entourez la bonne réponse.
\begin{enumerate}
\item \all{Die Haare der Witwe \ldots\ abgeschnitten.}\\
a) werden \quad b) wurden \quad c) sind worden
\item \all{Das Ritual \ldots\ von den Dorfbewohnern gepflegt.}\\
a) wird \quad b) wird geworden \quad c) wurde werden
\item \all{Die Witwe \ldots\ ein Jahr lang schwarze Kleidung.} (Präteritum de \all{tragen})\\
a) tragt \quad b) trug \quad c) hat getragen
\item \all{Die Trauerfeier \ldots\ gestern im Dorf \ldots\ .} (Perfekt, passif)\\
a) wurde \ldots\ gefeiert \quad b) ist \ldots\ gefeiert worden \quad c) hat \ldots\ gefeiert
\end{enumerate}
\end{exercice}

\begin{exercice}[Vrai ou faux ? Justifiez]
Répondez par \emph{vrai} ou \emph{faux} et justifiez par une phrase en allemand.
\begin{enumerate}
\item \all{Eine Witwe ist ein Mann, dessen Frau gestorben ist.}
\item Au passif, le verbe \all{werden} se conjugue et le participe passé se place à la fin de la phrase.
\item Le prétérit de \all{sterben} est \all{sterbte}.
\item Dans une proposition relative introduite par \all{den}, le verbe conjugué se place à la fin.
\end{enumerate}
\end{exercice}

\begin{exercice}[Vocabulaire : deuil et veuvage]
Complétez chaque phrase avec le mot qui convient : \all{die Witwe}, \all{der Witwer}, \all{der Brauch}, \all{die Trauer}, \all{das Ritual}, \all{die Beerdigung}.
\begin{enumerate}
\item Nach dem Tod ihres Mannes trug \ldots\ ein Jahr lang schwarze Kleidung.
\item Herr Mbarga ist seit zwei Jahren \ldots\ ; seine Frau ist bei einem Unfall gestorben.
\item In vielen Dörfern Kameruns ist das Schneiden der Haare ein alter \ldots\ .
\item Die Familie zeigt ihre \ldots\ durch Gebete und traditionelle Lieder.
\item Die \ldots\ fand am Samstag auf dem Friedhof des Dorfes statt.
\end{enumerate}
\end{exercice}

\begin{exercice}[Conjugaison : prétérit et participe]
Complétez le tableau avec le prétérit (3\up{e} personne du singulier) et le participe passé, en choisissant le bon auxiliaire (\all{haben} ou \all{sein}).

\begin{center}
\begin{tabular}{|l|l|l|l|}
\hline
\rowcolor{popblueL} \textbf{Infinitiv} & \textbf{Präteritum} & \textbf{Auxiliaire} & \textbf{Partizip II} \\
\hline
sterben & & & \\
\hline
tragen & & & \\
\hline
schneiden & & & \\
\hline
bleiben & & & \\
\hline
verbieten & & & \\
\hline
\end{tabular}
\end{center}
\end{exercice}

\courssub{Je m'entraîne}

\begin{exercice}[Thème : traduction en allemand]
Traduisez les phrases suivantes en allemand.
\begin{enumerate}
\item Après la mort de son mari, la veuve a porté des vêtements noirs pendant un an.
\item Les rites qui sont observés dans son village sont très anciens.
\item Les cheveux de la veuve ont été coupés par les femmes du quartier.
\item Le veuf est resté sept jours dans sa maison.
\end{enumerate}
\end{exercice}

\begin{exercice}[De l'actif au passif]
Mettez les phrases à la forme passive, au temps indiqué entre parenthèses.
\begin{enumerate}
\item \all{Die Frauen schneiden die Haare.} (Präteritum)
\item \all{Die Dorfbewohner pflegen den alten Brauch.} (Präsens)
\item \all{Die Familie hat die Trauerfeier organisiert.} (Perfekt)
\item \all{Man verbot der Witwe das Verlassen des Hauses.} (Präteritum)
\end{enumerate}
\end{exercice}

\begin{exercice}[Propositions relatives]
Reliez les deux phrases par un pronom relatif qui convient.
\begin{enumerate}
\item \all{Der Brauch ist sehr alt. Die Dorfbewohner pflegen ihn.}
\item \all{Die Witwe wohnt in Yaoundé. Ihr Mann ist letztes Jahr gestorben.}
\item \all{Das Haus steht am Rande des Dorfes. Die Trauerfeier findet in dem Haus statt.}
\item \all{Die Frauen singen traditionelle Lieder. Die Frauen tragen schwarze Kleider.}
\end{enumerate}
\end{exercice}

\begin{exercice}[Texte à trous : un rituel]
Complétez le texte avec les formes passives des verbes entre parenthèses, au Präsens.

\begin{textebox}[Ein Ritual im Dorf]
Nach dem Tod eines Mannes (1) \ldots\ die ganze Familie \ldots\ (informieren). Die Witwe (2) \ldots\ in ein besonderes Zimmer \ldots\ (bringen). Dort (3) \ldots\ ihre Haare von den ältesten Frauen \ldots\ (schneiden). Sieben Tage lang (4) \ldots\ Gebete \ldots\ (sprechen) und traditionelle Lieder (5) \ldots\ (singen). Am Ende der Trauerzeit (6) \ldots\ ein großes Fest \ldots\ (feiern), und die Witwe (7) \ldots\ wieder in die Gemeinschaft \ldots\ (aufnehmen).
\end{textebox}
\end{exercice}

\courssub{Je me prépare au Bac}

\begin{exercice}[Type Bac — Compréhension écrite]
Lisez le texte suivant, puis répondez aux questions.

\begin{textebox}[Witwenrituale in Kamerun und in Deutschland]
Als Frau Ngono aus einem Dorf in der Nähe von Ebolowa ihren Mann verlor, begann für sie eine lange Zeit der Trauer. Zuerst wurden ihre Haare von den ältesten Frauen der Familie abgeschnitten. Dann trug sie ein Jahr lang nur schwarze Kleidung und durfte das Haus nicht allein verlassen. Jeden Abend sprachen die Nachbarn mit ihr Gebete, und am Wochenende wurden traditionelle Lieder gesungen. Diese Bräuche, die seit Generationen gepflegt werden, sollen die Witwe schützen und ihr helfen, den Schmerz zu überwinden. Nach einem Jahr wurde ein großes Fest gefeiert : die Trauerzeit war beendet, und Frau Ngono wurde wieder in das Leben der Gemeinschaft aufgenommen.

In Deutschland ist die Situation anders. Die Witwe oder der Witwer trägt oft nur am Tag der Beerdigung schwarze Kleidung. Es gibt keine festen Rituale, die befolgt werden müssen. Viele Menschen suchen Hilfe bei einem Psychologen oder in einer Trauergruppe. Trotzdem bleibt die Erinnerung an den Verstorbenen wichtig : Fotos werden aufgestellt, und der Todestag wird jedes Jahr mit einer Kerze und einem Gebet begangen. Ob in Kamerun oder in Deutschland — überall sucht die Familie einen Weg, mit dem Verlust zu leben.
\end{textebox}

\textbf{A. Questions de compréhension.} Répondez en allemand, par des phrases complètes.
\begin{enumerate}
\item Was geschah mit den Haaren von Frau Ngono, und von wem?
\item Wie lange dauerte die Trauerzeit, und was durfte die Witwe während dieser Zeit nicht tun?
\item Welchen Zweck haben die Bräuche, die im Dorf gepflegt werden?
\item Wie unterscheidet sich die Trauer in Deutschland von der Trauer in Kamerun? Nennen Sie zwei Unterschiede.
\item Wie wird in Deutschland der Todestag des Verstorbenen begangen?
\end{enumerate}

\textbf{B. Vrai ou faux ?} Justifiez chaque réponse par une phrase du texte.
\begin{enumerate}
\item Frau Ngono konnte während der Trauerzeit allein aus dem Haus gehen.
\item In Deutschland gibt es strenge Rituale, die alle Witwen befolgen müssen.
\item In beiden Ländern versucht die Familie, mit dem Verlust zu leben.
\end{enumerate}

\textbf{C. Questions de langue.}
\begin{enumerate}
\item Relevez dans le texte trois phrases au passif et recopiez-les en indiquant le temps employé.
\item Relevez dans le texte une proposition relative et soulignez le pronom relatif.
\end{enumerate}
\end{exercice}

\begin{exercice}[Type Bac — Thème et version]
\textbf{A. Thème.} Traduisez en allemand le texte suivant.

\begin{textebox}[Texte à traduire]
Quand le vieux Tamo est mort, tout le village était triste. Sa femme a pleuré longtemps. Ses cheveux ont été coupés par les femmes de la famille, et elle a porté des vêtements noirs pendant plusieurs mois. Les traditions qui sont respectées dans cette région sont très anciennes. À la fin de la période de deuil, une grande fête a été organisée par les voisins.
\end{textebox}

\textbf{B. Version.} Traduisez en français le texte suivant.

\begin{textebox}[Text zum Übersetzen]
In vielen Dörfern Kameruns werden die Rituale des Todes noch heute gepflegt. Wenn ein Mann stirbt, wird die Witwe von den Frauen der Familie begleitet. Sie bleibt mehrere Tage im Haus und spricht mit niemandem. Die Haare werden abgeschnitten, und alte Lieder werden gesungen. Diese Bräuche, die von den Großeltern weitergegeben wurden, helfen der Familie, den schweren Verlust zu ertragen.
\end{textebox}
\end{exercice}

\begin{exercice}[Type Bac — Production écrite]
Rédigez un texte en allemand de \textbf{80 à 100 mots} sur le sujet suivant :

\begin{textebox}[Thema]
\all{Beschreiben Sie ein Trauerritual in Ihrer Region und vergleichen Sie es mit Deutschland. Geben Sie am Ende Ihre eigene Meinung zu diesen Traditionen.}
\end{textebox}

\textbf{Conseils de méthode :}
\begin{itemize}
\item Décrivez le rituel en employant au moins \textbf{trois phrases au passif} (\all{werden + Partizip II}).
\item Utilisez au moins \textbf{deux propositions relatives} (\all{der, die, das, den, dem}).
\item Employez le \textbf{prétérit} pour raconter les événements passés.
\item Terminez par votre opinion personnelle : \all{Ich finde diese Traditionen \ldots, weil \ldots}
\end{itemize}
\end{exercice}

\newpage

\coursec{Corrigés des exercices}

\begin{corrige}[Exercice 1 — QCM : le passif]
\textbf{Réponses :} 1. b) ; 2. a) ; 3. b) ; 4. b).

\textbf{Justifications :}
\begin{enumerate}
\item \all{Die Haare der Witwe wurden abgeschnitten.} — Le passif au Präteritum : \all{werden} au Präteritum (\all{wurden}) + participe passé en fin de phrase. « Les cheveux de la veuve furent coupés. »
\item \all{Das Ritual wird von den Dorfbewohnern gepflegt.} — Passif au Präsens : \all{wird} + \all{gepflegt}. « Le rituel est entretenu par les habitants du village. »
\item \all{Die Witwe trug ein Jahr lang schwarze Kleidung.} — Prétérit du verbe fort \all{tragen} : \all{trug} (3\up{e} pers. sg.). « La veuve porta des vêtements noirs pendant un an. »
\item \all{Die Trauerfeier ist gestern im Dorf gefeiert worden.} — Passif au Perfekt : \all{sein} (au Présens) + participe passé + \all{worden}. Attention : on dit \all{worden} et jamais \all{geworden} au passif. « La cérémonie de deuil a été célébrée hier au village. »
\end{enumerate}
\end{corrige}

\begin{corrige}[Exercice 2 — Vrai ou faux ? Justifiez]
\begin{enumerate}
\item \textbf{Faux.} \all{Eine Witwe ist eine Frau, deren Mann gestorben ist.} („Une veuve est une femme dont le mari est mort.“) Un homme dont la femme est morte est \all{ein Witwer} (un veuf).
\item \textbf{Vrai.} \all{Beim Passiv wird „werden“ konjugiert, und das Partizip II steht am Ende des Satzes.} Exemple : \all{Die Haare werden geschnitten.}
\item \textbf{Faux.} \all{Das Präteritum von „sterben“ ist „starb“.} \all{sterben} est un verbe fort : \all{sterben — starb — ist gestorben}. « Il mourut » = \all{er starb}.
\item \textbf{Vrai.} \all{In einem Relativsatz mit „den“ steht das konjugierte Verb am Ende.} Exemple : \all{Der Brauch, den die Dorfbewohner pflegen, ist sehr alt.}
\end{enumerate}
\end{corrige}

\begin{corrige}[Exercice 3 — Vocabulaire : deuil et veuvage]
\begin{enumerate}
\item \all{Nach dem Tod ihres Mannes trug die Witwe ein Jahr lang schwarze Kleidung.} (« Après la mort de son mari, la veuve porta des vêtements noirs pendant un an. »)
\item \all{Herr Mbarga ist seit zwei Jahren Witwer; seine Frau ist bei einem Unfall gestorben.} (« M. Mbarga est veuf depuis deux ans ; sa femme est morte dans un accident. »)
\item \all{In vielen Dörfern Kameruns ist das Schneiden der Haare ein alter Brauch.} (« Dans de nombreux villages du Cameroun, la coupe des cheveux est une vieille coutume. »)
\item \all{Die Familie zeigt ihre Trauer durch Gebete und traditionelle Lieder.} (« La famille montre son deuil par des prières et des chants traditionnels. »)
\item \all{Die Beerdigung fand am Samstag auf dem Friedhof des Dorfes statt.} (« L'enterrement a eu lieu samedi au cimetière du village. »)
\end{enumerate}
\textbf{Rappel lexical :}
\begin{center}
\begin{tabular}{|l|l|l|}
\hline
\rowcolor{popblueL} \textbf{Mot} & \textbf{Pluriel} & \textbf{Sens} \\
\hline
\all{die Witwe} & \all{die Witwen} & la veuve \\
\hline
\all{der Witwer} & \all{die Witwer} & le veuf \\
\hline
\all{der Brauch} & \all{die Bräuche} & la coutume \\
\hline
\all{die Trauer} & \textendash & le deuil (invariable) \\
\hline
\all{das Ritual} & \all{die Rituale} & le rite, le rituel \\
\hline
\all{die Beerdigung} & \all{die Beerdigungen} & l'enterrement \\
\hline
\end{tabular}
\end{center}
\end{corrige}

\begin{corrige}[Exercice 4 — Conjugaison : prétérit et participe]
\begin{center}
\begin{tabular}{|l|l|l|l|}
\hline
\rowcolor{popblueL} \textbf{Infinitiv} & \textbf{Präteritum} & \textbf{Auxiliaire (Perfekt)} & \textbf{Partizip II} \\
\hline
\all{sterben} & \all{starb} & \all{sein} & \all{gestorben} \\
\hline
\all{tragen} & \all{trug} & \all{haben} & \all{getragen} \\
\hline
\all{schneiden} & \all{schnitt} & \all{haben} & \all{geschnitten} \\
\hline
\all{bleiben} & \all{blieb} & \all{sein} & \all{geblieben} \\
\hline
\all{verbieten} & \all{verbot} & \all{haben} & \all{verboten} \\
\hline
\end{tabular}
\end{center}
\textbf{Points de vigilance :}
\begin{itemize}
\item \all{sterben} et \all{bleiben} expriment un changement d'état ou de lieu : auxiliaire \all{sein} au Perfekt (\all{er ist gestorben}, \all{er ist geblieben}).
\item \all{verbieten} a un préfixe inséparable (\all{ver-}) : le participe ne prend pas \all{ge-} : \all{verboten}.
\item \all{schneiden} : verbe fort, alternance vocalique \all{ei} $\to$ \all{i} au Prétérit (\all{schnitt}) et au Participe (\all{geschnitten}).
\item \all{tragen} : alternance \all{a} $\to$ \all{u} $\to$ \all{a} (\all{tragen — trug — getragen}).
\end{itemize}
\end{corrige}

\begin{corrige}[Exercice 5 — Thème : traduction en allemand]
\begin{enumerate}
\item \all{Nach dem Tod ihres Mannes trug die Witwe ein Jahr lang schwarze Kleidung.}\\
(\all{trug} : Prétérit de \all{tragen} ; « pendant un an » = \all{ein Jahr lang}.)
\item \all{Die Riten, die in ihrem Dorf beobachtet werden, sind sehr alt.}\\
(Proposition relative au nominatif pluriel : \all{die} ; verbe à la fin ; passif au Présens \all{werden beobachtet}.)
\item \all{Die Haare der Witwe wurden von den Frauen des Viertels geschnitten.}\\
(Passif au Prétérit : \all{wurden} + participe en fin ; complément d'agent introduit par \all{von} + datif pluriel \all{den Frauen}.)
\item \all{Der Witwer blieb sieben Tage in seinem Haus.}\\
(\all{blieb} : Prétérit de \all{bleiben} ; « le veuf » = \all{der Witwer}, ne pas confondre avec \all{die Witwe}.)
\end{enumerate}
\end{corrige}

\begin{corrige}[Exercice 6 — De l'actif au passif]
\begin{enumerate}
\item \all{Die Haare wurden von den Frauen geschnitten.} (Prétérit : \all{wurden} + \all{geschnitten}.) « Les cheveux furent coupés par les femmes. »
\item \all{Der alte Brauch wird von den Dorfbewohnern gepflegt.} (Présens : \all{wird} + \all{gepflegt} ; l'accusatif \all{den alten Brauch} devient nominatif \all{der alte Brauch}.) « La vieille coutume est entretenue par les habitants du village. »
\item \all{Die Trauerfeier ist von der Familie organisiert worden.} (Perfekt du passif : \all{ist} + participe + \all{worden}.) « La cérémonie de deuil a été organisée par la famille. »
\item \all{Der Witwe wurde das Verlassen des Hauses verboten.} (Prétérit ; le datif \all{der Witwe} reste au datif au passif.) « Il fut interdit à la veuve de quitter la maison. »
\end{enumerate}
\textbf{Point de vigilance :} seul le complément d'objet direct (accusatif) devient sujet au passif ; le complément d'objet indirect (datif) ne change pas de cas. Avec \all{verbieten} (interdire qqc à qqn), la chose interdite (accusatif) devient sujet, la personne (datif) reste au datif.
\end{corrige}

\begin{corrige}[Exercice 7 — Propositions relatives]
\begin{enumerate}
\item \all{Der Brauch, den die Dorfbewohner pflegen, ist sehr alt.}\\
(\all{den} : accusatif masculin singulier, car \all{pflegen} gouverne l'accusatif ; virgules obligatoires ; verbe conjugué \all{pflegen} à la fin.) « La coutume que les habitants du village entretiennent est très ancienne. »
\item \all{Die Witwe, deren Mann letztes Jahr gestorben ist, wohnt in Yaoundé.}\\
(\all{deren} : génitif féminin singulier, « dont le mari » ; antécédent \all{die Witwe} fém. sg.) « La veuve, dont le mari est mort l'année dernière, habite à Yaoundé. »
\item \all{Das Haus, in dem die Trauerfeier stattfindet, steht am Rande des Dorfes.}\\
(\all{in dem} : datif neutre singulier après la préposition \all{in} (lieu où) ; \all{stattfinden} gouverne \all{in} + datif.) « La maison dans laquelle la cérémonie de deuil a lieu se trouve au bord du village. »
\item \all{Die Frauen, die schwarze Kleider tragen, singen traditionelle Lieder.}\\
(\all{die} : nominatif pluriel, sujet de \all{tragen}.) « Les femmes qui portent des vêtements noirs chantent des chants traditionnels. »
\end{enumerate}
\end{corrige}

\begin{corrige}[Exercice 8 — Texte à trous : un rituel]
\begin{textebox}[Ein Ritual im Dorf — Lösung]
Nach dem Tod eines Mannes (1) \textbf{wird} die ganze Familie informiert. Die Witwe (2) \textbf{wird} in ein besonderes Zimmer gebracht. Dort (3) \textbf{werden} ihre Haare von den ältesten Frauen geschnitten. Sieben Tage lang (4) \textbf{werden} Gebete gesprochen und traditionelle Lieder (5) \textbf{werden} gesungen. Am Ende der Trauerzeit (6) \textbf{wird} ein großes Fest gefeiert, und die Witwe (7) \textbf{wird} wieder in die Gemeinschaft aufgenommen.
\end{textebox}
\textbf{Traduction :} « Après la mort d'un homme, toute la famille est informée. La veuve est conduite dans une pièce spéciale. Là, ses cheveux sont coupés par les femmes les plus âgées. Pendant sept jours, des prières sont dites et des chants traditionnels sont chantés. À la fin de la période de deuil, une grande fête est célébrée, et la veuve est de nouveau accueillie dans la communauté. »

\textbf{Points de vigilance :}
\begin{itemize}
\item Accord sujet-verbe : (1) et (2) sujet singulier (\all{die ganze Familie}, \all{die Witwe}) $\to$ \all{wird} ; (3), (4), (5) sujets pluriels (\all{Haare}, \all{Gebete}, \all{Lieder}) $\to$ \all{werden} ; (6) sujet singulier (\all{ein großes Fest}) $\to$ \all{wird} ; (7) sujet singulier (\all{die Witwe}) $\to$ \all{wird}.
\item (2) \all{bringen} (verbe faible mixte) : participe \all{gebracht}.
\item (7) \all{aufnehmen} (verbe à particule séparable) : participe \all{aufgenommen} (\all{ge-} inséré entre la particule \all{auf-} et le radical \all{nehm-}).
\end{itemize}
\end{corrige}

\begin{corrige}[Exercice 9 — Type Bac — Compréhension écrite]
\textbf{Barème indicatif (sur 20) :} A. 10 pts (2 pts par question) ; B. 6 pts (2 pts par item : 1 pt pour vrai/faux, 1 pt pour la justification) ; C. 4 pts (2 pts par question).

\textbf{A. Questions de compréhension — réponses modèles :}
\begin{enumerate}
\item \all{Ihre Haare wurden von den ältesten Frauen der Familie abgeschnitten.} (« Ses cheveux furent coupés par les femmes les plus âgées de la famille. »)
\item \all{Die Trauerzeit dauerte ein Jahr. Während dieser Zeit durfte die Witwe das Haus nicht allein verlassen.} (« La période de deuil dura un an. Pendant ce temps, la veuve ne pouvait pas quitter la maison seule. »)
\item \all{Die Bräuche sollen die Witwe schützen und ihr helfen, den Schmerz zu überwinden.} (« Ces coutumes doivent protéger la veuve et l'aider à surmonter la douleur. »)
\item \all{In Deutschland trägt die Witwe oft nur am Tag der Beerdigung schwarze Kleidung, und es gibt keine festen Rituale. Außerdem suchen viele Menschen Hilfe bei einem Psychologen oder in einer Trauergruppe.} (« En Allemagne, la veuve ne porte souvent du noir que le jour de l'enterrement, et il n'y a pas de rituels fixes. De plus, beaucoup cherchent de l'aide auprès d'un psychologue ou dans un groupe de deuil. »)
\item \all{Der Todestag wird jedes Jahr mit einer Kerze und einem Gebet begangen.} (« L'anniversaire de la mort est commémoré chaque année avec une bougie et une prière. »)
\end{enumerate}

\textbf{B. Vrai ou faux ?}
\begin{enumerate}
\item \textbf{Faux.} Le texte dit : \all{„Dann trug sie ein Jahr lang nur schwarze Kleidung und durfte das Haus nicht allein verlassen.“}
\item \textbf{Faux.} Le texte dit : \all{„Es gibt keine festen Rituale, die befolgt werden müssen.“}
\item \textbf{Vrai.} Le texte dit : \all{„Ob in Kamerun oder in Deutschland — überall sucht die Familie einen Weg, mit dem Verlust zu leben.“}
\end{enumerate}

\textbf{C. Questions de langue :}
\begin{enumerate}
\item Trois phrases au passif (au choix parmi le texte) :
\begin{itemize}
\item \all{Zuerst wurden ihre Haare von den ältesten Frauen der Familie abgeschnitten.} (Prétérit)
\item \all{Am Wochenende wurden traditionelle Lieder gesungen.} (Prétérit)
\item \all{Diese Bräuche, die seit Generationen gepflegt werden, \ldots} (Présens — relative avec passif)
\item \all{Nach einem Jahr wurde ein großes Fest gefeiert.} (Prétérit)
\item \all{Fotos werden aufgestellt.} (Présens)
\end{itemize}
\item Proposition relative : \all{Diese Bräuche, \underline{die} seit Generationen gepflegt werden, sollen die Witwe schützen.} (ou : \all{Es gibt keine festen Rituale, \underline{die} befolgt werden müssen.})
\end{enumerate}
\textbf{Points de vigilance :} répondre par des phrases complètes avec le verbe en deuxième position ; ne pas recopier le texte sans adapter le pronom (\all{sie} $\to$ \all{die Witwe} ou \all{Frau Ngono}).
\end{corrige}

\begin{corrige}[Exercice 10 — Type Bac — Thème et version]
\textbf{Barème indicatif (sur 20) :} A. Thème 10 pts ; B. Version 10 pts. Compter la justesse grammaticale (passif, relatives, prétérit), le vocabulaire et l'orthographe (majuscules, Umlaute).

\textbf{A. Thème — traduction modèle :}

\begin{textebox}[Traduction modèle]
Als der alte Tamo starb, war das ganze Dorf traurig. Seine Frau weinte lange. Ihre Haare wurden von den Frauen der Familie geschnitten, und sie trug mehrere Monate lang schwarze Kleidung. Die Traditionen, die in dieser Region respektiert werden, sind sehr alt. Am Ende der Trauerzeit wurde ein großes Fest von den Nachbarn organisiert.
\end{textebox}
\textbf{Points de vigilance :}
\begin{itemize}
\item « Quand le vieux Tamo est mort » : subordonnée temporelle au Prétérit $\to$ \all{Als der alte Tamo starb} (et non \all{quand}/\all{wenn}, réservé au présent ou à la répétition dans le passé).
\item « ont été coupés » : passif au Prétérit $\to$ \all{wurden \ldots\ geschnitten} ; « a été organisée » : \all{wurde \ldots\ organisiert}.
\item « qui sont respectées » : relative au passif Présens $\to$ \all{die \ldots\ respektiert werden}.
\item « pendant plusieurs mois » : \all{mehrere Monate lang}.
\end{itemize}

\textbf{B. Version — texte source et traduction modèle :}

\begin{textebox}[Texte source (allemand)]
In vielen Dörfern Kameruns werden die Todesrituale heute noch gepflegt. Wenn ein Mann stirbt, wird die Witwe von den Frauen der Familie begleitet. Sie bleibt mehrere Tage im Haus und spricht mit niemandem. Die Haare werden abgeschnitten, und alte Lieder werden gesungen. Diese Bräuche, die von den Großeltern weitergegeben wurden, helfen der Familie, den schweren Schmerz zu ertragen.
\end{textebox}

\begin{textebox}[Traduction modèle]
Dans de nombreux villages du Cameroun, les rituels de la mort sont encore entretenus aujourd'hui. Quand un homme meurt, la veuve est accompagnée par les femmes de la famille. Elle reste plusieurs jours dans la maison et ne parle à personne. Les cheveux sont coupés, et de vieux chants sont chantés. Ces coutumes, qui ont été transmises par les grands-parents, aident la famille à supporter le lourd deuil.
\end{textebox}
\textbf{Points de vigilance :}
\begin{itemize}
\item \all{werden \ldots\ gepflegt} (passif présent) $\to$ « sont entretenus » ; \all{werden \ldots\ begleitet / abgeschnitten / gesungen} $\to$ passif présent.
\item \all{Wenn} + présent pour une règle générale (si/quand = chaque fois que).
\item \all{mit niemandem} = « avec personne / à personne » (négation).
\item \all{weitergegeben wurden} (passif prétérit) $\to$ « ont été transmises » (passé composé passif).
\item \all{ertragen} = « supporter, endurer » ; rendre le passif allemand par le passif ou une tournure naturelle en français (\all{helfen, den Schmerz zu ertragen} $\to$ « aident à supporter la douleur »).
\end{itemize}
\end{corrige}

\begin{corrige}[Exercice 11 — Type Bac — Production écrite]
\textbf{Barème indicatif (sur 20) :} respect du sujet et de la longueur (80–100 mots) : 4 pts ; emploi d'au moins trois passifs : 4 pts ; deux propositions relatives : 4 pts ; prétérit correct : 4 pts ; opinion personnelle et cohérence : 4 pts.

\textbf{Proposition de rédaction modèle (environ 95 mots) :}

\begin{textebox}[Modelltext]
In meiner Region, die in der Nähe von Bafoussam liegt, gibt es alte Trauerrituale. Als letztes Jahr der Vater meines Freundes starb, wurde die ganze Familie sofort informiert. Die Witwe, die sehr traurig war, wurde von den ältesten Frauen begleitet. Ihre Haare wurden abgeschnitten, und sie trug mehrere Monate schwarze Kleidung. Jeden Abend wurden Gebete gesprochen, und traditionelle Lieder wurden gesungen. Am Ende der Trauerzeit wurde ein großes Fest gefeiert. In Deutschland gibt es solche festen Rituale nicht. Ich finde diese Traditionen wichtig, weil sie der Familie helfen, den Schmerz zu überwinden.
\end{textebox}

\textbf{Traduction :} « Dans ma région, qui se trouve près de Bafoussam, il existe de vieux rituels de deuil. Quand l'année dernière le père de mon ami est mort, toute la famille a été immédiatement informée. La veuve, qui était très triste, a été accompagnée par les femmes les plus âgées. Ses cheveux ont été coupés, et elle a porté des vêtements noirs pendant plusieurs mois. Chaque soir, des prières étaient dites et des chants traditionnels étaient chantés. À la fin de la période de deuil, une grande fête a été célébrée. En Allemagne, de tels rituels fixes n'existent pas. Je trouve ces traditions importantes, parce qu'elles aident la famille à surmonter la douleur. »

\textbf{Points de vigilance :}
\begin{itemize}
\item Vérifier l'accord sujet–verbe au passif : \all{die Haare wurden} (pluriel), mais \all{die Witwe wurde} (singulier), \all{ein großes Fest wurde} (singulier).
\item Participe des verbes séparables : \all{abgeschnitten} (\all{abschneiden}), \all{aufgenommen} (\all{aufnehmen}).
\item Relative après préposition : \all{in dem}, \all{mit der} — verbe conjugué en fin de relative.
\item Clôture d'opinion : \all{Ich finde diese Traditionen wichtig, weil sie der Familie helfen, \ldots} — attention, \all{weil} renvoie le verbe à la fin ; \all{helfen} + datif (\all{der Familie}).
\end{itemize}
\end{corrige}

\newpage

\coursec{Fiche bilan}

\begin{popfigure}
\begin{center}\tcbox[colback=popblue,colframe=popblue,coltext=white,arc=9pt,boxsep=4pt,left=6pt,right=6pt]{\bfseries\small Witwen\-rituale (Rites de veuvage)}\end{center}
\vspace{-2pt}
\begin{tcbraster}[raster columns=3,raster equal height=rows,raster column skip=6pt,raster row skip=6pt,enhanced,blankest]
\begin{tcolorbox}[enhanced,colback=popgreen,colframe=popgreen,coltext=white,boxrule=0.9pt,arc=3pt,left=4pt,right=4pt,top=3pt,bottom=3pt,boxsep=1pt,halign=left]\bfseries\scriptsize Texte \emph{Witwenrituale in} \emph{Kamerun und} \emph{Deutschland}\end{tcolorbox}
\begin{tcolorbox}[enhanced,colback=poporange,colframe=poporange,coltext=white,boxrule=0.9pt,arc=3pt,left=4pt,right=4pt,top=3pt,bottom=3pt,boxsep=1pt,halign=left]\bfseries\scriptsize Lexique die Witwe, der Witwer, die Trauer, der Brauch\end{tcolorbox}
\begin{tcolorbox}[enhanced,colback=poppurple,colframe=poppurple,coltext=white,boxrule=0.9pt,arc=3pt,left=4pt,right=4pt,top=3pt,bottom=3pt,boxsep=1pt,halign=left]\bfseries\scriptsize Passiv werden + Partizip II \emph{Die Witwe wird} \emph{geführt.}\end{tcolorbox}
\begin{tcolorbox}[enhanced,colback=poppink,colframe=poppink,coltext=white,boxrule=0.9pt,arc=3pt,left=4pt,right=4pt,top=3pt,bottom=3pt,boxsep=1pt,halign=left]\bfseries\scriptsize Relativsätze der, die, das verbe à la fin\end{tcolorbox}
\begin{tcolorbox}[enhanced,colback=popgold,colframe=popgold,coltext=white,boxrule=0.9pt,arc=3pt,left=4pt,right=4pt,top=3pt,bottom=3pt,boxsep=1pt,halign=left]\bfseries\scriptsize Präteritum verbes forts kam, trug, sprach, blieb\end{tcolorbox}
\begin{tcolorbox}[enhanced,colback=popblue!70,colframe=popblue!70,coltext=white,boxrule=0.9pt,arc=3pt,left=4pt,right=4pt,top=3pt,bottom=3pt,boxsep=1pt,halign=left]\bfseries\scriptsize Partizip II ge-...-t / ge-...-en haben/sein\end{tcolorbox}
\begin{tcolorbox}[enhanced,colback=popdark,colframe=popdark,coltext=white,boxrule=0.9pt,arc=3pt,left=4pt,right=4pt,top=3pt,bottom=3pt,boxsep=1pt,halign=left]\bfseries\scriptsize Expression résumé, questions, court texte\end{tcolorbox}
\end{tcbraster}

\legende{Carte mentale de la leçon : les rites de veuvage}
\end{popfigure}

\begin{bilan}
\textbf{1. Lexique clé (allemand -- français)}
\begin{itemize}
  \item \all{die Witwe, -n} : la veuve ; \all{der Witwer, -} : le veuf
  \item \all{der Brauch, \"-e} : la coutume ; \all{das Ritual, -e} : le rituel ; \all{die Tradition, -en} : la tradition
  \item \all{die Trauer} : le deuil ; \all{trauern um + Akk.} : porter le deuil de ; \all{die Trauerfeier, -n} : la cérémonie funéraire
  \item \all{die Erben} (pl.) : les héritiers ; \all{das Erbe} : l'héritage ; \all{der Schwager, \"-} : le beau-frère
  \item \all{reinigen} : purifier ; \all{rasieren} : raser ; \all{schneiden (schnitt, geschnitten)} : couper
  \item \all{respektieren / achten} : respecter ; \all{ablehnen} : rejeter ; \all{verzichten auf + Akk.} : renoncer à
  \item \all{die Witwenrente, -n} : la pension de veuvage ; \all{die Versorgung} : la prise en charge
\end{itemize}

\textbf{2. Grammaire à connaître par c\oe ur}
\begin{center}
\begin{tabularx}{\linewidth}{|X|X|}
\hline
\rowcolor{popblueL} \textbf{Règle} & \textbf{Exemple} \\
\hline
Passif présent : \emph{werden} (conjugué) + Partizip II & \all{Die Witwe wird in ein Zimmer geführt.} (On conduit la veuve dans une pièce.) \\
\hline
Passif prétérit : \emph{wurde(n)} + Partizip II & \all{Die Haare wurden geschnitten.} (Les cheveux furent coupés.) \\
\hline
Passif Perfekt : \emph{ist/sind} + Partizip II + \emph{worden} & \all{Das Ritual ist abgeschafft worden.} \\
\hline
Agent : \emph{von} + Datif ; moyen : \emph{mit/durch} & \all{Sie wird von den Verwandten betreut.} \\
\hline
Proposition relative : pronom relatif (der/die/das + cas) + verbe en fin & \all{Der Brauch, \emph{den} man pflegt, ist alt.} \\
\hline
Prétérit des verbes forts : radical modifié, sans \emph{-t} & \all{kommen -- kam ; tragen -- trug ; sprechen -- sprach ; bleiben -- blieb} \\
\hline
Partizip II : faibles \emph{ge...t}, forts \emph{ge...en}, en \emph{-ieren} sans \emph{ge-} & \all{gemacht, getragen, studiert} \\
\hline
\end{tabularx}
\end{center}

\textbf{3. Méthodes express}
\begin{itemize}
  \item \textbf{Comprendre le texte} : lire le titre, repérer les mots du champ lexical du deuil, puis répondre d'abord globalement (thème, pays, personnages), ensuite en détail (rites, durée, acteurs).
  \item \textbf{Traduire un passif} : repérer \emph{werden/wurden} + Partizip II, traduire par « être + participe » ou par « on » (\all{Es wird getanzt} = « on danse »).
  \item \textbf{Résumer} : 4 à 6 phrases au Präsens ou Perfekt, avec connecteurs (\all{zuerst, dann, danach, schließlich}), sans recopier le texte.
  \item \textbf{Produire} : décrire un rite au passif (\all{Zuerst wird ..., dann wird ...}) et donner son avis avec \all{Meiner Meinung nach ...}.
\end{itemize}
\end{bilan}

\begin{aretenir}[Checklist avant le Bac]
\begin{itemize}
  \item $\square$ Je connais l'article et le pluriel du vocabulaire du veuvage (\all{die Witwe, -n ; der Brauch, \"-e}).
  \item $\square$ Je sais former le passif au présent, au prétérit et au Perfekt (\emph{wird / wurde / ist worden}).
  \item $\square$ Je place correctement le Partizip II à la fin de la phrase au passif.
  \item $\square$ J'introduis l'agent avec \all{von} + Datif et le moyen avec \all{mit/durch}.
  \item $\square$ Je choisis le bon pronom relatif selon le genre, le nombre et le cas.
  \item $\square$ Je mets le verbe conjugué en dernière position dans la proposition relative.
  \item $\square$ Je connais le prétérit des verbes forts fréquents (\all{kam, trug, sprach, gab, blieb, nahm}).
  \item $\square$ Je forme les participes passés sans erreur (\all{ge...t / ge...en}, pas de \emph{ge-} avec \emph{-ieren}).
  \item $\square$ Je sais si l'auxiliaire du Perfekt est \emph{haben} ou \emph{sein} (mouvement, changement : \emph{sein}).
  \item $\square$ Je traduis un passif allemand par « on » quand le français l'exige.
  \item $\square$ Je rédige un résumé structuré avec connecteurs et majuscules à tous les noms.
  \item $\square$ Je peux comparer un rite camerounais et un rite allemand en quelques phrases.
\end{itemize}
\end{aretenir}

\findecours

\end{document}
